% Production et traduction de textes liés à la vie en famille et en société — Allemand 1ère A — Cours signé OnBuch+
% Programme officiel MINESEC. Compiler avec : tectonic ALA04-production-et-traduction-de-textes-lies-a-la-vie-en-famille-.tex
\newcommand{\DOCMATIERE}{Allemand}
\newcommand{\DOCNIVEAU}{1ère A}
\newcommand{\DOCMODULE}{Chapitre 1 : Vie familiale et sociale}
\newcommand{\DOCLECON}{ALA04}
\newcommand{\DOCTITRE}{Production et traduction de textes liés à la vie en famille et en société}
% OnBuch+ — préambule « cours signés » ENRICHI (compilé avec tectonic / XeLaTeX)
% Système visuel « Pop » d'OnBuch+ : fond crème (jamais blanc), bordures encre
% épaisses, ombres « dures » décalées, titres Archivo Black en capitales.
% Version MATHÉMATIQUES : graphes (pgfplots/tikz), boîtes pédagogiques
% (définition, propriété, méthode, attention, le savais-tu, exercice résolu,
% exercice, TP, bilan), page de garde et signature OnBuch+ ancrée en bas.
%
% Macros attendues dans main.tex AVANT \input{preamble} :
%   \DOCMATIERE  \DOCNIVEAU  \DOCMODULE  \DOCLECON (numéro)  \DOCTITRE
\documentclass[11pt]{article}

\usepackage{fontspec}
\usepackage[ngerman,french]{babel} % français = langue principale ; l'allemand via \all{..} / textebox
\usepackage[a4paper,margin=2cm,top=3.4cm,bottom=2.8cm,headheight=1.4cm,footskip=1.3cm]{geometry}
\usepackage{amsmath,amssymb,amsfonts}
\usepackage{mathtools} % \xmapsto, \coloneqq... souvent utilisés par les modèles
\usepackage{cancel} % simplifications barrées (bilans de forces)
% \abs{z} (module, valeur absolue) et \norm{u} (norme), utilisés par les modèles
\providecommand{\abs}{}\let\abs\relax\DeclarePairedDelimiter{\abs}{\lvert}{\rvert}
\providecommand{\norm}{}\let\norm\relax\DeclarePairedDelimiter{\norm}{\lVert}{\rVert}
\usepackage[table]{xcolor} % [table] : fournit aussi \rowcolors (alternance de lignes)
\usepackage{pagecolor}
\usepackage{tikz}
\usetikzlibrary{arrows.meta,positioning,calc,shapes.geometric,shapes.misc,shapes.arrows,decorations.pathmorphing,decorations.pathreplacing,decorations.markings,patterns,shadows,mindmap,trees,fit,backgrounds,matrix,intersections,angles,quotes}
\usepackage{pgfplots}
\usepackage[version=4]{mhchem} % équations nucléaires \ce{^{235}_{92}U}
\usepackage[european,siunitx]{circuitikz} % schémas électriques
\pgfplotsset{compat=1.18}
\usepgfplotslibrary{fillbetween,groupplots} % aires entre courbes (calcul intégral)
\usepackage{enumitem}
\usepackage{fancyhdr}
% [most] charge toutes les bibliothèques tcolorbox (listings, minted, raster,
% documentation...) et gonfle inutilement la mémoire TeX sur un document long
% (~300 boîtes) : on ne charge que ce qui est réellement utilisé.
\usepackage[skins,breakable]{tcolorbox}
\usepackage{graphicx}
\usepackage{colortbl}
\usepackage{booktabs}
\usepackage{tabularx}
\usepackage{diagbox} % case d'en-tête diagonale des tableaux à double entrée
% Colonnes extensibles centrée / gauche / droite (C, L, R), souvent utilisées par les modèles
\newcolumntype{C}{>{\centering\arraybackslash}X}
\newcolumntype{L}{>{\raggedright\arraybackslash}X}
\newcolumntype{R}{>{\raggedleft\arraybackslash}X}
\usepackage{array}
\usepackage{multicol}
\usepackage{multirow}
\usepackage{siunitx}
% parse-numbers=false : accepte aussi \SI{2,5 \times 10^{-3}}{...}
\sisetup{locale=FR,output-decimal-marker={,},group-minimum-digits=4,parse-numbers=false}
\DeclareSIUnit\ohm{\ensuremath{\symup{\Omega}}} % U+2126 absent de la police
\DeclareSIUnit\division{div} % réglages d'oscilloscope (ms/div, V/div)
\DeclareSIUnit\tr{tr} % tours (vitesse de rotation, ex. \SI{1500}{\tr\per\minute})
\DeclareSIUnit\molar{M} % concentration molaire (ex. \SI{3}{\molar}, \SI{1}{\micro\molar})
\DeclareSIUnit\an{an} % année (le modèle confond parfois avec \year, un compteur TeX) — ex. \SI{5}{\kilo\meter\per\an}
\DeclareSIUnit\FCFA{FCFA} % franc CFA (ex. \SI{500}{\FCFA\per\kilo\gram})
\DeclareSIUnit\degreeBrix{\textdegree Brix} % teneur en sucre d'un sirop/jus (réfractométrie)
\DeclareSIUnit\month{mois} % le modèle utilise parfois \month, un compteur TeX, comme unité de durée
\usepackage{qrcode}
% \degree : commande fréquemment utilisée par les modèles pour un angle ou une
% température en dehors de \SI{}{} (non fournie par les paquets chargés).
\providecommand{\degree}{^{\circ}}
% \qed (fin de démonstration), utilisé par les modèles sans amsthm
\providecommand{\qed}{\hfill$\square$}
% \barre : double barre des valeurs interdites dans un tableau de variations (tkz-tab)
\providecommand{\barre}{\ensuremath{\|}}
% environnement proof (amsthm non chargé) utilisé par les modèles
\ifdefined\proof\else\newenvironment{proof}[1][Démonstration]{\par\noindent\textit{#1.}\ }{\qed\par}\fi
\usepackage{lastpage}
\usepackage{hyperref}
\hypersetup{hidelinks}

% ============================================================
% Palette « Pop » OnBuch+ — mêmes hex que la classe P (plus_ui.dart)
% ============================================================
\definecolor{poporange}{HTML}{F59321}
\definecolor{popdark}{HTML}{E07A0C}
\definecolor{popcream}{HTML}{FFF6E8}
\definecolor{popink}{HTML}{1C1714}
\definecolor{poppurple}{HTML}{7A5AE0}
\definecolor{popgreen}{HTML}{2E9E6A}
\definecolor{poppink}{HTML}{E0457B}
\definecolor{popblue}{HTML}{2F7DE1}
\definecolor{popgold}{HTML}{C98A12}
\definecolor{popmuted}{HTML}{8A7F76}
\definecolor{popline}{HTML}{EADFD2}
\definecolor{popgray}{HTML}{8A7F76}
\colorlet{popgrey}{popgray}
% Alias tolérants (le modèle nomme parfois les couleurs différemment)
\colorlet{popwhite}{white}
\colorlet{popblack}{popink}
\colorlet{popred}{poppink}
\colorlet{popyellow}{popgold}
% Teintes claires (fonds de tableaux/encadrés) — le modèle nomme parfois
% les couleurs avec un suffixe L (ex. popblueL) sans les avoir définies.
\colorlet{poporangeL}{poporange!20}
\colorlet{popdarkL}{popdark!20}
\colorlet{poppurpleL}{poppurple!20}
\colorlet{popgreenL}{popgreen!20}
\colorlet{poppinkL}{poppink!20}
\colorlet{popblueL}{popblue!20}
\colorlet{popgoldL}{popgold!20}
\colorlet{popredL}{popred!20}
\colorlet{popgrayL}{popgray!20}
% Teintes claires pour fonds de tableaux / zones
\colorlet{popblueL}{popblue!10!white}
\colorlet{poporangeL}{poporange!14!white}
\colorlet{popgreenL}{popgreen!12!white}
\colorlet{poppurpleL}{poppurple!12!white}
\colorlet{poppinkL}{poppink!12!white}
\colorlet{popgoldL}{popgold!14!white}

\pagecolor{popcream}

% ============================================================
% Typographie
% ============================================================
\defaultfontfeatures{Ligatures=TeX}
\setmainfont{PlusJakartaSans-Medium.ttf}[Path=../../../fonts/,BoldFont=PlusJakartaSans-Bold.ttf]
% Symboles, API et emojis absents de la police principale : repli sur DejaVu Sans (caractères actifs)
\newfontface\symfont{DejaVuSans.ttf}[Path=../../../fonts/]
\makeatletter
\newcommand\symactive[1]{\catcode#1=13 \begingroup\lccode`\~=#1 \lowercase{\endgroup\def~}{{\symfont\char#1}}}
\newcommand\symblank[1]{\catcode#1=13 \begingroup\lccode`\~=#1 \lowercase{\endgroup\def~}{}}
\newcommand\symhyph[1]{\catcode#1=13 \begingroup\lccode`\~=#1 \lowercase{\endgroup\def~}{\mbox{-}}}
\newcount\symc
\newcommand\symrange[2]{\symc=#1 \@whilenum\symc<#2 \do{\expandafter\symactive\expandafter{\the\symc}\advance\symc by 1 }}
\newcommand\symblankrange[2]{\symc=#1 \@whilenum\symc<#2 \do{\expandafter\symblank\expandafter{\the\symc}\advance\symc by 1 }}
\makeatother
\symrange{"0250}{"02FF}\symactive{"2713}\symactive{"2714}\symactive{"2717}\symactive{"2718}\symactive{"2610}\symactive{"2611}\symactive{"2612}
\symhyph{"2011}\symhyph{"2E1A}\symblankrange{"1F300}{"1FAFF}\symblank{"FE0F}
% Maths en sans-serif (Fira Math) avec chiffres et lettres droites de Plus
% Jakarta, pour que les formules se fondent dans le texte.
\usepackage[math-style=ISO,bold-style=ISO]{unicode-math}
\setmathfont{FiraMath-Regular.otf}[Path=../../../fonts/]
\setmathfont{PlusJakartaSans-Medium.ttf}[Path=../../../fonts/,range={up/num,up/{Latin,latin},\mathup}]
\usepackage{icomma} % virgule décimale sans espace en mode math
% Caractères mathématiques Unicode absents des polices de texte (Plus Jakarta,
% Archivo Black) : rendus par leur équivalent mathématique, en texte comme en
% formule — sinon ils s'affichent en carré vide (ex. « Arithmétique dans ℤ »).
\usepackage{newunicodechar}
\newunicodechar{ℝ}{\ensuremath{\mathbb{R}}}
\newunicodechar{ℤ}{\ensuremath{\mathbb{Z}}}
\newunicodechar{ℂ}{\ensuremath{\mathbb{C}}}
\newunicodechar{ℕ}{\ensuremath{\mathbb{N}}}
\newunicodechar{ℚ}{\ensuremath{\mathbb{Q}}}
\newunicodechar{𝔻}{\ensuremath{\mathbb{D}}}
\newunicodechar{α}{\ensuremath{\alpha}}
\newunicodechar{β}{\ensuremath{\beta}}
\newunicodechar{γ}{\ensuremath{\gamma}}
\newunicodechar{δ}{\ensuremath{\delta}}
\newunicodechar{ε}{\ensuremath{\varepsilon}}
\newunicodechar{θ}{\ensuremath{\theta}}
\newunicodechar{λ}{\ensuremath{\lambda}}
\newunicodechar{μ}{\ensuremath{\mu}}
\newunicodechar{σ}{\ensuremath{\sigma}}
\newunicodechar{τ}{\ensuremath{\tau}}
\newunicodechar{φ}{\ensuremath{\varphi}}
\newunicodechar{ψ}{\ensuremath{\psi}}
\newunicodechar{ω}{\ensuremath{\omega}}
\newunicodechar{Σ}{\ensuremath{\Sigma}}
% majuscules produites par \MakeUppercase dans les titres (α-aminés -> Α)
\newunicodechar{Α}{\ensuremath{\alpha}}
\newunicodechar{Β}{\ensuremath{\beta}}
\newunicodechar{Γ}{\ensuremath{\Gamma}}
\newunicodechar{Δ}{\ensuremath{\Delta}}
\newunicodechar{Ε}{\ensuremath{\varepsilon}}
\newunicodechar{Θ}{\ensuremath{\Theta}}
\newunicodechar{Λ}{\ensuremath{\Lambda}}
\newunicodechar{Μ}{\ensuremath{\mu}}
\newunicodechar{Τ}{\ensuremath{\tau}}
\newunicodechar{Φ}{\ensuremath{\Phi}}
\newunicodechar{Ψ}{\ensuremath{\Psi}}
\newunicodechar{Ω}{\ensuremath{\Omega}}
\newunicodechar{↦}{\ensuremath{\mapsto}}
\newunicodechar{∈}{\ensuremath{\in}}
\newunicodechar{∉}{\ensuremath{\notin}}
\newunicodechar{∧}{\ensuremath{\wedge}}
\newunicodechar{⇔}{\ensuremath{\Leftrightarrow}}
\newunicodechar{⟺}{\ensuremath{\Longleftrightarrow}}
\newunicodechar{⇒}{\ensuremath{\Rightarrow}}
\newunicodechar{⟹}{\ensuremath{\Longrightarrow}}
\newunicodechar{∘}{\ensuremath{\circ}}
\newunicodechar{∪}{\ensuremath{\cup}}
\newunicodechar{∩}{\ensuremath{\cap}}
\newunicodechar{≡}{\ensuremath{\equiv}}
\newunicodechar{⊂}{\ensuremath{\subset}}
\newunicodechar{⊥}{\ensuremath{\perp}}
\newunicodechar{∃}{\ensuremath{\exists}}
\newunicodechar{∀}{\ensuremath{\forall}}
\newunicodechar{∛}{\ensuremath{\sqrt[3]{\,}}}
\newunicodechar{∜}{\ensuremath{\sqrt[4]{\,}}}
\newunicodechar{⃗}{\ensuremath{{}^{\to}}}
\newunicodechar{ⁿ}{\ifmmode^{n}\else\textsuperscript{\ensuremath{n}}\fi}
\newunicodechar{ˣ}{\ifmmode^{x}\else\textsuperscript{\ensuremath{x}}\fi}
\newunicodechar{ᵇ}{\ifmmode^{b}\else\textsuperscript{\ensuremath{b}}\fi}
\newunicodechar{ʳ}{\ifmmode^{r}\else\textsuperscript{\ensuremath{r}}\fi}
\newunicodechar{ᵘ}{\ifmmode^{u}\else\textsuperscript{\ensuremath{u}}\fi}
\newunicodechar{ʸ}{\ifmmode^{y}\else\textsuperscript{\ensuremath{y}}\fi}
\newunicodechar{ᵃ}{\ifmmode^{a}\else\textsuperscript{\ensuremath{a}}\fi}
\newunicodechar{ᵉ}{\ifmmode^{e}\else\textsuperscript{\ensuremath{e}}\fi}
\newunicodechar{ᵏ}{\ifmmode^{k}\else\textsuperscript{\ensuremath{k}}\fi}
\newunicodechar{ᵖ}{\ifmmode^{p}\else\textsuperscript{\ensuremath{p}}\fi}
\newunicodechar{ᴺ}{\ifmmode^{N}\else\textsuperscript{\ensuremath{N}}\fi}
\newunicodechar{ᶜ}{\ifmmode^{c}\else\textsuperscript{\ensuremath{c}}\fi}
\newunicodechar{ʰ}{\ifmmode^{h}\else\textsuperscript{\ensuremath{h}}\fi}
\newunicodechar{ᵠ}{\ifmmode^{q}\else\textsuperscript{\ensuremath{q}}\fi}
\newunicodechar{ₙ}{\ifmmode_{n}\else\textsubscript{\ensuremath{n}}\fi}
\newunicodechar{ₐ}{\ifmmode_{a}\else\textsubscript{\ensuremath{a}}\fi}
\newunicodechar{ᵢ}{\ifmmode_{i}\else\textsubscript{\ensuremath{i}}\fi}
\newunicodechar{ᵣ}{\ifmmode_{r}\else\textsubscript{\ensuremath{r}}\fi}
\newunicodechar{ₖ}{\ifmmode_{k}\else\textsubscript{\ensuremath{k}}\fi}
\newunicodechar{ᵦ}{\ifmmode_{\beta}\else\textsubscript{\ensuremath{\beta}}\fi}
\newunicodechar{ₓ}{\ifmmode_{x}\else\textsubscript{\ensuremath{x}}\fi}
\newfontfamily\popdisplay{ArchivoBlack-Regular.ttf}[Path=../../../fonts/]
\newfontfamily\poplogo{SpaceGrotesk-Bold.ttf}[Path=../../../fonts/]
\newfontfamily\popsemi{PlusJakartaSans-SemiBold.ttf}[Path=../../../fonts/]
\newfontfamily\popxbold{PlusJakartaSans-ExtraBold.ttf}[Path=../../../fonts/]
\newfontfamily\popxboldls{PlusJakartaSans-ExtraBold.ttf}[Path=../../../fonts/,LetterSpace=3.0]

\setlist[enumerate]{leftmargin=1.7em,itemsep=5pt,topsep=4pt}
\setlist[itemize]{leftmargin=1.7em,itemsep=4pt,topsep=4pt,label={\color{poporange}\textbullet}}
\setlength{\parindent}{0pt}
\setlength{\parskip}{5pt}
\renewcommand{\arraystretch}{1.25}

% pgfplots : style Pop par défaut
\pgfplotsset{
  every axis/.append style={axis line style={popink,line width=1pt},tick style={popink},
    grid style={popline,line width=0.5pt},label style={font=\small},tick label style={font=\footnotesize}},
  popcurve/.style={poporange,line width=1.8pt,no markers,smooth},
}
% Même style disponible dans un \draw TikZ simple (hors pgfplots)
\tikzset{popcurve/.style={poporange,line width=1.8pt}}
\tikzset{popgrid/.style={popline,very thin}}
\colorlet{popcurve}{poporange} % le nom du style est parfois utilisé comme couleur

% ============================================================
% Pill / squiggle
% ============================================================
\newcommand{\poppill}[3]{%
  \tikz[baseline=-0.55ex]\node[draw=popink,line width=1.2pt,rounded corners=2.6pt,%
    fill=#2,inner xsep=6.5pt,inner ysep=3pt]{%
    \popxboldls\fontsize{7}{7}\selectfont\color{#3}\MakeUppercase{#1}};%
}
\newcommand{\popsquiggle}[1]{%
  \tikz[baseline=-0.4ex]{
    \draw[#1,line width=2.6pt,line cap=round]
      plot [smooth] coordinates {(0,0.09) (0.34,0.01) (0.68,0.21) (1.02,0.01) (1.36,0.10)};
  }%
}
% Mot-clé surligné (marqueur orange sous le texte)
\newcommand{\cle}[1]{{\popxbold\color{popdark}#1}}

% ============================================================
% En-tête / pied
% ============================================================
\pagestyle{fancy}
\fancyhf{}
\renewcommand{\headrulewidth}{0pt}
\renewcommand{\footrulewidth}{0pt}
\lhead{\raisebox{-0.15ex}{\poplogo\fontsize{13}{13}\selectfont\color{popink}OnBuch\color{poporange}+}}
\rhead{\poppill{\DOCMATIERE\ \textbullet\ \DOCNIVEAU\ \textbullet\ Leçon \DOCLECON}{white}{popink}}
\lfoot{\popxbold\fontsize{7.6}{7.6}\selectfont\color{popmuted}\MakeUppercase{Cours signé OnBuch+}\ \textbullet\ {\popsemi\DOCTITRE}}
\rfoot{\tikz[baseline=-0.6ex]\node[rounded corners=8.5pt,fill=popink,minimum height=17pt,minimum width=17pt,inner xsep=5pt,inner sep=0]{\popxbold\fontsize{8}{8}\selectfont\color{popcream}\thepage\,/\,\pageref*{LastPage}};}

% ============================================================
% Titres
% ============================================================
\newcommand{\chaptitre}[2]{%
  \begin{tcolorbox}[enhanced,colback=poporange,colframe=popink,boxrule=2.6pt,arc=11pt,
      drop shadow=popink,left=15pt,right=15pt,top=12pt,bottom=11pt,before skip=3pt,after skip=18pt]
    \poppill{#2}{popink}{popcream}\par\vspace{9pt}
    {\raggedright\hyphenpenalty=10000\exhyphenpenalty=10000\popdisplay\fontsize{22}{24}\selectfont\color{popcream}\MakeUppercase{#1}\par}
  \end{tcolorbox}
}
% Section numérotée : pastille orange avec le numéro + titre Archivo
\newcounter{popsec}
\newcommand{\coursec}[1]{%
  \stepcounter{popsec}\setcounter{popsub}{0}%
  \par\vspace{16pt}\noindent
  \begin{minipage}{\linewidth}
  \tikz[baseline=-0.7ex]\node[draw=popink,line width=1.6pt,fill=poporange,rounded corners=4pt,minimum size=22pt,inner sep=0]{\popdisplay\fontsize{12}{12}\selectfont\color{popcream}\thepopsec};%
  \hspace{8pt}{\popdisplay\fontsize{15}{17}\selectfont\color{popink}\MakeUppercase{#1}}\par
  \vspace{4pt}\noindent{\color{poporange}\rule{36pt}{3.6pt}}
  \end{minipage}\par\nopagebreak\vspace{8pt}
}
\newcounter{popsub}
\newcommand{\courssub}[1]{%
  \stepcounter{popsub}%
  \par\vspace{9pt}\noindent{\popxbold\fontsize{11.5}{13}\selectfont\color{popink}{\color{poporange}\thepopsec.\thepopsub}\hspace{6pt}#1}\par\nopagebreak\vspace{4pt}
}

% ============================================================
% Boîtes pédagogiques — même recette : fond blanc, bordure épaisse colorée,
% ombre dure encre, badge plein. Titre optionnel [#1].
% ============================================================
\tcbset{popbase/.style={breakable,enhanced,colback=white,boxrule=2.3pt,arc=9pt,
  shadow={3pt}{-3pt}{0pt}{popink},left=14pt,right=14pt,top=11pt,bottom=11pt,before skip=11pt,after skip=15pt,
  fontupper=\popsemi\fontsize{10.6}{14.5}\selectfont\color{popink},
  fontlower=\popsemi\fontsize{10.6}{14.5}\selectfont\color{popink}}}
% Coche dessinée (le glyphe ✓ n'existe pas dans les polices du document)
\newcommand{\popcheck}[1]{\tikz[baseline=-0.5ex]\draw[#1,line width=1.6pt,line cap=round,line join=round](-0.13,0.0)--(-0.04,-0.09)--(0.13,0.1);}
\providecommand{\coche}{\popcheck{popgreen}}
\providecommand{\resultat}[1]{\par\smallskip\noindent\fcolorbox{popgreen}{popgreenL}{\parbox{\dimexpr\linewidth-2\fboxsep-2\fboxrule\relax}{\textbf{Résultat.} #1}}\par\smallskip}
\newcommand{\popboxhead}[3]{\poppill{#1}{#2}{white}\ifx\relax#3\relax\else\hspace{6pt}{\popxbold\fontsize{10.6}{12}\selectfont\color{popink}#3}\fi\par\vspace{7pt}}

\newtcolorbox{aretenir}[1][]{popbase,colframe=popgreen,before upper={\popboxhead{À retenir}{popgreen}{#1}}}
% Texte en allemand (document de lecture, dialogue, extrait) : cadre
% d'étude. Un titre optionnel (source, auteur) s'affiche dans l'en-tête.
\newtcolorbox{textebox}[1][]{popbase,colframe=popblue,before upper={\popboxhead{Text}{popblue}{#1}\begin{otherlanguage}{ngerman}},after upper={\end{otherlanguage}}}
% Marques ✓ / ✗ (forme correcte / fautive) souvent utilisées par les modèles
\providecommand{\cross}{\textcolor{poppink}{\ensuremath{\times}}}
\providecommand{\xmark}{\cross}\providecommand{\crossmark}{\cross}\providecommand{\cmark}{\ensuremath{\checkmark}}
% Ligne de réponse / justification à compléter (inventée par les modèles)
\providecommand{\justification}[1]{\par\noindent\textit{Justification~:} \dotfill\par}
% \textipa (package tipa non chargé) : la police ne couvre pas l'API -> simple passage
\providecommand{\textipa}[1]{#1}
% Soulignement (ulem non chargé) : \uline, \ul
\providecommand{\uline}[1]{\underline{#1}}\providecommand{\ul}[1]{\underline{#1}}
% \glossaire{\item ...} inventée par les modèles : liste de glossaire
\providecommand{\glossaire}[1]{\begin{itemize}#1\end{itemize}}
% \alert (beamer) inventée par les modèles : texte mis en évidence
\providecommand{\alert}[1]{\textbf{\textcolor{poppink}{#1}}}
% variantes ASCII d'ordinaux écrites en macro
\providecommand{\iere}{\textsuperscript{re}}\providecommand{\ieme}{\textsuperscript{e}}\providecommand{\ere}{\textsuperscript{re}}\providecommand{\eme}{\textsuperscript{e}}\providecommand{\er}{\textsuperscript{er}}\providecommand{\ie}{\textsuperscript{e}}
% Ordinaux français écrits en macro par les modèles : 1\ière, 3\ième
\newcommand{\ière}{\textsuperscript{re}}\newcommand{\ième}{\textsuperscript{e}}
% \exemple (inventée par les modèles) : étiquette « Exemple : »
\providecommand{\exemple}{\textbf{Exemple~:}\ }
% \square utilisable aussi hors mode math (cases à cocher)
\AtBeginDocument{\let\squareMath\square\renewcommand{\square}{\ensuremath{\squareMath}}}
% Mot ou phrase en allemand à l'intérieur d'un paragraphe français
\newcommand{\all}[1]{\ifmmode\text{\textit{#1}}\else\begingroup\selectlanguage{ngerman}\itshape #1\endgroup\fi}
\let\esp\all % alias toléré
\newtcolorbox{exemplebox}[1][]{popbase,colframe=poporange,before upper={\popboxhead{Exemple}{poporange}{#1}}}
\newtcolorbox{definition}[1][]{popbase,colframe=popblue,before upper={\popboxhead{Définition}{popblue}{#1}}}
\newtcolorbox{propriete}[1][]{popbase,colframe=poppurple,before upper={\popboxhead{Propriété}{poppurple}{#1}}}
\newtcolorbox{methode}[1][]{popbase,colframe=popgold,before upper={\popboxhead{Méthode}{popgold}{#1}}}
\newtcolorbox{attention}[1][]{popbase,colframe=poppink,before upper={\popboxhead{Attention}{poppink}{#1}}}
\newtcolorbox{experience}[1][]{popbase,colframe=popblue,colback=popblueL,before upper={\popboxhead{Expérience}{popblue}{#1}}}
% « Le savais-tu ? » : carte violette pleine (contexte, culture, Cameroun)
\newtcolorbox{savaistu}[1][]{popbase,colback=poppurple,colframe=popink,
  fontupper=\popsemi\fontsize{10.4}{14.2}\selectfont\color{white},
  before upper={\poppill{Le savais-tu\,?}{white}{poppurple}\ifx\relax#1\relax\else\hspace{6pt}{\popxbold\color{white}#1}\fi\par\vspace{7pt}}}
% Exercice résolu : énoncé en haut, \tcblower puis la solution sur fond crème
\newcounter{popexo}\newcounter{popexr}
\newtcolorbox{exoresolu}[1][]{popbase,colframe=poporange,colbacklower=popcream,
  segmentation style={popink,line width=1.2pt,dashed},
  before upper={\stepcounter{popexr}\popboxhead{Exercice résolu \thepopexr}{poporange}{#1}},
  before lower={\poppill{Solution}{popink}{popcream}\par\vspace{6pt}}}
% Exercice (non résolu en place) — la correction est dans la partie Corrigés
\newtcolorbox{exercice}[1][]{popbase,colframe=popink,
  before upper={\stepcounter{popexo}\popboxhead{Exercice \thepopexo}{popink}{#1}}}
% Corrigé
\newtcolorbox{corrige}[1][]{popbase,colframe=popgreen,colback=popgreenL,
  before upper={\popboxhead{Corrigé}{popgreen}{#1}}}
% Objectifs / prérequis / bilan
\newtcolorbox{objectifs}{popbase,colframe=popink,colback=poporangeL,
  before upper={\popboxhead{Objectifs de la leçon}{popink}{}}}
\newtcolorbox{prerequis}{popbase,colframe=popink,colback=popblueL,
  before upper={\popboxhead{Prérequis}{popblue}{}}}
\newtcolorbox{bilan}{popbase,colframe=popink,colback=white,boxrule=2.8pt,
  before upper={\popboxhead{Fiche bilan}{poporange}{L'essentiel à connaître pour l'examen}}}
% Figure encadrée avec légende
\newtcolorbox{popfigure}[1][]{enhanced,breakable=false,colback=white,colframe=popline,boxrule=1.4pt,arc=8pt,
  left=10pt,right=10pt,top=10pt,bottom=8pt,before skip=10pt,after skip=12pt,halign=center,
  lower separated=false,halign lower=center,
  fontlower=\popsemi\fontsize{9}{11.5}\selectfont\color{popmuted},
  #1}
\newcommand{\legende}[1]{\par\vspace{6pt}{\popsemi\fontsize{9}{11.5}\selectfont\color{popmuted}#1}}

% ============================================================
% Signature OnBuch+
% ============================================================
\newcommand{\poplogoinline}[1]{{\poplogo\fontsize{#1}{#1}\selectfont\color{popink}OnBuch\color{poporange}+}}
\newcommand{\signatureonbuch}{%
  \begin{tcolorbox}[enhanced,colback=popink,colframe=popink,boxrule=2.2pt,arc=10pt,
      drop shadow=poporange,left=14pt,right=14pt,top=10pt,bottom=10pt,before skip=0pt,after skip=0pt]
    \tikz[baseline=-0.7ex]\node[circle,fill=popgreen,draw=popcream,line width=1.3pt,minimum size=22pt,inner sep=0]{\popcheck{white}};%
    \hspace{9pt}\begin{minipage}[c]{0.62\linewidth}
      {\popxbold\fontsize{12}{13}\selectfont\color{popcream}Signé\ \textbullet\ {\poplogo OnBuch\color{poporange}+}}\par\vspace{2pt}
      {\raggedright\popsemi\fontsize{8}{10.5}\selectfont\color{popline}\MakeUppercase{Équipe pédagogique OnBuch+}\\\MakeUppercase{Conforme au programme officiel MINESEC}\par}
    \end{minipage}\hfill
    {\poplogo\fontsize{22}{22}\selectfont\color{popcream}OnBuch\color{poporange}+}
  \end{tcolorbox}%
}

% ============================================================
% Page de garde
% #1 titre · #2 matière · #3 niveau · #4 module · #5 n° leçon · #6 durée ·
% #7 description / objectifs (texte ou itemize)
% ============================================================
\newcommand{\pagegarde}[7]{%
  \thispagestyle{empty}
  \newgeometry{a4paper,margin=1.7cm,top=1.6cm,bottom=1.4cm}%
  \begin{tikzpicture}[remember picture,overlay]
    \fill[poporange!18] ([xshift=-3.2cm,yshift=-3.6cm]current page.north east) circle (4.6cm);
    \fill[poppurple!14] ([xshift=2.4cm,yshift=7.6cm]current page.south west) circle (3.8cm);
  \end{tikzpicture}%
  {\poplogo\fontsize{30}{30}\selectfont\color{popink}OnBuch\color{poporange}+}\hfill
  \raisebox{0.6ex}{\poppill{Cours signé \textbullet\ Premium}{popink}{popcream}}\par
  \vspace{1pt}\popsquiggle{poppurple}\par
  \vspace{8pt}
  \begin{tcolorbox}[enhanced,colback=poporange,colframe=popink,boxrule=2.8pt,arc=13pt,
      drop shadow=popink,left=18pt,right=18pt,top=12pt,bottom=13pt,after skip=10pt]
    \poppill{#2 \textbullet\ #3}{popink}{popcream}\hspace{5pt}\poppill{#4}{white}{popink}\par\vspace{8pt}
    {\popdisplay\fontsize{14}{14}\selectfont\color{popink}LEÇON #5}\par\vspace{4pt}
    {\raggedright\hyphenpenalty=10000\exhyphenpenalty=10000\popdisplay\fontsize{25}{27}\selectfont\color{popcream}\MakeUppercase{#1}\par}
    \vspace{8pt}
    {\popxbold\fontsize{9.5}{9.5}\selectfont\color{popink}\MakeUppercase{Durée conseillée : #6}}
  \end{tcolorbox}
  \begin{tcolorbox}[enhanced,colback=white,colframe=popink,boxrule=2.2pt,arc=10pt,
      drop shadow=popink,left=16pt,right=16pt,top=10pt,bottom=10pt,after skip=8pt]
    \poppill{Dans cette leçon}{poppurple}{white}\par\vspace{5pt}
    \popsemi\fontsize{9.3}{12.6}\selectfont\color{popink}#7
  \end{tcolorbox}
  \noindent\begin{minipage}[t]{0.487\linewidth}
    \begin{tcolorbox}[enhanced,colback=popgreen,colframe=popink,boxrule=2.2pt,arc=10pt,
        drop shadow=popink,left=11pt,right=11pt,top=10pt,bottom=10pt,height=3.1cm,valign=center]
      \tcbox[colback=white,colframe=popink,boxrule=1.3pt,arc=5pt,boxsep=0pt,nobeforeafter,
        left=3pt,right=3pt,top=3pt,bottom=3pt,valign=center]{\qrcode[height=1.9cm]{https://play.google.com/store/apps/details?id=com.luvvix.onbuch&hl=fr}}%
      \hspace{8pt}\begin{minipage}[c]{0.5\linewidth}
        {\popdisplay\fontsize{11}{12.5}\selectfont\color{white}\MakeUppercase{Télécharge OnBuch}}\par\vspace{4pt}
        {\raggedright\popsemi\fontsize{8.4}{11}\selectfont\color{white}Gratuit \textbullet\ Google Play\\Tuteur IA, annales, concours, résultats\par}
      \end{minipage}
    \end{tcolorbox}
  \end{minipage}\hfill
  \begin{minipage}[t]{0.487\linewidth}
    \begin{tcolorbox}[enhanced,colback=poppurple,colframe=popink,boxrule=2.2pt,arc=10pt,
        drop shadow=popink,left=11pt,right=11pt,top=10pt,bottom=10pt,height=3.1cm,valign=center]
      \tikz[baseline=-0.7ex]\node[circle,fill=white,draw=popink,line width=1.3pt,minimum size=1.6cm,inner sep=0]{\popdisplay\fontsize{24}{24}\selectfont\color{poppurple}+};%
      \hspace{8pt}\begin{minipage}[c]{0.6\linewidth}
        {\popdisplay\fontsize{11}{12.5}\selectfont\color{white}\MakeUppercase{Passe à OnBuch+}}\par\vspace{4pt}
        {\raggedright\popsemi\fontsize{8.4}{11}\selectfont\color{white}Tous les cours signés, Tuteur IA illimité, fiches PDF — onglet OnBuch+ de l'app\par}
      \end{minipage}
    \end{tcolorbox}
  \end{minipage}
  \vspace{8pt}
  \popcontenu
  \vfill
  \signatureonbuch
  \restoregeometry
  \newpage
}

% Rangée « Ce cours contient » (page de garde)
\newcommand{\poptile}[3]{% #1 couleur #2 gros texte #3 légende
  \begin{tcolorbox}[enhanced,colback=#1,colframe=popink,boxrule=2pt,arc=9pt,shadow={2.5pt}{-2.5pt}{0pt}{popink},
      left=6pt,right=6pt,top=9pt,bottom=9pt,halign=center,valign=center,height=1.75cm,width=\linewidth,nobeforeafter]
    {\popdisplay\fontsize{12.5}{13}\selectfont\color{white}\MakeUppercase{#2}}\par\vspace{3pt}
    {\centering\popsemi\fontsize{7.8}{9.5}\selectfont\color{white}#3\par}
  \end{tcolorbox}}
\newcommand{\popcontenu}{%
  \par\noindent{\popxboldls\fontsize{7.5}{7.5}\selectfont\color{popmuted}CE COURS SIGNÉ CONTIENT}\par\vspace{7pt}
  \noindent\begin{minipage}[t]{0.235\linewidth}\poptile{popblue}{Cours}{complet, illustré, pas à pas}\end{minipage}\hfill
  \begin{minipage}[t]{0.235\linewidth}\poptile{poporange}{Méthodes}{et exercices résolus}\end{minipage}\hfill
  \begin{minipage}[t]{0.235\linewidth}\poptile{poppink}{Exercices}{type examen + corrigés}\end{minipage}\hfill
  \begin{minipage}[t]{0.235\linewidth}\poptile{popgreen}{Fiche bilan}{l'essentiel à retenir}\end{minipage}\par}

% Sommaire
\newtcolorbox{sommaire}{enhanced,colback=white,colframe=popink,boxrule=2.2pt,arc=10pt,
  drop shadow=popink,left=18pt,right=18pt,top=13pt,bottom=13pt,before skip=6pt,after skip=20pt,
  before upper={\poppill{Sommaire}{poppurple}{white}\par\vspace{9pt}\popsemi\fontsize{10.6}{14.5}\selectfont\color{popink}}}

% Signature de fin de cours — poussée tout en bas de la dernière page
\newcommand{\findecours}{%
  \par\vfill\nopagebreak
  \begin{center}\popsquiggle{poporange}\end{center}\vspace{4pt}
  \signatureonbuch
}
\AtBeginDocument{\def\llbracket{\mathopen{[\mkern-3.5mu[}}\def\rrbracket{\mathclose{]\mkern-3.5mu]}}} % intervalles d'entiers (glyphe ⟦ absent de la police)
% Glyphes absents de Fira Math : redéfinis à partir de glyphes disponibles
\AtBeginDocument{%
  \def\setminus{\mathbin{\backslash}}\let\smallsetminus\setminus
  \def\vdots{\mathord{\vbox{\baselineskip3.2pt\lineskiplimit0pt\kern4pt\hbox{.}\hbox{.}\hbox{.}}}}%
  \def\ddots{\mathinner{\mkern1mu\raise6.5pt\hbox{.}\mkern2mu\raise3.6pt\hbox{.}\mkern2mu\raise0.7pt\hbox{.}\mkern1mu}}%
  \def\triangle{\mathord{\scalebox{0.9}{$\symup{\Delta}$}}}%
  \def\nabla{\mathord{\raisebox{\depth}{\scalebox{1}[-1]{$\symup{\Delta}$}}}}%
  \def\checkmark{\popcheck{popdark}}%
  \def\star{\mathbin{\ast}}\def\bigstar{\mathord{\ast}}%
  \def\diamond{\mathbin{\scalebox{0.6}{$\Diamond$}}}%
  \def\bigcup{\mathop{\mathchoice{\raisebox{-0.3em}{\scalebox{1.7}{$\cup$}}}{\scalebox{1.25}{$\cup$}}{\cup}{\cup}}\displaylimits}%
  \def\bigcap{\mathop{\mathchoice{\raisebox{-0.3em}{\scalebox{1.7}{$\cap$}}}{\scalebox{1.25}{$\cap$}}{\cap}{\cap}}\displaylimits}%
  \def\lesseqqgtr{\mathrel{\lessgtr}}\def\bigoplus{\mathop{\oplus}\displaylimits}%
}
\providecommand{\ssi}{si et seulement si} % abréviation fréquente
\AtBeginDocument{\providecommand{\wideparen}{\overparen}} % arc de cercle

%% FIGFIX-BEGIN v5 — correctifs globaux des figures (docs/onbuch-generation/tools/figfix)
\usepackage{adjustbox}\usepackage{etoolbox}\tcbuselibrary{raster}
% toute figure TikZ/pgfplots plus large que la ligne est réduite à la largeur disponible
\BeforeBeginEnvironment{tikzpicture}{\begin{adjustbox}{max width=\linewidth}}
\AfterEndEnvironment{tikzpicture}{\end{adjustbox}}
% légendes pgfplots lisibles par-dessus les courbes
\pgfplotsset{every axis/.append style={legend style={fill=white,fill opacity=0.92,text opacity=1,draw=black!15,inner sep=3pt}}}
% étiquettes d'arêtes (syntaxe edge["..."]) avec fond blanc
\tikzset{every edge quotes/.append style={fill=white,inner sep=1.2pt}}
%% FIGFIX-END
\begin{document}

\pagegarde{Production et traduction de textes liés à la vie en famille et en société}{Allemand}{1ère A}{Chapitre 1 : Vie familiale et sociale}{ALA04}{2 h}{Cette leçon mixte entraîne les élèves à comprendre, produire et traduire des textes courts (lettre, description, dialogue) sur la vie familiale et sociale. Elle consolide les révisions de grammaire essentielles : présent, parfait et ordre des mots, à travers le thème et la version.
\vspace{4pt}
\begin{multicols}{2}\small
\begin{itemize}
\item À la fin de cette leçon, je sais comprendre un texte allemand simple sur la vie en famille et en société.
\item À la fin de cette leçon, je sais utiliser le lexique de la famille, des fêtes et de la vie sociale dans mes productions.
\item À la fin de cette leçon, je sais conjuguer correctement au présent les verbes réguliers, irréguliers et modaux.
\item À la fin de cette leçon, je sais raconter un événement familial au parfait avec le bon auxiliaire (haben/sein) et le bon participe.
\item À la fin de cette leçon, je sais placer correctement le verbe en position 2 dans la phrase déclarative et gérer l'inversion et les verbes à particule.
\item À la fin de cette leçon, je sais rédiger une courte lettre personnelle, une description de famille et un dialogue simple.
\end{itemize}\end{multicols}}

\begin{sommaire}
\begin{multicols}{2}
\begin{enumerate}
\item Texte support : Ein Familienfest in München
\item Lexique thématique : famille et vie sociale
\item Grammaire 1 : révision du présent
\item Grammaire 2 : le parfait (Perfekt)
\item Grammaire 3 : l'ordre des mots
\item Production : lettre, description, dialogue
\item Traduction : thème et version
\item Activité d'intégration
\item Méthodes et exercices résolus
\item Exercices
\item Corrigés des exercices
\item Fiche bilan
\end{enumerate}
\end{multicols}
\end{sommaire}

\begin{objectifs}
\begin{itemize}
\item À la fin de cette leçon, je sais comprendre un texte allemand simple sur la vie en famille et en société.
\item À la fin de cette leçon, je sais utiliser le lexique de la famille, des fêtes et de la vie sociale dans mes productions.
\item À la fin de cette leçon, je sais conjuguer correctement au présent les verbes réguliers, irréguliers et modaux.
\item À la fin de cette leçon, je sais raconter un événement familial au parfait avec le bon auxiliaire (haben/sein) et le bon participe.
\item À la fin de cette leçon, je sais placer correctement le verbe en position 2 dans la phrase déclarative et gérer l'inversion et les verbes à particule.
\item À la fin de cette leçon, je sais rédiger une courte lettre personnelle, une description de famille et un dialogue simple.
\item À la fin de cette leçon, je sais traduire de courtes phrases et un texte simple du français vers l'allemand (thème) et de l'allemand vers le français (version).
\item À la fin de cette leçon, je sais relire et corriger ma production : orthographe, majuscules, ordre des mots et cas.
\end{itemize}
\end{objectifs}

\begin{prerequis}
\begin{itemize}
\item Conjugaison au présent des verbes réguliers et des verbes irréguliers fréquents (sein, haben, werden, modaux) vue en Seconde.
\item Formation du parfait (Perfekt) avec haben et sein vue en Seconde.
\item Ordre des mots de base : verbe conjugué en 2e position, verbe à l'infinitif ou participe en fin de phrase.
\item Vocabulaire de base de la famille (die Familie, die Eltern, die Geschwister...) acquis en Seconde.
\item Déclinaisons de base des articles et des cas (nominatif, accusatif, datif) vues en début d'année.
\end{itemize}
\end{prerequis}

\newpage

\coursec{Texte support : Ein Familienfest in München}

\begin{center}
\emph{« Paul, élève de Première à Yaoundé, vient de recevoir un e-mail de sa correspondante allemande Lena : elle lui raconte sa vie de famille à Munich et lui pose des questions sur la sienne. Pour lui répondre, Paul doit savoir décrire sa famille, raconter un événement familial au passé et rédiger une courte lettre en allemand correct. Il doit aussi traduire un petit texte que Lena a joint à son message. C'est exactement ce que cette leçon va vous apprendre : produire et traduire des textes sur la vie en famille et en société, en réinvestissant le présent, le parfait et l'ordre des mots. »}
\end{center}

\textbf{Problématique :} comment lire, comprendre et exploiter un texte allemand qui raconte un événement familial, afin de pouvoir ensuite en produire un soi-même et le traduire ?

\courssub{Lecture globale : le texte de Lena}

Lisons d'abord le texte que Lena a envoyé à Paul. Lisez-le une première fois sans dictionnaire : essayez seulement de saisir le sens général.

\begin{textebox}[Ein Familienfest in München]
Letztes Wochenende haben wir den 70. Geburtstag meiner Großmutter gefeiert. Meine Eltern haben alle Verwandten eingeladen: Onkel, Tanten und Cousins. Zuerst haben wir zusammen gegessen, dann hat meine Mutter einen großen Kuchen gebracht. Meine Großmutter hat viele Geschenke bekommen und hat sich sehr gefreut. Danach haben wir gesungen und getanzt. Schließlich hat mein Vater Fotos gemacht. Am Abend sind wir spät nach Hause gefahren. Es war ein schönes Fest!
\end{textebox}

\textbf{Traduction :} Le week-end dernier, nous avons fêté les 70 ans de ma grand-mère. Mes parents ont invité tous les membres de la famille : oncles, tantes et cousins. D'abord, nous avons mangé ensemble, puis ma mère a apporté un grand gâteau. Ma grand-mère a reçu beaucoup de cadeaux et a été très contente. Ensuite, nous avons chanté et dansé. Finalement, mon père a pris des photos. Le soir, nous sommes rentrés tard à la maison. C'était une belle fête!

\begin{definition}[Le texte narratif (die Erzählung)]
Un \cle{texte narratif} (en allemand \all{die Erzählung}) raconte des événements réels ou imaginaires dans un ordre chronologique. En allemand, le récit au passé se fait le plus souvent au \cle{parfait} (\all{Perfekt}) dans la langue parlée et les lettres personnelles. Le texte est organisé par des \cle{connecteurs temporels} comme \all{zuerst} (d'abord), \all{dann} (puis), \all{danach} (ensuite), \all{schließlich} (finalement).
\end{definition}

\courssub{Glossaire : les mots nouveaux}

Avant la deuxième lecture, voici le vocabulaire essentiel du texte. Apprenez chaque nom \textbf{avec son article et son pluriel} : c'est indispensable en allemand.

\begin{center}
\begin{tabularx}{\linewidth}{|X|X|X|}
\hline
\rowcolor{popblueL} \textbf{Deutsch} & \textbf{Français} & \textbf{Exemple} \\
\hline
das Fest, -e & la fête & \all{Es war ein schönes Fest.} \\
\hline
die Feier, -n & la fête, la célébration & \all{Die Feier beginnt um 18 Uhr.} \\
\hline
die Großeltern (Pl.) & les grands-parents & \all{Meine Großeltern wohnen in München.} \\
\hline
sich freuen & se réjouir, être content & \all{Sie hat sich sehr gefreut.} \\
\hline
einladen (lädt ein, hat eingeladen) & inviter & \all{Wir haben alle Verwandten eingeladen.} \\
\hline
das Geschenk, -e & le cadeau & \all{Sie hat viele Geschenke bekommen.} \\
\hline
die Verwandten (Pl.) & les membres de la famille, les parents & \all{Alle Verwandten kommen.} \\
\hline
zubereiten (hat zubereitet) & préparer (un plat) & \all{Lena hat einen Salat zubereitet.} \\
\hline
\end{tabularx}
\end{center}

\begin{attention}[Faux amis et pièges]
\begin{itemize}
\item \all{die Verwandten} ne signifie pas « les voisins » mais « les membres de la famille, les parents ».
\item \all{bekommen} ne veut pas dire « devenir » mais « recevoir » : \all{Sie hat viele Geschenke bekommen.} = Elle a reçu beaucoup de cadeaux. (« Devenir » se dit \all{werden}.)
\item Ne traduisez jamais mot à mot : \all{sich freuen} est un \cle{verbe pronominal} : \all{ich freue mich, du freust dich, sie freut sich...} Le pronom réfléchi fait partie du verbe, comme « se réjouir » en français.
\item \all{einladen} est un \cle{verbe à particule séparable} : au présent \all{ich lade ein}, au parfait \all{hat eingeladen} (particule \all{ein-} + \all{ge-} + radical).
\end{itemize}
\end{attention}

\begin{methode}[Lire un texte en 3 étapes]
\begin{enumerate}
\item \textbf{Lecture globale} : lire le texte entier une fois, sans s'arrêter sur les mots inconnus, pour saisir le sens général (de quoi parle-t-on ?).
\item \textbf{Lecture analytique} : relire en répondant aux questions : qui ? quoi ? où ? quand ? pourquoi ? Souligner les verbes et les connecteurs.
\item \textbf{Relevé} : noter le lexique utile et les structures (temps, ordre des mots, connecteurs) que l'on pourra réutiliser dans sa propre production.
\end{enumerate}
\end{methode}

\courssub{Première lecture : thème, personnages, lieu et moment}

Appliquons l'étape 2 de la méthode au texte de Lena.

\begin{itemize}
\item \textbf{Le thème :} une fête de famille — l'anniversaire de la grand-mère (\all{der 70. Geburtstag}).
\item \textbf{Les personnages :} Lena (la narratrice), sa grand-mère, ses parents, et les \all{Verwandten} (oncles, tantes, cousins).
\item \textbf{Le lieu :} à Munich, chez la famille de Lena (le titre l'indique : \all{in München}).
\item \textbf{Le moment :} le week-end dernier (\all{letztes Wochenende}) — donc un événement \textbf{passé}.
\end{itemize}

\begin{exemplebox}[Repérage]
\all{Letztes Wochenende haben wir den 70. Geburtstag meiner Großmutter gefeiert.}

Cette première phrase donne à elle seule le moment (\all{letztes Wochenende}), l'événement (\all{Geburtstag ... gefeiert}) et le personnage central (\all{meine Großmutter}). C'est une excellente phrase d'introduction à imiter dans vos propres récits : « Le week-end dernier, nous avons fêté les 70 ans de ma grand-mère. »
\end{exemplebox}

\courssub{Deuxième lecture : les temps du récit}

Soulignons maintenant \textbf{tous} les verbes conjugués du texte et justifions le choix de chaque temps.

\begin{center}
\begin{tabularx}{\linewidth}{|X|X|X|}
\hline
\rowcolor{popblueL} \textbf{Verbe} & \textbf{Temps} & \textbf{Justification} \\
\hline
haben ... gefeiert & Perfekt & événement passé, récit oral/épistolaire \\
\hline
haben ... eingeladen & Perfekt & action passée achevée \\
\hline
haben ... gegessen & Perfekt & action passée achevée \\
\hline
hat ... gebracht & Perfekt & action passée achevée \\
\hline
hat ... bekommen / hat sich gefreut & Perfekt & action et sentiment passés \\
\hline
haben ... gesungen / haben ... getanzt & Perfekt & actions successives \\
\hline
hat ... gemacht & Perfekt & action achevée \\
\hline
sind ... gefahren & Perfekt (avec \all{sein}) & verbe de déplacement \\
\hline
war & Präteritum de \all{sein} & bilan final, impression \\
\hline
\end{tabularx}
\end{center}

\begin{aretenir}[Présent ou parfait ?]
Dans un récit d'événement passé (lettre, e-mail, conversation), l'allemand utilise le \cle{Perfekt} : auxiliaire \all{haben} ou \all{sein} au présent + participe passé \textbf{à la fin de la phrase}. Le \cle{présent} sert à décrire des faits actuels ou habituels (par exemple pour décrire sa famille : \all{Meine Familie wohnt in Yaoundé.}). Les verbes \all{sein} et \all{haben} apparaissent souvent au Präteritum (\all{war}, \all{hatte}) même dans la langue parlée.
\end{aretenir}

\courssub{Questions de compréhension}

\textbf{Auf Deutsch (répondez en phrases complètes) :}

\begin{enumerate}
\item Wer feiert? (Qui fête ?)
\item Warum feiert die Familie? (Pourquoi la famille fête-t-elle ?)
\item Wer kommt zum Fest? (Qui vient à la fête ?)
\item Was hat die Mutter gebracht? (Qu'est-ce que la mère a apporté ?)
\item Wie war das Fest? (Comment était la fête ?)
\end{enumerate}

\textbf{En français (compréhension fine) :}

\begin{enumerate}
\item Quel âge la grand-mère de Lena a-t-elle ?
\item Qui a invité les membres de la famille ?
\item Que reçoit la grand-mère et comment réagit-elle ?
\item Quelles activités la famille fait-elle après le repas ?
\item Comment la famille rentre-t-elle, et quand ?
\item Quel sentiment général Lena exprime-t-elle à la fin du texte ?
\end{enumerate}

\begin{exoresolu}[Compréhension du texte]
Énoncé : répondez en allemand, par une phrase complète, à la question : \all{Warum feiert die Familie?}
\tcblower
Solution détaillée : la question porte sur la \textbf{cause} (\all{warum} = pourquoi). On repère dans la première phrase : \all{den 70. Geburtstag meiner Großmutter}. On construit une réponse avec \all{weil} (parce que), en plaçant le verbe conjugué à la fin de la subordonnée :

\all{Die Familie feiert, weil die Großmutter 70 Jahre alt wird.} — La famille fête parce que la grand-mère a 70 ans.

Autre réponse possible sans subordonnée : \all{Die Familie feiert den 70. Geburtstag der Großmutter.}
\end{exoresolu}

\courssub{Les connecteurs temporels}

Relevons les mots qui organisent la chronologie du récit. Ce sont des outils précieux pour vos propres productions.

\begin{center}
\begin{tabularx}{\linewidth}{|X|X|X|}
\hline
\rowcolor{popblueL} \textbf{Connecteur} & \textbf{Français} & \textbf{Fonction} \\
\hline
zuerst & d'abord & première étape \\
\hline
dann & puis, ensuite & étape suivante \\
\hline
danach & après cela, ensuite & étape suivante \\
\hline
schließlich & finalement & dernière étape \\
\hline
am Abend & le soir & repère du moment de la journée \\
\hline
\end{tabularx}
\end{center}

\begin{attention}[L'ordre des mots après un connecteur]
Quand la phrase \textbf{commence} par un connecteur temporel (ou tout autre complément), le verbe conjugué reste en \textbf{deuxième position} et le sujet passe après le verbe (inversion) :

\all{Am Abend sind wir spät nach Hause gefahren.} — Le soir, nous sommes rentrés tard à la maison.

On n'écrit jamais \all{Am Abend wir sind...} : c'est l'erreur la plus fréquente des francophones. Règle d'or : \textbf{une seule idée (un seul constituant) avant le verbe}, et le verbe toujours en position 2.
\end{attention}

\begin{exemplebox}[Exemples traduits]
\all{Meine Großmutter hat sich sehr gefreut.} — Ma grand-mère a été très contente (litt. « s'est beaucoup réjouie »).

\all{Am Abend sind wir spät nach Hause gefahren.} — Le soir, nous sommes rentrés tard à la maison. (Auxiliaire \all{sein} car \all{fahren} exprime ici un déplacement.)

\all{Zuerst haben wir zusammen gegessen, dann hat meine Mutter einen Kuchen gebracht.} — D'abord nous avons mangé ensemble, puis ma mère a apporté un gâteau.
\end{exemplebox}

\courssub{La structure du texte narratif}

Observons comment le texte de Lena est construit : il suit le plan classique du récit, en trois parties.

\begin{itemize}
\item \textbf{Introduction (Situation)} : l'occasion est présentée — \all{Letztes Wochenende haben wir den 70. Geburtstag ... gefeiert.}
\item \textbf{Développement (Événements)} : le déroulement de la fête, étape par étape, grâce aux connecteurs — manger, le gâteau, les cadeaux, chanter, danser, les photos.
\item \textbf{Conclusion (Bilan)} : le retour et le sentiment final — \all{Es war ein schönes Fest!}
\end{itemize}

\begin{popfigure}
\begin{tikzpicture}[
  node distance=0.9cm,
  box/.style={draw=popblue, fill=popblueL, rounded corners, thick, align=center, text width=4.2cm, minimum height=1.1cm},
  arr/.style={-{Stealth[length=3mm]}, thick, poporange}
]
\node[box, align=center] (a){\textbf{Introduction}\\Situation initiale\\(occasion, lieu, moment)};
\node[box, below=of a, align=center] (b){\textbf{Développement}\\Événements\\(zuerst, dann, danach, schließlich)};
\node[box, below=of b, align=center] (c){\textbf{Conclusion}\\Bilan et sentiment\\(Es war ...!)};
\draw[arr] (a) -- (b);
\draw[arr] (b) -- (c);
\end{tikzpicture}
\legende{Organigramme de la structure d'un texte narratif : Situation initiale, Événements, Bilan.}
\end{popfigure}

\begin{popfigure}
\begin{tikzpicture}[
  grow=down,
  level distance=1.6cm,
  level 1/.style={sibling distance=4.5cm},
  level 2/.style={sibling distance=2.6cm},
  every node/.style={draw=popgreen, fill=popgreenL, rounded corners, thick, align=center, text width=2.6cm},
  edge from parent/.style={draw=popdark, thick}
]
\node[align=center]{\textbf{die Großeltern}\\(grands-parents)}
  child { node[align=center]{\textbf{der Vater}\\(le père)}
    child { node[align=center]{\textbf{Lena}\\(la narratrice)} }
    child { node[align=center]{\textbf{der Bruder}\\(le frère)} }
  }
  child { node[align=center]{\textbf{die Tante}\\(la tante)}
    child { node[align=center]{\textbf{die Cousins}\\(les cousins)} }
  };
\end{tikzpicture}
\legende{Arbre généalogique simplifié de la famille de Lena : Großeltern, Eltern, Kinder.}
\end{popfigure}

\begin{savaistu}[Les fêtes de famille en Allemagne]
En Allemagne, les anniversaires « ronds » (70., 80. Geburtstag) sont de grandes occasions de réunion familiale. On offre des fleurs et des cadeaux, on partage un \all{Kuchen} (gâteau) l'après-midi lors du traditionnel \all{Kaffee und Kuchen}. Au Cameroun aussi, les anniversaires et les fêtes de famille rassemblent parents, oncles, tantes et cousins autour d'un grand repas : un beau point commun entre les deux cultures, idéal pour vos sujets de production écrite!
\end{savaistu}

\begin{exercice}[À vous de jouer]
\begin{enumerate}
\item Recopiez le texte de Lena en soulignant tous les verbes au Perfekt et en entourant les connecteurs temporels.
\item Traduisez en français : \all{Zuerst haben wir zusammen gegessen, dann hat meine Mutter einen großen Kuchen gebracht.}
\item Répondez en allemand : \all{Wie war das Fest?}
\item Conjuguez \all{sich freuen} au présent aux trois personnes du singulier.
\end{enumerate}
\end{exercice}

\begin{corrige}[Exercice : À vous de jouer]
\begin{enumerate}
\item Verbes au Perfekt : haben ... gefeiert, haben ... eingeladen, haben ... gegessen, hat ... gebracht, hat ... bekommen, hat sich ... gefreut, haben ... gesungen, haben ... getanzt, hat ... gemacht, sind ... gefahren. Connecteurs : zuerst, dann, danach, schließlich, am Abend.
\item « D'abord nous avons mangé ensemble, puis ma mère a apporté un grand gâteau. »
\item \all{Das Fest war sehr schön.} (La fête était très belle.)
\item \all{ich freue mich, du freust dich, er/sie/es freut sich.}
\end{enumerate}
\end{corrige}

\begin{aretenir}[Ce qu'il faut retenir de cette section]
\begin{itemize}
\item Un texte narratif se lit en trois étapes : sens général, questions analytiques (qui ? quoi ? où ? quand ?), relevé du lexique et des structures.
\item Le récit au passé en allemand utilise le \cle{Perfekt} ; le verbe conjugué est en position 2, le participe passé en fin de phrase.
\item Les connecteurs \all{zuerst, dann, danach, schließlich, am Abend} organisent la chronologie ; placés en tête de phrase, ils provoquent l'inversion du sujet.
\item La structure Situation initiale → Événements → Bilan servira de modèle pour votre propre production (lettre, description, dialogue) dans la suite de la leçon.
\end{itemize}
\end{aretenir}

\coursec{Lexique thématique : famille et vie sociale}

Dans la section précédente, nous avons lu le texte \emph{Ein Familienfest in München} : la famille Schneider se réunit pour l'anniversaire du grand-père. Ce texte contenait déjà de nombreux mots de la famille et de la fête : \all{der Großvater}, \all{die Enkel}, \all{das Geschenk}, \all{die Feier}. Il est maintenant temps de \cle{structurer ce vocabulaire} de manière systématique. Un bon lexique s'apprend en \emph{champs thématiques}, jamais en liste désordonnée : c'est la condition pour produire ensuite une lettre, une description ou un dialogue sans chercher ses mots.

\begin{methode}[Comment apprendre un nom allemand ?]
\begin{enumerate}
\item Apprends toujours le nom avec son \cle{article} : \all{der Bruder}, jamais « Bruder » seul.
\item Apprends en même temps son \cle{pluriel} : \all{der Bruder, die Brüder}.
\item Note le nom dans une phrase complète : \all{Mein Bruder wohnt in Douala.}
\item Regroupe les mots par famille de sens (membres, fêtes, sentiments...).
\end{enumerate}
\end{methode}

\courssub{Champ 1 : les membres de la famille}

Observons d'abord quelques phrases tirées de notre texte support et de la vie quotidienne :

\all{Mein Vater arbeitet in Yaoundé, meine Mutter ist Lehrerin.} « Mon père travaille à Yaoundé, ma mère est enseignante. »

\all{Ich habe einen Bruder und zwei Schwestern.} « J'ai un frère et deux sœurs. »

\all{Meine Eltern besuchen oft meine Großeltern.} « Mes parents rendent souvent visite à mes grands-parents. »

Voici le tableau complet à mémoriser, avec article et pluriel :

\begin{center}
\begin{tabularx}{\linewidth}{|l|l|X|}
\hline
\rowcolor{popblueL} \textbf{Deutsch} & \textbf{Pluriel} & \textbf{Français} \\
\hline
der Vater & die Väter & le père \\
\hline
die Mutter & die Mütter & la mère \\
\hline
die Eltern & (pluriel seul) & les parents \\
\hline
der Bruder & die Brüder & le frère \\
\hline
die Schwester & die Schwestern & la sœur \\
\hline
die Geschwister & (pluriel seul) & les frères et sœurs \\
\hline
der Sohn & die Söhne & le fils \\
\hline
die Tochter & die Töchter & la fille \\
\hline
das Enkelkind & die Enkelkinder & le petit-enfant \\
\hline
der Onkel & die Onkel & l'oncle \\
\hline
die Tante & die Tanten & la tante \\
\hline
der Cousin & die Cousins & le cousin \\
\hline
die Cousine & die Cousinen & la cousine \\
\hline
der Ehemann & die Ehemänner & le mari (époux) \\
\hline
die Ehefrau & die Ehefrauen & la femme (épouse) \\
\hline
\end{tabularx}
\end{center}

\begin{attention}[die Eltern et die Geschwister : pluriels obligatoires]
\all{die Eltern} (les parents) et \all{die Geschwister} (les frères et sœurs) n'existent qu'au pluriel. On dit donc : \all{Meine Eltern \textbf{wohnen} in Bamenda.} (verbe au pluriel !) et non \all{*meine Eltern wohnt}. Pour dire « un parent », on précise : \all{der Vater} ou \all{die Mutter}.
\end{attention}

\begin{exemplebox}[La famille en phrases]
\all{Meine Schwester heiratet nächsten Monat.} « Ma sœur se marie le mois prochain. »

\all{Seine Tochter lernt Deutsch in der Schule.} « Sa fille apprend l'allemand à l'école. »

\all{Unsere Enkelkinder besuchen uns jeden Sonntag.} « Nos petits-enfants nous rendent visite chaque dimanche. »

\all{Ihr Onkel lebt in Deutschland, in München.} « Leur oncle vit en Allemagne, à Munich. »
\end{exemplebox}

\courssub{Champ 2 : la vie quotidienne en famille}

\begin{definition}[Verbes de la vie quotidienne]
\all{zusammen leben} : vivre ensemble ; \all{sich verstehen} : s'entendre ; \all{sich streiten} : se disputer ; \all{helfen (+ datif)} : aider ; \all{sich kümmern um (+ accusatif)} : s'occuper de ; \all{die Hausarbeit, -en} : les tâches ménagères ; \all{kochen} : cuisiner ; \all{einkaufen} : faire les courses.
\end{definition}

Remarque la construction de deux verbes réfléchis :

\all{Wir verstehen uns gut.} « Nous nous entendons bien. »

\all{Die Kinder streiten sich oft, aber sie helfen sich auch.} « Les enfants se disputent souvent, mais ils s'entraident aussi. »

\all{Sie kümmert sich um ihre kranke Großmutter.} « Elle s'occupe de sa grand-mère malade. »

\all{Am Samstag kauft die ganze Familie auf dem Markt ein.} « Le samedi, toute la famille fait ses courses au marché. »

\begin{propriete}[Verbes construits avec le DATIF]
Certains verbes allemands se construisent avec le \cle{datif}, alors que leurs équivalents français prennent un complément direct. C'est le cas de : \all{helfen} (aider), \all{danken} (remercier), \all{gratulieren} (féliciter), \all{antworten} (répondre).

\all{Ich helfe \textbf{meiner Mutter}.} « J'aide ma mère. »

\all{Wir danken \textbf{unseren Nachbarn}.} « Nous remercions nos voisins. »

\all{Er antwortet \textbf{dem Lehrer}.} « Il répond au professeur. »

Règle : après ces verbes, la personne est au datif (\all{dem}, \all{der}, \all{den} + nom au pluriel avec -n).
\end{propriete}

\courssub{Champ 3 : fêtes et événements sociaux}

\begin{textebox}[Tableau lexique : les fêtes]
\all{der Geburtstag, -e} : l'anniversaire \\
\all{die Hochzeit, -en} : le mariage \\
\all{die Taufe, -n} : le baptême \\
\all{die Beerdigung, -en} : l'enterrement \\
\all{die Feier, -n} : la fête, la célébration \\
\all{das Fest, -e} : la fête, le festival \\
\all{einladen (lädt ein, hat eingeladen)} : inviter \\
\all{die Einladung, -en} : l'invitation \\
\all{das Geschenk, -e} : le cadeau \\
\all{gratulieren (+ datif)} : féliciter
\end{textebox}

\begin{exemplebox}[Fêtes en contexte]
\all{Er gratuliert seiner Schwester zum Geburtstag.} « Il félicite sa sœur pour son anniversaire. »

\all{Wir laden unsere Freunde zur Hochzeit ein.} « Nous invitons nos amis au mariage. »

\all{Danke für die Einladung zu deinem Fest!} « Merci pour l'invitation à ta fête ! »

\all{Die Kinder bekommen viele Geschenke.} « Les enfants reçoivent beaucoup de cadeaux. »

\all{Die Taufe findet in der Kirche statt.} « Le baptême a lieu à l'église. »
\end{exemplebox}

\begin{attention}[bekommen n'est pas « devenir » !]
Faux ami classique : \all{bekommen} signifie \cle{recevoir}, jamais « devenir ».

\all{Ich habe ein Geschenk bekommen.} « J'ai reçu un cadeau. »

Pour « devenir », l'allemand utilise \all{werden} : \all{Mein Bruder wird Arzt.} « Mon frère devient médecin. »
\end{attention}

\courssub{Champ 4 : sentiments et relations}

\begin{definition}[Sentiments et relations]
\all{sich freuen über (+ acc.)} : se réjouir de (quelque chose de présent ou passé) ; \all{sich freuen auf (+ acc.)} : se réjouir à l'idée de (quelque chose de futur) ; \all{lieben} : aimer ; \all{dankbar sein} : être reconnaissant ; \all{stolz sein (auf + acc.)} : être fier de ; \all{die Freundschaft, -en} : l'amitié ; \all{der Freund, -e / die Freundin, -nen} : l'ami / l'amie ; \all{die Nachbarn} : les voisins.
\end{definition}

La distinction \all{über / auf} après \all{sich freuen} est essentielle :

\begin{center}
\begin{tabularx}{\linewidth}{|X|X|}
\hline
\rowcolor{popgreenL} \textbf{sich freuen über} (présent/passé) & \textbf{sich freuen auf} (futur) \\
\hline
\all{Ich freue mich über das Geschenk.} « Je me réjouis du cadeau. » & \all{Ich freue mich auf das Fest.} « Je me réjouis à l'idée de la fête. » \\
\hline
\all{Wir freuen uns über deinen Besuch.} « Nous sommes heureux de ta visite. » & \all{Wir freuen uns auf die Ferien.} « Nous attendons les vacances avec joie. » \\
\hline
\end{tabularx}
\end{center}

\all{Sie ist stolz auf ihren Sohn.} « Elle est fière de son fils. »

\all{Wir sind unseren Eltern dankbar.} « Nous sommes reconnaissants envers nos parents. » (datif !)

\all{Unsere Nachbarn sind sehr freundlich.} « Nos voisins sont très aimables. »

\courssub{Champ 5 : la société}

\begin{definition}[Vivre en société]
\all{die Gesellschaft, -en} : la société ; \all{die Gemeinde, -n} : la communauté, la commune ; \all{das Dorf, -¨er} : le village ; \all{die Stadt, -¨e} : la ville ; \all{die Tradition, -en} : la tradition ; \all{der Brauch, -¨e} : la coutume ; \all{höflich} : poli ; \all{respektvoll} : respectueux.
\end{definition}

\all{In unserem Dorf gibt es viele alte Traditionen.} « Dans notre village, il y a beaucoup de vieilles traditions. »

\all{Man muss höflich und respektvoll zu den Älteren sein.} « Il faut être poli et respectueux envers les aînés. »

\all{Die Gemeinde organisiert jedes Jahr ein großes Fest.} « La commune organise chaque année une grande fête. »

\courssub{Carte mentale : le champ lexical de la famille}

\begin{popfigure}
\begin{center}
\begin{tikzpicture}[mindmap, grow cyclic, every node/.style={concept, align=center, font=\small}, root concept/.append style={fill=popblue, text=white, font=\bfseries}, level 1/.append style={level distance=3.4cm, sibling angle=72}, concept color=popblue!40]
\node[root concept]{Die Familie}
  child[concept color=poporange!50]{ node[align=center]{Mitglieder\\ \all{der Vater}\\ \all{die Mutter}\\ \all{die Geschwister}} }
  child[concept color=popgreen!50]{ node[align=center]{Alltag\\ \all{zusammen leben}\\ \all{helfen}\\ \all{die Hausarbeit}} }
  child[concept color=popgold!60]{ node[align=center]{Feste\\ \all{der Geburtstag}\\ \all{die Hochzeit}\\ \all{das Geschenk}} }
  child[concept color=poppink!50]{ node[align=center]{Gefühle\\ \all{sich freuen}\\ \all{lieben}\\ \all{stolz sein}} }
  child[concept color=poppurple!40]{ node[align=center]{Gesellschaft\\ \all{das Dorf}\\ \all{die Tradition}\\ \all{höflich}} };
\end{tikzpicture}
\end{center}
\legende{Carte mentale du champ lexical : la famille et la vie sociale.}
\end{popfigure}

\courssub{Comparaison culturelle : Cameroun et Allemagne}

\begin{savaistu}[Großfamilie ou Kleinfamilie ?]
Au Cameroun, la \all{Großfamilie} (la famille élargie) joue un rôle central : oncles, tantes, cousins et grands-parents participent à la vie quotidienne, aux deuils et aux mariages, souvent célébrés avec tout le village. En Allemagne, on vit plus souvent en \all{Kleinfamilie} (famille nucléaire : parents et enfants), mais les grandes fêtes réunissent aussi toute la famille. Dans certaines régions catholiques d'Allemagne (Bavière, Rhénanie), on fête même le \all{Namenstag} (la fête du saint dont on porte le prénom). Et le 18e anniversaire, \all{die Volljährigkeit} (la majorité), est une très grande fête — tout comme la majorité est un événement important au Cameroun.
\end{savaistu}

\begin{center}
\begin{tabularx}{\linewidth}{|X|X|X|}
\hline
\rowcolor{popblueL} \textbf{Français} & \textbf{Deutsch} & \textbf{Remarque} \\
\hline
la famille nucléaire & die Kleinfamilie, -n & parents + enfants \\
\hline
la famille élargie & die Großfamilie, -n & avec oncles, tantes, cousins \\
\hline
la majorité (18 ans) & die Volljährigkeit & grande fête en Allemagne \\
\hline
la fête traditionnelle & das traditionelle Fest & souvent au village \\
\hline
\end{tabularx}
\end{center}

\courssub{Exercice résolu et entraînement}

\begin{exoresolu}[Complète avec le mot qui convient]
Complète : \all{a) Meine \dots\ wohnen in Douala.} (les parents) \all{b) Er \dots\ seinem Freund zum Geburtstag.} (félicite) \all{c) Wir haben ein schönes \dots\ bekommen.} (cadeau) \all{d) Sie hilft \dots\ bei der Hausarbeit.} (à sa tante)
\tcblower
\textbf{Solution détaillée :}

a) \all{Meine \textbf{Eltern} wohnen in Douala.} — \all{die Eltern} est toujours pluriel, le verbe \all{wohnen} est donc correctement au pluriel.

b) \all{Er \textbf{gratuliert} seinem Freund zum Geburtstag.} — \all{gratulieren} est un verbe faible régulier : \all{er gratulier-t}. Attention : \all{seinem Freund} est au datif (verbe à datif).

c) \all{Wir haben ein schönes \textbf{Geschenk} bekommen.} — \all{das Geschenk} est neutre : adjectif en \all{-es} à l'accusatif après \all{ein}. Rappel : \all{bekommen} = recevoir.

d) \all{Sie hilft \textbf{ihrer Tante} bei der Hausarbeit.} — \all{helfen} exige le datif : \all{die Tante} devient \all{der Tante}, possessif \all{ihrer}.
\end{exoresolu}

\begin{exercice}[À toi de jouer]
\begin{enumerate}
\item Traduis en allemand : « Ma sœur invite ses amis à son anniversaire. »
\item Traduis : « Nous nous réjouissons à l'idée du mariage de notre cousin. »
\item Complète au cas correct : \all{Ich danke mein\dots\ Großeltern für das Geschenk.}
\item Trouve le mot : le contraire de \all{sich verstehen} est \dots
\item Classe ces mots dans les cinq champs de la carte mentale : \all{die Tante, kochen, die Einladung, die Freundschaft, das Dorf, stolz sein, der Sohn, die Feier, einkaufen, respektvoll}.
\end{enumerate}
\end{exercice}

\begin{corrige}[Exercice « À toi de jouer »]
1. \all{Meine Schwester lädt ihre Freunde zu ihrem Geburtstag ein.} (particule \all{ein-} en fin de phrase !)

2. \all{Wir freuen uns auf die Hochzeit unseres Cousins.} (\all{sich freuen auf} car le mariage est futur.)

3. \all{Ich danke mein\textbf{en} Großeltern für das Geschenk.} (datif pluriel : \all{meinen}.)

4. \all{sich streiten} (se disputer).

5. Membres : \all{die Tante, der Sohn} ; Quotidien : \all{kochen, einkaufen} ; Fêtes : \all{die Einladung, die Feier} ; Sentiments : \all{die Freundschaft, stolz sein} ; Société : \all{das Dorf, respektvoll}.
\end{corrige}

\begin{aretenir}[L'essentiel du lexique]
\begin{itemize}
\item Chaque nom s'apprend avec article et pluriel : \all{der Bruder, die Brüder} ; \all{die Schwester, -n} ; \all{das Geschenk, -e}.
\item \all{helfen, danken, gratulieren, antworten} + \textbf{datif} : \all{Ich helfe meiner Mutter.}
\item \all{sich freuen über} = se réjouir de (présent/passé) ; \all{sich freuen auf} = se réjouir à l'idée de (futur).
\item \all{bekommen} = recevoir ; « devenir » = \all{werden}.
\item \all{die Eltern} et \all{die Geschwister} : toujours au pluriel.
\end{itemize}
\end{aretenir}

Ce lexique est maintenant ton outil de travail : dans les sections suivantes, nous réviserons le présent et le parfait, puis tu réinvestiras tous ces mots dans une lettre, une description et un dialogue sur la vie en famille et en société.

\coursec{Grammaire 1 : Révision du présent}

Après avoir découvert le vocabulaire de la famille et de la vie sociale dans la section précédente, nous allons maintenant consolider les bases grammaticales indispensables pour produire et traduire des textes sur ce thème. Le présent de l'indicatif (\cle{Präsens}) est le temps le plus utilisé en allemand : il exprime non seulement l'action en cours, mais aussi l'habitude, la vérité générale et, très souvent, le futur proche. Une maîtrise parfaite de ses formes (régulières, irrégulières, modaux, verbes à particule) et de l'ordre des mots est donc la condition \emph{sine qua non} pour réussir les exercices de thème, de version et de production écrite qui suivront.

\courssub{Observation à partir du texte support}

Reprenons quelques phrases extraites du texte \emph{Ein Familienfest in München} (section 1) pour observer le présent en action :

\begin{textebox}[Extraits du texte support]
Meine Familie \textbf{feiert} heute ein großes Fest. Mein Onkel \textbf{kommt} aus Berlin. Meine Tante \textbf{bringt} einen Kuchen. Wir \textbf{sprechen} viel und \textbf{lachen} zusammen. Mein Cousin \textbf{kann} gut Gitarre \textbf{spielen}. Ich \textbf{möchte} ihn später \textbf{besuchen}.
\end{textebox}

\textbf{Observation :} On remarque que le verbe conjugué occupe toujours la \cle{deuxième position} dans la phrase principale (ex. \all{Meine Familie feiert ...}, \all{Mein Onkel kommt ...}). L'infinitif des verbes modaux (\all{spielen}, \all{besuchen}) ou la particule séparable (\all{mitbringen} → \all{bringt ... mit}) se placent en \cle{dernière position}.

\courssub{Les verbes réguliers (schwache Verben)}

La grande majorité des verbes allemands sont réguliers. Leur radical ne change pas ; on ajoute les terminaisons standards au radical (ici \all{mak-} pour \all{machen}) :

\begin{center}
\begin{tabular}{|l|l|l|}
\hline
\rowcolor{popblueL}
\textbf{Personne} & \textbf{Terminaison} & \textbf{Exemple : machen (faire)} \\
\hline
ich & -e & ich mache \\
du & -st & du machst \\
er / sie / es & -t & er macht \\
wir & -en & wir machen \\
ihr & -t & ihr macht \\
sie / Sie & -en & sie / Sie machen \\
\hline
\end{tabular}
\end{center}

\begin{exemplebox}[Verbes réguliers du thème « famille »]
\all{wohnen} (habiter) → \all{Ich wohne in Yaoundé.} « J'habite à Yaoundé. » \\
\all{arbeiten} (travailler) → \all{Mein Vater arbeitet viel.} « Mon père travaille beaucoup. » \\
\all{lernen} (apprendre) → \all{Wir lernen Deutsch.} « Nous apprenons l'allemand. » \\
\all{besuchen} (rendre visite) → \all{Besuchst du deine Oma?} « Rends-tu visite à ta grand-mère ? »
\end{exemplebox}

\courssub{Les verbes à changement de voyelle (Vokalwechsel)}

Certains verbes dits « forts » ou « irréguliers » modifient leur voyelle radicale \textbf{uniquement à la 2e et 3e personnes du singulier} (\all{du}, \all{er/sie/es}). Les autres personnes gardent la voyelle de l'infinitif.

\begin{propriete}[Règle du changement de voyelle]
Le changement de voyelle ne concerne que \textbf{du} et \textbf{er/sie/es}. Les formes \textbf{ich}, \textbf{wir}, \textbf{ihr}, \textbf{sie/Sie} restent régulières.
\end{propriete}

Trois séries principales existent :

\begin{center}
\begin{tabularx}{\linewidth}{|X|X|X|X|}
\hline
\rowcolor{popblueL}
\textbf{Type} & \textbf{Infinitif} & \textbf{du / er-sie-es} & \textbf{Autres personnes} \\
\hline
e → i & \all{sprechen} (parler) & du sprichst / er spricht & ich spreche, wir sprechen \\
& \all{helfen} (aider) & du hilfst / er hilft & ich helfe, wir helfen \\
& \all{essen} (manger) & du isst / er isst & ich esse, wir essen \\
\hline
e → ie & \all{lesen} (lire) & du liest / er liest & ich lese, wir lesen \\
& \all{sehen} (voir) & du siehst / er sieht & ich sehe, wir sehen \\
& \all{empfehlen} (recommander) & du empfiehlst / er empfiehlt & ich empfehle, wir empfehlen \\
\hline
a → ä & \all{fahren} (conduire/aller) & du fährst / er fährt & ich fahre, wir fahren \\
& \all{schlafen} (dormir) & du schläfst / er schläft & ich schlafe, wir schlafen \\
& \all{tragen} (porter) & du trägst / er trägt & ich trage, wir tragen \\
\hline
\end{tabularx}
\end{center}

\begin{exemplebox}[Exemples traduits]
\all{Mein Bruder liest ein Buch.} « Mon frère lit un livre. » (e → ie) \\
\all{Wir müssen unseren Eltern helfen.} « Nous devons aider nos parents. » (e → i à la 3e pers. sg. : \all{er hilft}) \\
\all{Lädst du deine Cousins ein?} « Invites-tu tes cousins ? » (a → ä + particule séparable)
\end{exemplebox}

\courssub{Les verbes irréguliers essentiels : sein, haben, werden, wissen}

Ces quatre verbes sont des piliers de la langue ; ils sont totalement irréguliers au présent et doivent être connus par cœur.

\begin{center}
\begin{tabular}{|l|l|l|l|l|}
\hline
\rowcolor{popblueL}
\textbf{Personne} & \textbf{sein (être)} & \textbf{haben (avoir)} & \textbf{werden (devenir)} & \textbf{wissen (savoir)} \\
\hline
ich & bin & habe & werde & weiß \\
du & bist & hast & wirst & weißt \\
er/sie/es & ist & hat & wird & weiß \\
wir & sind & haben & werden & wissen \\
ihr & seid & habt & werdet & wisst \\
sie/Sie & sind & haben & werden & wissen \\
\hline
\end{tabular}
\end{center}

\begin{attention}[Pièges fréquents pour francophones]
\begin{itemize}
\item \textbf{sein vs haben} : \all{Ich bin müde} (Je suis fatigué) — \all{Ich habe Hunger/Durst} (J'ai faim/soif). En français on « a » faim, en allemand on « est » faim (\all{Hunger haben}).
\item \textbf{wissen vs kennen} : \all{wissen} = savoir une information (\all{Ich weiß das.}) ; \all{kennen} = connaître une personne/lieu (\all{Ich kenne ihn.}).
\item \textbf{werden} : auxiliaire du futur (\all{Ich werde kommen.}) et de la voix passive, mais aussi verbe plein « devenir » (\all{Er wird Arzt.} = Il devient médecin).
\end{itemize}
\end{attention}

\courssub{Les verbes modaux au présent (Modalverben)}

Les six verbes modaux expriment la capacité, l'obligation, la permission, la volonté, le conseil ou le désir. Ils se conjuguent de manière irrégulière (souvent changement de voyelle aux singulières) et gouvernent un \textbf{infinitif à la fin de la phrase}.

\begin{center}
\begin{tabular}{|l|l|l|l|}
\hline
\rowcolor{popblueL}
\textbf{Modal} & \textbf{Sens} & \textbf{ich / er/sie/es} & \textbf{wir / sie/Sie} \\
\hline
können & pouvoir, savoir & kann / kann & können / können \\
müssen & devoir (obligation) & muss / muss & müssen / müssen \\
dürfen & avoir la permission & darf / darf & dürfen / dürfen \\
wollen & vouloir & will / will & wollen / wollen \\
sollen & devoir (conseil) & soll / soll & sollen / sollen \\
mögen & aimer & mag / mag & mögen / mögen \\
\hline
\end{tabular}
\end{center}

\begin{methode}[Conjuguer une phrase avec un verbe modal]
\begin{enumerate}
\item Identifier le \textbf{sujet} (position 1).
\item Conjuguer le \textbf{verbe modal} à la personne du sujet → \textbf{position 2}.
\item Placer le \textbf{complément d'objet / circonstanciel} au milieu.
\item Mettre l'\textbf{infinitif du verbe principal} tout à la \textbf{fin} (dernière position).
\end{enumerate}
\textbf{Exemple :} \all{Ich (1) kann (2) gut Deutsch (3) sprechen (4).}
\end{methode}

\begin{exemplebox}[Verbes modaux en contexte familial]
\all{Meine Schwester kann gut kochen.} « Ma sœur sait bien cuisiner. » \\
\all{Wir müssen das Geschirr spülen.} « Nous devons faire la vaisselle. » \\
\all{Darf ich dein Handy benutzen?} « Puis-je utiliser ton portable ? » \\
\all{Willst du mit uns ins Kino gehen?} « Veux-tu aller au cinéma avec nous ? » \\
\all{Man soll seine Eltern respektieren.} « On doit respecter ses parents. » \\
\all{Ich mag keine laute Musik.} « Je n'aime pas la musique forte. »
\end{exemplebox}

\courssub{Le subjonctif II de politesse : \all{möchte}}

\begin{savaistu}[\all{möchte} — « je voudrais »]
\all{möchte} n'est pas un infinitif, c'est la forme du \textbf{Konjunktiv II} (subjonctif II) de \all{mögen} à la 1e et 3e personnes du singulier. Il s'utilise comme un modal : \all{Ich möchte dich einladen.} « Je voudrais t'inviter. » C'est la forme standard pour exprimer un souhait poli, bien plus courante que \all{ich will} (trop direct).
\end{savaistu}

\courssub{Les verbes à particule séparable (trennbare Verben)}

Beaucoup de verbes courants sont composés d'un radical + une particule (préfixe) qui se détache et va en \textbf{fin de phrase} à l'indicatif présent (et au prétérit). La particule porte l'accent tonique.

\begin{propriete}[Position de la particule séparable]
À l'indicatif présent (et prétérit) en phrase principale : \textbf{Sujet + Verbe conjugué (sans particule) + ... + Particule (fin de phrase)}.
\end{propriete}

\begin{center}
\begin{tabular}{|l|l|l|}
\hline
\rowcolor{popblueL}
\textbf{Verbe composé} & \textbf{Radical + Particule} & \textbf{Exemple au présent} \\
\hline
einladen & lad- + ein & \all{Ich lade meine Freunde ein.} \\
anrufen & ruf- + an & \all{Rufst du deine Mutter an?} \\
aufstehen & steh- + auf & \all{Wir stehen um 7 Uhr auf.} \\
mitbringen & bring- + mit & \all{Er bringt einen Kuchen mit.} \\
abholen & hol- + ab & \all{Holt ihr die Großeltern ab?} \\
\hline
\end{tabular}
\end{center}

\begin{attention}[Erreur classique des francophones]
✗ \all{Ich einlade meine Freunde.} (ordre français : je invite mes amis) \\
✓ \all{Ich lade meine Freunde \textbf{ein}.} \\
\textbf{Règle d'or :} La particule \textbf{ne reste jamais collée au verbe} en position 2 ; elle \textbf{doit} aller à la fin.
\end{attention}

\courssub{Emplois du présent en allemand : comparaison avec le français}

L'allemand n'a \textbf{pas de forme continue} (pas de \emph{être en train de} + infinitif). Le présent simple couvre tous ces sens :

\begin{center}
\begin{tabularx}{\linewidth}{|X|X|X|}
\hline
\rowcolor{popblueL}
\textbf{Emploi} & \textbf{Allemand (Präsens)} & \textbf{Français} \\
\hline
Action en cours & \all{Ich lese gerade ein Buch.} & Je suis en train de lire un livre. \\
Habitude / vérité générale & \all{Meine Mutter kocht gut.} & Ma mère cuisine bien. / Ma mère est une bonne cuisinière. \\
Futur proche (avec indicateur de temps) & \all{Morgen besuche ich meine Tante.} & Demain, je rendrai visite à ma tante / je vais rendre visite à ma tante. \\
Action débutée dans le passé et qui se poursuit (avec \all{seit}) & \all{Ich wohne seit 2020 in Douala.} & J'habite à Douala depuis 2020. \\
\hline
\end{tabularx}
\end{center}

\begin{exemplebox}[Le présent pour le futur : marqueurs temporels]
\all{Heute Abend spielen wir Karten.} « Ce soir, nous jouons aux cartes. » \\
\all{Nächste Woche fährt mein Vater nach Deutschland.} « La semaine prochaine, mon père va en Allemagne. » \\
\all{In einer Stunde essen wir zu Mittag.} « Dans une heure, nous déjeunons. »
\end{exemplebox}

\courssub{Tableau récapitulatif TikZ : conjugaison complète de verbes clés}

\begin{popfigure}
\begin{tikzpicture}[
  node distance=0.3cm and 0.5cm,
  verb/.style={rectangle, draw=popblue, thick, fill=popblueL, rounded corners, minimum width=2.2cm, minimum height=0.8cm, align=center, font=\small},
  form/.style={rectangle, draw=popmuted, fill=white, minimum width=2.2cm, minimum height=0.7cm, align=center, font=\small\ttfamily},
  pers/.style={rectangle, draw=popdark, fill=popgoldL, minimum width=1.8cm, minimum height=0.7cm, align=center, font=\small\bfseries},
  title/.style={font=\bfseries\large, text=popdark}
]
% Titre
\node[title] (titre) at (0,7.5) {Conjugaison au présent : verbes clés du chapitre};

% Colonnes : Personnes
\foreach \p/\y in {ich/6.5, du/5.5, er/4.5, wir/3.5, ihr/2.5, sie/1.5}
  \node[pers] at (-4.5,\y) {\p};

% Verbe sprechen
\node[verb, align=center] (sp) at (0,6.5){sprechen\\(e→i)};
\foreach \f/\y in {spreche/6.5, sprichst/5.5, spricht/4.5, sprechen/3.5, sprecht/2.5, sprechen/1.5}
  \node[form] at (0,\y) {\f};

% Verbe fahren
\node[verb, align=center] (fa) at (3,6.5){fahren\\(a→ä)};
\foreach \f/\y in {fahre/6.5, fährst/5.5, fährt/4.5, fahren/3.5, fahrt/2.5, fahren/1.5}
  \node[form] at (3,\y) {\f};

% Verbe sein
\node[verb, align=center] (se) at (6,6.5){sein\\(irrégulier)};
\foreach \f/\y in {bin/6.5, bist/5.5, ist/4.5, sind/3.5, seid/2.5, sind/1.5}
  \node[form] at (6,\y) {\f};

% Verbe haben
\node[verb, align=center] (ha) at (9,6.5){haben\\(irrégulier)};
\foreach \f/\y in {habe/6.5, hast/5.5, hat/4.5, haben/3.5, habt/2.5, haben/1.5}
  \node[form] at (9,\y) {\f};

% Verbe können
\node[verb, align=center] (ko) at (12,6.5){können\\(modal)};
\foreach \f/\y in {kann/6.5, kannst/5.5, kann/4.5, können/3.5, könnt/2.5, können/1.5}
  \node[form] at (12,\y) {\f};

% Flèches de liaison
\draw[poporange, thick, -{Stealth[length=3mm]}] (sp.east) -- (fa.west) node[midway, above, font=\tiny, text=poporange] {chgt voyelle};
\draw[poporange, thick, -{Stealth[length=3mm]}] (fa.east) -- (se.west) node[midway, above, font=\tiny, text=poporange] {irrégulier};
\draw[poporange, thick, -{Stealth[length=3mm]}] (se.east) -- (ha.west) node[midway, above, font=\tiny, text=poporange] {irrégulier};
\draw[poporange, thick, -{Stealth[length=3mm]}] (ha.east) -- (ko.west) node[midway, above, font=\tiny, text=poporange] {modal};
\end{tikzpicture}
\legende{Tableau de conjugaison au présent : sprechen, fahren, sein, haben, können. Les cases colorées indiquent les formes à changement de voyelle ou irrégulières.}
\end{popfigure}

\courssub{Exercice résolu : conjugaison et ordre des mots}

\begin{exoresolu}[Mettre les verbes entre parenthèses au présent et placer la particule le cas échéant]
\textbf{Énoncé :}
\begin{enumerate}
\item Meine Eltern \_\_\_\_\_\_ (kommen) heute aus Yaoundé.
\item Mein Bruder \_\_\_\_\_\_ (lesen) gerade eine Zeitung.
\item Wir \_\_\_\_\_\_ (müssen) unserer Oma \_\_\_\_\_\_ (helfen).
\item \_\_\_\_\_\_ du deine Tante \_\_\_\_\_\_ (anrufen)?
\item Ich \_\_\_\_\_\_ (möchte) dich zu meinem Geburtstag \_\_\_\_\_\_ (einladen).
\end{enumerate}

\textbf{Solution détaillée :}
\begin{enumerate}
\item \textbf{kommen} → 3e pers. pl. \textbf{kommen} (verbe régulier). \\
      \all{Meine Eltern \textbf{kommen} heute aus Yaoundé.}
\item \textbf{lesen} → 3e pers. sg. changement e→ie : \textbf{liest}. \all{gerade} marque l'action en cours. \\
      \all{Mein Bruder \textbf{liest} gerade eine Zeitung.}
\item \textbf{müssen} → 1e pers. pl. \textbf{müssen} (modal) + infinitif \textbf{helfen} à la fin. \\
      \all{Wir \textbf{müssen} unserer Oma \textbf{helfen}.} (Datif \all{unserer Oma} après \all{helfen}).
\item \textbf{anrufen} → particule \textbf{an} séparable. Question : verbe en position 1 → \textbf{rufst}, particule \textbf{an} à la fin. \\
      \all{\textbf{Rufst} du deine Tante \textbf{an}?}
\item \textbf{möchte} → 1e pers. sg. de \all{mögen} (Konjunktiv II) + \textbf{einladen} (particule \textbf{ein} à la fin). \\
      \all{Ich \textbf{möchte} dich zu meinem Geburtstag \textbf{einladen}.}
\end{enumerate}
\end{exoresolu}

\courssub{Texte d'entraînement : « Ein Sonntag bei der Familie Ngo »}

Le texte suivant réinvestit le lexique familial (section 2) et les formes du présent étudiées ici. Il servira de base aux exercices de traduction (section 7) et de production (section 6).

\begin{textebox}[Ein Sonntag bei der Familie Ngo]
Am Morgen \textbf{stehen} alle früh \textbf{auf}. Der Vater \textbf{ruft} die Großeltern \textbf{an}. An einem sonnigen Sonntag \textbf{besucht} die Familie Ngo ihre Großeltern in einem Vorort von Yaoundé. Der Großvater, Herr Ngo, \textbf{sitzt} auf der Veranda und \textbf{liest} die Zeitung. Die Großmutter \textbf{kocht} in der Küche ein traditionelles Gericht: Ndolé mit Reis. Der Duft \textbf{verbreitet} sich im ganzen Haus.

Die Kinder, Paul und Marie, \textbf{spielen} im Garten mit ihren Cousins. Sie \textbf{können} gut Fußball \textbf{spielen}. Paul \textbf{will} später Profi-Fußballer \textbf{werden}. Marie \textbf{mag} lieber tanzen; sie \textbf{tanzt} oft zu Makossa-Musik.

Zur Mittagszeit \textbf{ruft} die Großmutter: „Das Essen ist fertig!“ Alle \textbf{kommen} zusammen, \textbf{waschen} sich die Hände und \textbf{setzen} sich an den großen Tisch. Herr Ngo \textbf{bedankt} sich für das Essen und \textbf{erzählt} Geschichten aus seiner Jugend. Die Enkel \textbf{hören} aufmerksam \textbf{zu}.

Am Nachmittag \textbf{müssen} die Eltern wieder \textbf{fahren}. Die Großeltern \textbf{winken} zum Abschied. „Wir \textbf{kommen} bald wieder“, \textbf{verspricht} der Vater. Die Kinder \textbf{winken} zurück und \textbf{rufen}: „Tschüss, Oma! Tschüss, Opa!“
\end{textebox}

\begin{center}
\begin{tabular}{|l|l|}
\hline
\rowcolor{popblueL}
\textbf{Allemand} & \textbf{Français} \\
\hline
der Vorort, die Vororte & la banlieue, la périphérie \\
die Veranda, die Verandas & la véranda \\
das Gericht, die Gerichte & le plat \\
der Duft, die Düfte & le parfum, l'odeur \\
sich verbreiten (verbreitet sich) & se répandre \\
der Profi-Fußballer & le footballeur professionnel \\
die Makossa-Musik & la musique makossa \\
zur Mittagszeit & à l'heure du déjeuner \\
sich bedanken (bedankt sich) & remercier \\
erzählen & raconter, narrer \\
zu-hören (hört ... zu) & écouter attentivement \\
winken & faire un signe de la main \\
versprechen & promettre \\
zurück-winken (winkt ... zurück) & répondre au signe \\
an-rufen (ruft ... an) & appeler (par téléphone) \\
auf-stehen (stehen ... auf) & se lever \\
\hline
\end{tabular}
\end{center}

\textbf{Questions de compréhension (oral ou écrit) :}
\begin{enumerate}
\item Wo besucht die Familie Ngo ihre Großeltern? (Dans un quartier périphérique de Yaoundé.)
\item Was kocht die Großmutter? (Du Ndolé avec du riz.)
\item Was wollen Paul und Marie später werden? (Paul veut devenir footballeur pro ; Marie préfère danser.)
\item Welches Modalverb drückt einen Wunsch aus? (\all{möchte} / \all{will} — ici \all{will} pour Paul.)
\item Finde drei trennbare Verben im Text und nenne ihre Partikeln. (\textbf{Réponse :} \all{auf-stehen} → \all{auf}, \all{an-rufen} → \all{an}, \all{zu-hören} → \all{zu}, \all{zurück-winken} → \all{zurück}.)
\end{enumerate}

\courssub{Mini-production guidée : décrire sa famille au présent}

\begin{exercice}[Ma famille au présent]
Rédigez 6 à 8 phrases en allemand pour décrire votre famille un dimanche typique. Utilisez :
\begin{itemize}
\item au moins 2 verbes réguliers (\all{wohnen}, \all{arbeiten}, \all{spielen}…)
\item au moins 1 verbe à changement de voyelle (\all{lesen}, \all{sehen}, \all{fahren}…)
\item au moins 1 verbe modal (\all{können}, \all{müssen}, \all{wollen}, \all{möchte}…)
\item au moins 1 verbe à particule séparable (\all{anrufen}, \all{aufstehen}, \all{mitbringen}…)
\item les verbes \all{sein} et \all{haben}
\end{itemize}
Respectez l'ordre des mots (verbe position 2, infinitif/particule à la fin). Relisez-vous : majuscules aux noms, Umlaute, ß.
\end{exercice}

\begin{corrige}[Exercice « Ma famille au présent » — proposition de correction]
\textbf{Modèle :}
\all{Meine Familie wohnt in Douala. Mein Vater arbeitet in einer Bank. Meine Mutter kocht gerne. Am Sonntag stehen wir spät auf. Mein Bruder liest gerne Comics. Ich kann gut Gitarre spielen. Wir müssen oft unserer Oma helfen. Ich möchte sie bald besuchen.}

\textbf{Points de vigilance :}
\begin{itemize}
\item \all{stehen ... auf} (particule \all{auf} à la fin), \all{kann ... spielen} (infinitif à la fin), \all{helfen} (changement e→i : \all{er hilft} mais \all{wir helfen}), \all{möchte ... besuchen}.
\item Vérifiez les majuscules : \all{Familie, Vater, Mutter, Sonntag, Bruder, Comics, Gitarre, Oma}.
\end{itemize}
\end{corrige}

\courssub{Synthèse : checklist du présent pour l'auto-évaluation}

\begin{aretenir}[Checklist : suis-je prêt(e) pour la suite ?]
\begin{itemize}
\item[☐] Je conjugue correctement un verbe régulier aux 6 personnes.
\item[☐] J'applique le changement de voyelle \textbf{uniquement} à \all{du} et \all{er/sie/es} (e→i, e→ie, a→ä).
\item[☐] Je connais par cœur \all{sein, haben, werden, wissen} au présent.
\item[☐] Je construis une phrase avec un modal : \textbf{Sujet + Modal (pos. 2) + ... + Infinitif (fin)}.
\item[☐] Je place la particule séparable \textbf{à la fin} de la phrase principale.
\item[☐] J'utilise le présent pour exprimer le futur proche (\all{Morgen fahre ich ...}).
\item[☐] Je ne traduis pas « je suis en train de » par une forme continue, mais par \all{gerade} + présent.
\item[☐] Je mets une majuscule à \textbf{tous} les noms (\all{die Familie, der Sonntag, das Essen}).
\end{itemize}
\end{aretenir}

Cette révision du présent constitue la charpente grammaticale sur laquelle nous appuierons la leçon suivante (Perfekt) et les activités de production/traduction. Prenez le temps de mémoriser les tableaux et de refaire l'exercice résolu à l'oral avant de passer à la section 4.

\coursec{Grammaire 2 : le parfait (Perfekt)}

Dans la section précédente, nous avons révisé le présent de l'indicatif, le temps de la description et des habitudes : \all{Meine Familie wohnt in Yaoundé.} Mais dans notre texte support \all{Ein Familienfest in München}, la narratrice \emph{racontait} un événement passé : la fête de famille du week-end dernier. Pour raconter le passé en allemand parlé et dans les lettres personnelles, on utilise un temps composé : le \cle{Perfekt}. C'est le temps que tu utiliseras dans toutes tes productions de ce chapitre (lettre, description, dialogue).

\courssub{Observation : repérer le Perfekt dans un texte}

\begin{textebox}[Am Sonntag bei Oma — un dimanche chez grand-mère]
Am Samstag haben wir unsere Großeltern in Douala besucht. Wir sind um 8 Uhr mit dem Bus gefahren. Meine Mutter hat einen Kuchen gebacken, und mein Vater hat Getränke gekauft. Um 12 Uhr sind wir angekommen. Oma hat uns an der Tür empfangen und alle haben sich gefreut. Zuerst haben wir zusammen gegessen, dann haben die Kinder im Garten gespielt. Am Abend sind wir müde nach Hause gefahren. Es war ein schöner Tag!
\end{textebox}

\textbf{Glossaire :} \all{der Bus, -se} (le bus) ; \all{der Kuchen, -} (le gâteau) ; \all{das Getränk, -e} (la boisson) ; \all{empfangen} (accueillir) ; \all{sich freuen} (se réjouir) ; \all{der Garten, die Gärten} (le jardin) ; \all{müde} (fatigué).

\textbf{Observe :} souligne dans le texte toutes les formes verbales. Que remarques-tu ? Chaque verbe conjugué est en \emph{deux morceaux} : un petit verbe (\all{haben} ou \all{sein}) en deuxième position, et une forme en \all{ge-...-t} ou \all{ge-...-en} tout à la fin de la phrase. C'est exactement la structure du Perfekt.

\courssub{Formation générale du Perfekt}

\begin{propriete}[Formation du Perfekt]
Le Perfekt se forme avec deux éléments :
\begin{enumerate}
\item l'auxiliaire \cle{haben} ou \cle{sein}, conjugué au \textbf{présent}, en \textbf{position 2} de la phrase ;
\item le \cle{participe II} (Partizip II) du verbe principal, placé \textbf{à la fin} de la phrase.
\end{enumerate}
\all{Wir \textbf{haben} unsere Großeltern \textbf{besucht}.} « Nous avons rendu visite à nos grands-parents. »
\end{propriete}

\begin{center}
\begin{tabular}{|l|l|l|}
\hline
\rowcolor{popblueL} \textbf{Personne} & \textbf{haben} & \textbf{sein} \\
\hline
ich & habe & bin \\
\hline
du & hast & bist \\
\hline
er/sie/es & hat & ist \\
\hline
wir & haben & sind \\
\hline
ihr & habt & seid \\
\hline
sie/Sie & haben & sind \\
\hline
\end{tabular}
\end{center}

\begin{attention}[Le participe II toujours en fin de phrase]
Contrairement au français (« j'ai \emph{fait} les devoirs »), le participe allemand se place \textbf{tout à la fin}, après tous les compléments :
\begin{itemize}
\item \texttimes{} \all{Ich habe gemacht die Hausarbeit.}
\item \checkmark{} \all{Ich habe die Hausarbeit gemacht.} « J'ai fait le ménage. »
\end{itemize}
Entre l'auxiliaire (position 2) et le participe (fin), on peut placer autant de compléments qu'on veut : \all{Meine Mutter hat gestern Abend für die ganze Familie einen großen Kuchen gebacken.} « Hier soir, ma mère a fait un grand gâteau pour toute la famille. »
\end{attention}

\courssub{Le participe II des verbes réguliers (faibles)}

\begin{propriete}[Participe des verbes faibles]
Pour les verbes réguliers : \cle{ge- + radical + -t}.
\begin{itemize}
\item \all{machen} $\rightarrow$ \all{ge\textbf{mach}t} (fait)
\item \all{feiern} $\rightarrow$ \all{ge\textbf{feier}t} (fêté)
\item \all{kaufen} $\rightarrow$ \all{ge\textbf{kauf}t} (acheté)
\item \all{spielen} $\rightarrow$ \all{ge\textbf{spiel}t} (joué)
\end{itemize}
Exception phonétique : les verbes dont le radical se termine en \all{-t}, \all{-d} ou en groupe de consonnes difficile prennent \all{-et} : \all{arbeiten} $\rightarrow$ \all{gearbeitet} (travaillé), \all{öffnen} $\rightarrow$ \all{geöffnet} (ouvert).
\end{propriete}

\courssub{Le participe II des verbes forts}

\begin{propriete}[Participe des verbes forts]
Pour les verbes forts (irréguliers) : \cle{ge- + radical souvent modifié + -en}. Il faut les apprendre par cœur, comme les verbes irréguliers anglais :
\begin{itemize}
\item \all{gehen} $\rightarrow$ \all{gegangen} (allé)
\item \all{essen} $\rightarrow$ \all{gegessen} (mangé)
\item \all{trinken} $\rightarrow$ \all{getrunken} (bu)
\item \all{kommen} $\rightarrow$ \all{gekommen} (venu)
\item \all{sehen} $\rightarrow$ \all{gesehen} (vu)
\item \all{geben} $\rightarrow$ \all{gegeben} (donné)
\item \all{schreiben} $\rightarrow$ \all{geschrieben} (écrit)
\item \all{backen} $\rightarrow$ \all{gebacken} (cuit au four)
\end{itemize}
\end{propriete}

\courssub{Cas particuliers : pas de « ge- » ou « ge » à l'intérieur}

\begin{propriete}[Trois cas particuliers du participe II]
\begin{enumerate}
\item \textbf{Verbes à préfixe inséparable} (\all{be-}, \all{ver-}, \all{er-}, \all{ent-}, \all{emp-}, \all{ge-}, \all{zer-}) : \textbf{pas de « ge- »} : \all{besuchen} $\rightarrow$ \all{besucht} ; \all{verkaufen} $\rightarrow$ \all{verkauft} ; \all{erzählen} $\rightarrow$ \all{erzählt}.
\item \textbf{Verbes en \all{-ieren}} (souvent d'origine française) : \textbf{pas de « ge- »} : \all{studieren} $\rightarrow$ \all{studiert} ; \all{telefonieren} $\rightarrow$ \all{telefoniert} ; \all{reparieren} $\rightarrow$ \all{repariert}.
\item \textbf{Verbes à particule séparable} : le « ge » se place \textbf{entre la particule et le radical} : \all{einladen} $\rightarrow$ \all{ein\textbf{ge}laden} (invité) ; \all{vorbereiten} $\rightarrow$ \all{vor\textbf{ge}bereitet} (préparé) ; \all{ankommen} $\rightarrow$ \all{an\textbf{ge}kommen} (arrivé) ; \all{aufstehen} $\rightarrow$ \all{auf\textbf{ge}standen} (levé).
\end{enumerate}
\end{propriete}

\begin{exemplebox}[Exemples traduits]
\begin{itemize}
\item \all{Meine Mutter hat einen Kuchen gebacken.} « Ma mère a fait un gâteau. »
\item \all{Die Gäste sind um 18 Uhr gekommen.} « Les invités sont arrivés à 18 heures. »
\item \all{Ich habe meine Tante besucht.} « J'ai rendu visite à ma tante. » (pas de « ge- » avec \all{be-} !)
\item \all{Wir haben die Verwandten eingeladen.} « Nous avons invité les membres de la famille. »
\item \all{Mein Bruder hat in Bafoussam studiert.} « Mon frère a étudié à Bafoussam. »
\end{itemize}
\end{exemplebox}

\courssub{Le choix de l'auxiliaire : haben ou sein ?}

C'est LA grande question du Perfekt. En français, on hésite aussi (« j'ai monté / je suis monté »), mais en allemand la règle est plus simple et plus régulière.

\begin{propriete}[Règle de l'auxiliaire]
\begin{itemize}
\item \cle{sein} s'emploie avec les verbes exprimant un \textbf{déplacement} (d'un point A vers un point B) ou un \textbf{changement d'état}, ainsi qu'avec \all{sein} et \all{bleiben} : \all{gehen, fahren, kommen, fliegen, laufen, reisen, aufstehen, sterben, einschlafen, werden, bleiben, sein}.
\all{Wir sind nach Hause gefahren.} « Nous sommes rentrés à la maison. »
\item \cle{haben} s'emploie \textbf{partout ailleurs} : verbes transitifs, verbes d'activité, verbes réfléchis, verbes modaux.
\all{Wir haben ein Fest gefeiert.} « Nous avons fêté une fête. »
\end{itemize}
\end{propriete}

\begin{popfigure}
\begin{center}
\begin{tikzpicture}[
  node distance=0.9cm and 1.6cm,
  question/.style={diamond, draw=popblue, fill=popblueL, text width=4.6cm, align=center, inner sep=2pt, aspect=2.2},
  result/.style={rectangle, rounded corners, draw=popdark, text width=4.2cm, align=center, inner sep=5pt},
  arrow/.style={-{Stealth[length=3mm]}, thick, popdark}
]
\node[question] (q) {Le verbe exprime-t-il un \textbf{déplacement} ou un \textbf{changement d'état} ? (ou \all{sein/bleiben})};
\node[result, fill=popgreenL, draw=popgreen, below left=1.1cm and -0.4cm of q, align=center] (sein){\textbf{Auxiliaire SEIN}\\ \all{Wir \textbf{sind} gekommen.}\\ \all{Sie \textbf{ist} gestorben.}};
\node[result, fill=poporangeL, draw=poporange, below right=1.1cm and -0.4cm of q, align=center] (haben){\textbf{Auxiliaire HABEN}\\ \all{Wir \textbf{haben} gefeiert.}\\ \all{Ich \textbf{habe} besucht.}};
\draw[arrow] (q) -- node[left, align=center]{OUI} (sein);
\draw[arrow] (q) -- node[right, align=center]{NON} (haben);
\end{tikzpicture}
\end{center}
\legende{Arbre de décision : \all{haben} ou \all{sein} au Perfekt ?}
\end{popfigure}

\begin{center}
\begin{tabular}{|l|l|l|}
\hline
\rowcolor{popblueL} \textbf{Verbe} & \textbf{Auxiliaire} & \textbf{Exemple au Perfekt} \\
\hline
\all{gehen} (aller) & sein & \all{Ich bin in die Stadt gegangen.} \\
\hline
\all{fahren} (aller en véhicule) & sein & \all{Wir sind nach Douala gefahren.} \\
\hline
\all{kommen} (venir) & sein & \all{Die Gäste sind gekommen.} \\
\hline
\all{fliegen} (voler) & sein & \all{Er ist nach Deutschland geflogen.} \\
\hline
\all{aufstehen} (se lever) & sein & \all{Ich bin um 6 Uhr aufgestanden.} \\
\hline
\all{sterben} (mourir) & sein & \all{Der Großvater ist gestorben.} \\
\hline
\all{bleiben} (rester) & sein & \all{Wir sind zu Hause geblieben.} \\
\hline
\all{feiern} (fêter) & haben & \all{Wir haben gefeiert.} \\
\hline
\all{besuchen} (rendre visite) & haben & \all{Ich habe meine Tante besucht.} \\
\hline
\all{essen} (manger) & haben & \all{Wir haben zusammen gegessen.} \\
\hline
\end{tabular}
\end{center}

\begin{attention}[Jamais d'accord du participe II en allemand]
En français, on accorde : « elle est allé\emph{e} », « ils sont arrivé\emph{s} ». En allemand, le participe II est \textbf{invariable}, même avec \all{sein} :
\begin{itemize}
\item \texttimes{} \all{sie ist gegangene}
\item \checkmark{} \all{sie ist gegangen} « elle est allée »
\item \checkmark{} \all{die Gäste sind gekommen} « les invités sont arrivés »
\end{itemize}
\end{attention}

\courssub{Emploi du Perfekt : le passé de la communication}

\begin{aretenir}[Quand utilise-t-on le Perfekt ?]
Le Perfekt est le temps du passé dans la \textbf{langue parlée}, les \textbf{dialogues} et les \textbf{lettres personnelles} — exactement les types de textes de cette leçon. Le \all{Präteritum} (passé simple) est réservé au récit écrit littéraire, aux journaux et aux verbes \all{sein} (\all{war}), \all{haben} (\all{hatte}) et les modaux. C'est pourquoi notre texte support dit \all{Es war ein schöner Tag!} et non \all{Es ist ein schöner Tag gewesen} : avec \all{sein}, le Präteritum est toujours préféré, même à l'oral.
\end{aretenir}

\begin{center}
\begin{tabularx}{\linewidth}{|X|X|X|}
\hline
\rowcolor{popblueL} \textbf{Français} & \textbf{Deutsch} & \textbf{Remarque} \\
\hline
Nous avons fêté l'anniversaire de papa. & \all{Wir haben den Geburtstag von Papa gefeiert.} & passé composé = Perfekt \\
\hline
Elle est arrivée à 18 heures. & \all{Sie ist um 18 Uhr gekommen.} & auxiliaire \all{sein}, pas d'accord \\
\hline
J'ai téléphoné à ma sœur. & \all{Ich habe meine Schwester angerufen.} & particule séparable : \all{an\textbf{ge}rufen} \\
\hline
Ils sont restés trois jours chez nous. & \all{Sie sind drei Tage bei uns geblieben.} & \all{bleiben} prend \all{sein} \\
\hline
\end{tabularx}
\end{center}

\begin{popfigure}
\begin{center}
\begin{tikzpicture}[scale=1, every node/.style={align=center}]
\draw[-{Stealth[length=3mm]}, thick, popdark] (0,0) -- (11,0);
\foreach \x/\label in {1/{Vergangenheit\\ (passé)}, 5.5/{JETZT\\ (présent)}, 10/{Zukunft\\ (futur)}} {
  \draw[thick, popdark] (\x,-0.15) -- (\x,0.15);
  \node[below] at (\x,-0.2) {\label};
}
\node[rectangle, rounded corners, draw=poporange, fill=poporangeL, text width=3.4cm, inner sep=4pt] at (2.6,1.3) {\textbf{Perfekt}\\ \all{Wir haben gefeiert.}\\ \all{Wir sind gefahren.}};
\draw[-{Stealth[length=2.5mm]}, poporange, thick] (2.6,0.9) -- (1.6,0.15);
\node[rectangle, rounded corners, draw=popblue, fill=popblueL, text width=3.2cm, inner sep=4pt] at (7.6,1.3) {\textbf{Präsens}\\ \all{Wir feiern.}\\ \all{Wir fahren.}};
\draw[-{Stealth[length=2.5mm]}, popblue, thick] (7.6,0.9) -- (5.9,0.15);
\end{tikzpicture}
\end{center}
\legende{Frise chronologique : le Perfekt situe l'action avant le moment présent.}
\end{popfigure}

\begin{methode}[Raconter un événement familial au Perfekt]
Pour raconter un week-end ou une fête en famille, construis chaque phrase selon la chaîne :
\begin{enumerate}
\item \textbf{Connecteur temporel} en position 1 : \all{zuerst} (d'abord), \all{dann} (ensuite), \all{danach} (après), \all{am Abend} (le soir), \all{schließlich} (finalement) ;
\item \textbf{auxiliaire} \all{haben/sein} conjugué en position 2 ;
\item \textbf{compléments} (sujet si besoin, temps, lieu, objets) ;
\item \textbf{participe II} en toute fin.
\end{enumerate}
Modèle : \all{Dann haben wir zusammen gegessen.} « Ensuite nous avons mangé ensemble. »
Attention : si le connecteur occupe la position 1, le sujet passe après l'auxiliaire (inversion) : \all{Am Abend \textbf{sind wir} nach Hause gefahren.}
\end{methode}

\begin{exoresolu}[Mets les verbes au Perfekt]
Énoncé — Conjugue au Perfekt : 1. \all{wir / besuchen / unsere Großeltern} ; 2. \all{die Gäste / kommen / um 18 Uhr} ; 3. \all{ich / einladen / meine Freunde} ; 4. \all{meine Mutter / backen / einen Kuchen} ; 5. \all{wir / fahren / nach Douala} ; 6. \all{er / studieren / in Yaoundé}.
\tcblower
Solution détaillée :
\begin{enumerate}
\item \all{Wir haben unsere Großeltern besucht.} — verbe régulier avec préfixe inséparable \all{be-} : pas de « ge- » ; \all{besuchen} n'est pas un déplacement, donc \all{haben}.
\item \all{Die Gäste sind um 18 Uhr gekommen.} — \all{kommen} = déplacement $\rightarrow$ \all{sein} ; verbe fort : \all{gekommen}.
\item \all{Ich habe meine Freunde eingeladen.} — particule séparable \all{ein-} : « ge » à l'intérieur $\rightarrow$ \all{eingeladen} ; auxiliaire \all{haben}.
\item \all{Meine Mutter hat einen Kuchen gebacken.} — verbe fort, auxiliaire \all{haben} (activité, pas déplacement).
\item \all{Wir sind nach Douala gefahren.} — déplacement $\rightarrow$ \all{sein} ; participe fort \all{gefahren}.
\item \all{Er hat in Yaoundé studiert.} — verbe en \all{-ieren} : pas de « ge- » $\rightarrow$ \all{studiert} ; auxiliaire \all{haben}.
\end{enumerate}
\end{exoresolu}

\begin{exercice}[À ton tour : un samedi en famille]
Mets le texte suivant au Perfekt. Verbes donnés à l'infinitif entre parenthèses :

\all{Am Samstag (aufstehen) wir um 7 Uhr. Meine Mutter (kochen) Ndolé. Wir (essen) zusammen. Dann (fahren) wir zum Markt in Mokolo. Mein Vater (kaufen) Gemüse und Fisch. Am Nachmittag (besuchen) wir unsere Tante. Am Abend (kommen) wir müde nach Hause.}
\end{exercice}

\begin{corrige}[Exercice « Un samedi en famille »]
\all{Am Samstag \textbf{sind} wir um 7 Uhr \textbf{aufgestanden}. Meine Mutter \textbf{hat} Ndolé \textbf{gekocht}. Wir \textbf{haben} zusammen \textbf{gegessen}. Dann \textbf{sind} wir zum Markt in Mokolo \textbf{gefahren}. Mein Vater \textbf{hat} Gemüse und Fisch \textbf{gekauft}. Am Nachmittag \textbf{haben} wir unsere Tante \textbf{besucht}. Am Abend \textbf{sind} wir müde nach Hause \textbf{gekommen}.}

Traduction : « Samedi, nous nous sommes levés à 7 heures. Ma mère a cuisiné du ndolé. Nous avons mangé ensemble. Ensuite nous sommes allés au marché de Mokolo. Mon père a acheté des légumes et du poisson. L'après-midi, nous avons rendu visite à notre tante. Le soir, nous sommes rentrés fatigués à la maison. »

Points de vigilance : \all{aufstehen, fahren, kommen} prennent \all{sein} (déplacement/changement d'état) ; \all{besucht} sans « ge- » ; après un connecteur en position 1, inversion du sujet (\all{Dann sind wir...}).
\end{corrige}

\begin{attention}[Les trois erreurs préférées des francophones]
\begin{enumerate}
\item Oublier la fin de phrase : \texttimes{} \all{Ich habe gegessen Reis} $\rightarrow$ \checkmark{} \all{Ich habe Reis gegessen.}
\item Mettre « ge- » partout : \texttimes{} \all{gebesucht}, \texttimes{} \all{gestudiert} $\rightarrow$ \checkmark{} \all{besucht}, \all{studiert}.
\item Copier l'auxiliaire français : « j'ai allé » n'existe pas, mais attention au réflexe inverse : \all{ich habe gegangen} est faux $\rightarrow$ \checkmark{} \all{ich bin gegangen} (déplacement = \all{sein}).
\end{enumerate}
\end{attention}

\begin{savaistu}[Pourquoi « Perfekt » ?]
Le nom vient du latin \emph{perfectum}, « achevé » : le Perfekt présente l'action comme terminée, avec un résultat visible au moment où l'on parle. En allemand du Sud (Bavière, Autriche, Suisse), on l'utilise encore plus qu'au Nord : un Bavarois dira presque toujours \all{ich habe gesagt} là où un Nord-Allemand écrira \all{ich sagte}. Bonne nouvelle pour toi : dans tes lettres et dialogues, le Perfekt est toujours correct !
\end{savaistu}

\coursec{Grammaire 3 : l'ordre des mots}

Dans la section précédente, nous avons étudié le \cle{Perfekt} et nous avons déjà remarqué un phénomène étrange : le verbe conjugué (\all{habe}, \all{bin}) reste en deuxième position pendant que le participe passé (\all{gemacht}, \all{geblieben}) s'en va tout à la fin de la phrase. Ce n'est pas un hasard : c'est l'une des lois fondamentales de la phrase allemande. Dans cette section, nous allons observer, comprendre et appliquer la règle la plus importante de la syntaxe allemande : la place du verbe.

\courssub{Observation : où est le verbe ?}

Lisons d'abord ce petit texte. Soulignons mentalement le verbe conjugué de chaque phrase et comptons sa position.

\begin{textebox}[Ein Besuch bei den Großeltern]
\all{Meine Familie wohnt in Yaoundé. Am Wochenende besuchen wir oft unsere Großeltern. Mein Großvater erzählt uns gern Geschichten aus seiner Jugend. Gestern sind wir mit dem Bus zu ihnen gefahren. Um acht Uhr haben wir das Haus verlassen. Meine Schwester hat einen Kuchen gebacken, und ich habe Früchte gekauft. Bei meinen Großeltern gibt es immer viel zu essen. Nach dem Essen spielen wir Karten oder wir hören Musik. Am Abend fahren wir müde, aber glücklich nach Hause. Nächste Woche feiert meine Großmutter ihren Geburtstag. Dann kommen auch meine Onkel und Tanten aus Douala. Wir werden sicher viel lachen und tanzen.}
\par\medskip
\emph{Glossaire :} der Großvater, die Großväter (le grand-père) ; die Großmutter, die Großmütter (la grand-mère) ; die Jugend (la jeunesse) ; der Kuchen, die Kuchen (le gâteau) ; die Frucht, die Früchte (le fruit) ; die Tante, die Tanten (la tante) ; der Geburtstag, die Geburtstage (l'anniversaire) ; lachen (rire) ; tanzen (danser).
\end{textebox}

Comptons ensemble les positions dans quelques phrases :

\begin{center}
\begin{tabularx}{\linewidth}{|X|l|l|}\hline
\rowcolor{popblueL} Phrase & Élément en 1\textsuperscript{re} position & Verbe conjugué \\ \hline
\all{Meine Familie wohnt in Yaoundé.} & \all{Meine Familie} (sujet) & \all{wohnt} : 2\textsuperscript{e} \\ \hline
\all{Am Wochenende besuchen wir oft unsere Großeltern.} & \all{Am Wochenende} (complément) & \all{besuchen} : 2\textsuperscript{e} \\ \hline
\all{Gestern sind wir mit dem Bus gefahren.} & \all{Gestern} (complément) & \all{sind} : 2\textsuperscript{e} \\ \hline
\all{Nächste Woche feiert meine Großmutter ihren Geburtstag.} & \all{Nächste Woche} (complément) & \all{feiert} : 2\textsuperscript{e} \\ \hline
\end{tabularx}
\end{center}

Observation capitale : que la phrase commence par le sujet ou par un complément de temps, le verbe conjugué reste toujours à la \textbf{deuxième place}. C'est le sujet qui se déplace, jamais le verbe !

\courssub{La règle d'or : le verbe en deuxième position (V2)}

\begin{propriete}[La règle V2]
Dans une phrase déclarative allemande, le \cle{verbe conjugué} occupe \textbf{toujours la deuxième position}, quelle que soit la nature de l'élément placé en première position (sujet, complément de temps, complément de lieu, etc.).
\par\medskip
\textbf{Attention :} la « première position » peut contenir \emph{plusieurs mots} si ces mots forment un seul bloc (un seul élément de sens) : \all{Am Wochenende} = un seul bloc temporel ; \all{Meine Familie} = un seul bloc sujet.
\end{propriete}

\textbf{Cas 1 : la phrase standard (sujet d'abord).} C'est l'ordre qui ressemble au français : sujet + verbe + compléments.

\all{Ich besuche morgen meine Großeltern.} « Je rends visite à mes grands-parents demain. »

\begin{center}
\begin{tabular}{|l|l|l|l|}\hline
\rowcolor{popgreenL} Position 1 & Position 2 & \multicolumn{2}{l|}{Le reste de la phrase} \\ \hline
\all{Ich} (sujet) & \all{besuche} (verbe) & \all{morgen meine Großeltern} & \\ \hline
\end{tabular}
\end{center}

\textbf{Cas 2 : l'inversion (un complément d'abord).} Si l'on veut insister sur un complément (de temps, de lieu), on le place en tête de phrase. Mais alors, pour que le verbe reste en deuxième position, \textbf{le sujet passe après le verbe}. C'est ce qu'on appelle l'\cle{inversion}.

\all{Morgen besuche ich meine Großeltern.} « Demain, je rends visite à mes grands-parents. »

\begin{exemplebox}[Inversion : exemples traduits]
\begin{itemize}
\item \all{Am Wochenende feiern wir ein Fest.} « Le week-end, nous fêtons une fête. »
\item \all{Gestern bin ich zu Hause geblieben.} « Hier, je suis resté à la maison. »
\item \all{In Douala wohnen meine Onkel.} « À Douala habitent mes oncles. »
\item \all{Nach dem Essen spielen wir Karten.} « Après le repas, nous jouons aux cartes. »
\item \all{Jetzt lerne ich Deutsch.} « Maintenant, j'apprends l'allemand. »
\end{itemize}
\end{exemplebox}

\begin{attention}[L'erreur n°1 des francophones]
En français, on dit : « Demain, \textbf{je vais} chez mes grands-parents » : le sujet reste avant le verbe. Le francophone traduit donc spontanément : \textbf{✗} \all{Morgen ich gehe zu meinen Großeltern.} C'est FAUX. En allemand, dès qu'un complément ouvre la phrase, le sujet doit passer après le verbe :
\par\medskip
\textbf{✓} \all{Morgen gehe ich zu meinen Großeltern.}
\par\medskip
Autocorrection : après un complément initial, demande-toi toujours : « Où est mon verbe ? » S'il n'est pas en deuxième position, ta phrase est fautive.
\end{attention}

\begin{popfigure}
\begin{tikzpicture}[
  box/.style={draw=popblue, thick, rounded corners=3pt, fill=popblueL, minimum height=1.1cm, align=center, font=\small},
  vbox/.style={draw=poporange, very thick, rounded corners=3pt, fill=poporangeL, minimum height=1.1cm, align=center, font=\small\bfseries},
  fbox/.style={draw=poppurple, thick, rounded corners=3pt, fill=poppurpleL, minimum height=1.1cm, align=center, font=\small},
  posS/.style={font=\footnotesize\bfseries, text=popdark}
]
\node[box, align=center] (p1) at (0,0){Position 1\\ sujet \emph{ou} complément};
\node[vbox, right=0.35cm of p1, align=center] (p2){Position 2\\ VERBE\\ conjugué};
\node[box, right=0.35cm of p2, align=center] (p3){Position 3\\ sujet (si inversion)\\ + compléments};
\node[fbox, right=0.35cm of p3, align=center] (p4){Fin de phrase\\ participe / infinitif\\ / particule};
\node[posS, above=0.12cm of p1] {1};
\node[posS, above=0.12cm of p2] {2};
\node[posS, above=0.12cm of p3] {3};
\node[posS, above=0.12cm of p4] {dernière};
\draw[-{Stealth}, very thick, poporange] (p1.south) to[bend right=25] node[below, font=\footnotesize, text=popdark]{Le sujet glisse en position 3} (p3.south);
\end{tikzpicture}
\legende{Le squelette de la phrase déclarative allemande : la position 2 appartient au verbe conjugué, la fin de phrase à la seconde partie du verbe.}
\end{popfigure}

\courssub{La cadre verbale (Satzklammer)}

Quand le verbe possède deux parties — auxiliaire + participe au \all{Perfekt}, modal + infinitif, verbe à particule séparable —, ces deux parties forment une \cle{cadre verbale} (en allemand \all{Satzklammer}, littéralement « parenthèse de la phrase ») : la première partie se place en position 2, la seconde \textbf{tout à la fin} de la phrase. Tous les compléments se logent \emph{à l'intérieur} de cette cadre.

\begin{exemplebox}[La cadre verbale en action]
\begin{itemize}
\item Perfekt : \all{Wir \textbf{haben} gestern unsere Großeltern \textbf{besucht}.} « Nous avons rendu visite à nos grands-parents hier. »
\item Modal + infinitif : \all{Ich \textbf{muss} morgen meiner Mutter \textbf{helfen}.} « Je dois aider ma mère demain. »
\item Particule séparable : \all{Das Fest \textbf{fängt} um 18 Uhr \textbf{an}.} « La fête commence à 18 heures. »
\item Futur : \all{Wir \textbf{werden} nächste Woche viel \textbf{lachen}.} « Nous rirons beaucoup la semaine prochaine. »
\end{itemize}
\end{exemplebox}

\begin{attention}[Ne rien mettre après la fin de la cadre]
Le francophone a tendance à ajouter un complément après le participe : \textbf{✗} \all{Ich habe besucht meine Großeltern gestern.} En allemand, le participe (ou l'infinitif, ou la particule) ferme la phrase : \textbf{✓} \all{Ich habe gestern meine Großeltern besucht.} Rien ne se place après, sauf rares compléments introduits par une préposition dans un style soutenu.
\end{attention}

\courssub{L'ordre des compléments : TE-KA-MO-LO}

À l'intérieur de la cadre, quand plusieurs compléments se succèdent, l'allemand suit un ordre préférentiel que l'on retient grâce au mot-valise \cle{TE-KA-MO-LO} :

\begin{center}
\begin{tabularx}{\linewidth}{|l|l|X|X|}\hline
\rowcolor{popblueL} Sigle & Type de complément & Question posée & Exemple \\ \hline
\textbf{TE} & Temporal (temps) & \all{Wann?} (quand ?) & \all{morgen, heute, am Abend} \\ \hline
\textbf{KA} & Kausal (cause) & \all{Warum?} (pourquoi ?) & \all{wegen des Wetters} \\ \hline
\textbf{MO} & Modal (manière) & \all{Wie?} (comment ?) & \all{mit dem Bus, gern} \\ \hline
\textbf{LO} & Lokal (lieu) & \all{Wohin?/Wo?} (où ?) & \all{nach Douala, zu Hause} \\ \hline
\end{tabularx}
\end{center}

\all{Ich fahre \textbf{morgen} (TE) \textbf{mit meinem Bruder} (MO) \textbf{nach Douala} (LO).} « Je vais demain avec mon frère à Douala. »

\all{Wir feiern \textbf{am Samstag} (TE) \textbf{wegen des Geburtstags} (KA) \textbf{mit der ganzen Familie} (MO) \textbf{im Garten} (LO).} « Nous fêtons samedi, à cause de l'anniversaire, avec toute la famille, dans le jardin. »

Remarque : le français préfère souvent l'ordre inverse (lieu, manière, temps) ou place le temps en tête ; ne calque jamais l'ordre français !

\courssub{La place des pronoms}

Règle simplifiée pour le niveau Première : \textbf{les pronoms personnels se placent juste après le verbe conjugué}, avant les noms communs et les compléments longs.

\all{Ich gebe \textbf{ihm} das Geschenk.} « Je lui donne le cadeau. »

Quand deux pronoms se suivent, l'ordre est : \textbf{accusatif avant datif} (contrairement à deux noms communs, où le datif précède l'accusatif).

\begin{center}
\begin{tabularx}{\linewidth}{|X|X|X|}\hline
\rowcolor{popblueL} Français & Deutsch & Remarque \\ \hline
Je le lui donne. & \all{Ich gebe es ihm.} & pronom acc. \all{es} avant pronom dat. \all{ihm} \\ \hline
Je donne le livre à mon frère. & \all{Ich gebe meinem Bruder das Buch.} & nom au datif avant nom à l'accusatif \\ \hline
Je le donne à mon frère. & \all{Ich gebe es meinem Bruder.} & le pronom passe avant le nom \\ \hline
Je le lui montre. & \all{Ich zeige es ihr.} & deux pronoms : acc. puis dat. \\ \hline
\end{tabularx}
\end{center}

\courssub{L'ordre des mots dans les questions}

\textbf{Question sans mot interrogatif} (réponse oui/non) : le verbe conjugué occupe la \textbf{première position}, exactement comme l'inversion sujet-verbe du français « Viens-tu ? ».

\all{\textbf{Kommst} du mit?} « Viens-tu avec (nous) ? »

\all{\textbf{Hast} du deine Hausaufgaben gemacht?} « As-tu fait tes devoirs ? » (cadre verbale : \all{hast} en position 1, \all{gemacht} à la fin).

\textbf{Question avec mot interrogatif} (\all{wann, wo, wer, wie, warum, was}...) : le mot interrogatif occupe la position 1, le verbe reste en \textbf{position 2}, puis vient le sujet. C'est une application directe de la règle V2 !

\all{\textbf{Wann} \textbf{besuchst} du deine Familie?} « Quand rends-tu visite à ta famille ? »

\all{\textbf{Wo} \textbf{wohnen} deine Großeltern?} « Où habitent tes grands-parents ? »

\all{\textbf{Warum} \textbf{bist} du gestern zu Hause geblieben?} « Pourquoi es-tu resté à la maison hier ? »

\begin{center}
\begin{tabularx}{\linewidth}{|l|X|X|}\hline
\rowcolor{popblueL} Type de phrase & Position 1 & Position 2 \\ \hline
Déclarative standard & sujet : \all{Ich} & verbe : \all{komme} \\ \hline
Déclarative avec inversion & complément : \all{Morgen} & verbe : \all{komme} (puis sujet) \\ \hline
Question oui/non & \textbf{verbe} : \all{Kommst} & sujet : \all{du} \\ \hline
Question avec mot interrogatif & \all{Wann} & verbe : \all{kommst} (puis sujet) \\ \hline
\end{tabularx}
\end{center}

\courssub{Comparaison français / allemand}

\begin{center}
\begin{tabularx}{\linewidth}{|X|X|}\hline
\rowcolor{popgreenL} Français & Allemand \\ \hline
Le sujet précède presque toujours le verbe : « Hier, \textbf{je suis resté} à la maison. » & Le verbe conjugué est toujours en 2\textsuperscript{e} position : \all{\textbf{Gestern bin ich} zu Hause geblieben.} \\ \hline
L'auxiliaire et le participe sont voisins : « j'\textbf{ai visité} mes grands-parents ». & L'auxiliaire et le participe encadrent la phrase : \all{Ich \textbf{habe} meine Großeltern \textbf{besucht}.} \\ \hline
Ordre des compléments libre, souvent lieu-manière-temps. & Ordre préférentiel TE-KA-MO-LO : temps, cause, manière, lieu. \\ \hline
Question oui/non : intonation ou « est-ce que ». & Question oui/non : verbe en position 1, \all{Kommst du?} \\ \hline
\end{tabularx}
\end{center}

\begin{methode}[Vérifier l'ordre des mots en quatre étapes]
\begin{enumerate}
\item \textbf{Compter les positions} : identifie le premier bloc (un seul élément : sujet ou complément) = position 1.
\item \textbf{Placer le verbe conjugué} en position 2, immédiatement après, sans rien intercaler.
\item \textbf{Placer le sujet} juste après le verbe si la phrase commence par un complément (inversion), puis les compléments dans l'ordre TE-KA-MO-LO, les pronoms d'abord.
\item \textbf{Fermer la cadre} : participe passé, infinitif ou particule séparable à la toute fin. Relis : rien ne doit se trouver après.
\end{enumerate}
\end{methode}

\begin{exoresolu}[Remets la phrase dans l'ordre]
Énoncé : remets les éléments suivants dans le bon ordre pour former une phrase déclarative correcte : \all{meine Tante / am Sonntag / besucht / haben / wir / in Douala}.
\tcblower
Solution détaillée :
\begin{enumerate}
\item Position 1 : choisissons le complément de temps pour insister : \all{Am Sonntag}.
\item Position 2 : le verbe conjugué (l'auxiliaire du Perfekt) : \all{haben}.
\item Inversion : le sujet suit le verbe : \all{wir}. Puis les compléments : l'objet accusatif \all{meine Tante} (personne visitée) avant le lieu \all{in Douala} (ordre standard : compléments d'objet avant compléments de lieu).
\item Fin de phrase : le participe passé ferme la cadre : \all{besucht}.
\end{enumerate}
\textbf{Phrase correcte (ordre neutre) :} \all{Am Sonntag haben wir meine Tante in Douala besucht.} « Dimanche, nous avons rendu visite à ma tante à Douala. »
\par
Variante également correcte (sujet d'abord) : \all{Wir haben am Sonntag meine Tante in Douala besucht.} En revanche, \all{Am Sonntag wir haben...} est faux : l'inversion a été oubliée.
\end{exoresolu}

\begin{exercice}[À toi de jouer]
\begin{enumerate}
\item Traduis en allemand : « Hier, j'ai aidé ma mère à la maison. »
\item Remets dans l'ordre : \all{mit dem Bus / fährt / morgen / nach Bamenda / mein Bruder}.
\item Forme une question avec \all{wann} : « Quand fêtez-vous l'anniversaire de votre grand-mère ? »
\item Corrige la faute : \all{Am Abend wir hören Musik.}
\end{enumerate}
\end{exercice}

\begin{corrige}[Exercice : À toi de jouer]
\begin{enumerate}
\item \all{Gestern habe ich meiner Mutter zu Hause geholfen.} (inversion après \all{Gestern} ; \all{meiner Mutter} au datif car \all{helfen} exige le datif ; participe \all{geholfen} en fin de phrase.)
\item \all{Mein Bruder fährt morgen mit dem Bus nach Bamenda.} (ordre TE : \all{morgen} — MO : \all{mit dem Bus} — LO : \all{nach Bamenda}.) Variante : \all{Morgen fährt mein Bruder mit dem Bus nach Bamenda.}
\item \all{Wann feiert ihr den Geburtstag eurer Großmutter?} (mot interrogatif en position 1, verbe \all{feiert} en position 2, sujet \all{ihr} ensuite.)
\item \all{Am Abend hören wir Musik.} Le complément \all{Am Abend} ouvre la phrase : le verbe \all{hören} doit passer en position 2 et le sujet \all{wir} après le verbe.
\end{enumerate}
\end{corrige}

\begin{attention}[Récapitulatif des pièges]
\begin{itemize}
\item Oublier l'inversion après un complément initial : ✗ \all{Morgen ich gehe...} → ✓ \all{Morgen gehe ich...}
\item Placer le participe juste après l'auxiliaire comme en français : ✗ \all{Ich habe besucht meine Tante.} → ✓ \all{Ich habe meine Tante besucht.}
\item Oublier la particule séparable en fin de phrase : ✗ \all{Der Zug kommt um 8 Uhr.} (pour « arriver ») → ✓ \all{Der Zug kommt um 8 Uhr an.}
\item Commencer une question oui/non par le sujet : ✗ \all{Du kommst mit?} (possible à l'oral avec l'intonation, mais à éviter à l'écrit) → ✓ \all{Kommst du mit?}
\end{itemize}
\end{attention}

\begin{savaistu}[L'ordre des mots au baccalauréat]
Au Bac, en épreuve de thème (traduction français → allemand), l'ordre des mots est un critère majeur de notation : une inversion oubliée après un complément initial coûte des points, même si le vocabulaire, les cas et la conjugaison sont parfaits. Les correcteurs repèrent immédiatement la phrase qui commence par un complément suivi du sujet : c'est la « signature » du francophone qui ne maîtrise pas encore la syntaxe allemande. Entraîne-toi donc, pour chaque phrase de thème, à compter à voix basse : « un... deux : verbe ! ». Ce réflexe, acquis en Première, fera la différence en Terminale.
\end{savaistu}

\coursec{Production : lettre, description, dialogue}

Dans la section précédente, nous avons révisé l'ordre des mots : le verbe conjugué en deuxième position, le participe et l'infinitif en fin de phrase. Cette grammaire n'a de sens que si elle sert à \textbf{communiquer}. Nous allons maintenant l'utiliser pour produire trois types de textes très fréquents à l'examen et dans la vie réelle : la \cle{courte lettre personnelle}, la \cle{description de famille} et le \cle{dialogue}. Ces trois genres tournent autour du même thème : la vie en famille et en société.

\courssub{Observer d'abord : une lettre modèle}

Lisons d'abord une lettre authentique. Paul, un élève de Première à Yaoundé, répond à son amie allemande Lena.

\begin{textebox}[Brief von Paul an Lena]
Yaoundé, den 15. Mai

Liebe Lena,

vielen Dank für deinen Brief. Bei uns ist alles gut. Letzten Sonntag haben wir den Geburtstag meines Vaters gefeiert. Die ganze Familie ist gekommen und wir haben zusammen gegessen und getanzt. Meine Mutter hat einen großen Kuchen gebacken und mein Onkel hat Musik gemacht. Es war sehr schön!

Und du, was hast du am Wochenende gemacht? Schreib mir bald!

Viele Grüße

Paul

\emph{Glossaire :} der Brief, -e (la lettre) ; der Geburtstag, -e (l'anniversaire) ; feiern (fêter) ; der Kuchen, - (le gâteau) ; backen (cuire, faire un gâteau) ; der Onkel, - (l'oncle) ; bald (bientôt).
\end{textebox}

\textbf{Questions de compréhension :}
\begin{enumerate}
\item Wer schreibt den Brief und an wen? (Qui écrit la lettre et à qui ?)
\item Was hat die Familie gefeiert? (Qu'est-ce que la famille a fêté ?)
\item Welche Zeitform benutzt Paul für das Fest? (Quel temps Paul utilise-t-il pour raconter la fête ?)
\end{enumerate}

\emph{Réponses :} 1. Paul schreibt an Lena. 2. Den Geburtstag seines Vaters. 3. Das Perfekt (\all{haben gefeiert}, \all{ist gekommen}, \all{haben gegessen}).

\courssub{La courte lettre personnelle : structure et règles}

\begin{propriete}[Les 5 blocs de la lettre personnelle]
Une lettre personnelle allemande (\all{der Brief}) suit toujours le même plan :
\begin{enumerate}
\item \cle{Lieu et date} en haut à droite : \all{Yaoundé, den 15. Mai} (accusatif avec \all{den} + article, sans préposition).
\item \cle{Salutation} suivie d'une virgule : \all{Liebe Lena,} (à une fille) / \all{Lieber Paul,} (à un garçon).
\item \cle{Einleitung} (introduction) : remerciement ou nouvelles — \all{vielen Dank für deinen Brief. Bei uns ist alles gut.}
\item \cle{Hauptteil} (corps) : on raconte l'événement, souvent au parfait.
\item \cle{Schluss und Gruß} (conclusion et formule de politesse) : \all{Schreib mir bald! Viele Grüße} ou \all{Herzliche Grüße}, puis la signature sans virgule.
\end{enumerate}
\end{propriete}

\begin{popfigure}
\begin{center}
\begin{tikzpicture}[node distance=0.45cm,
block/.style={rectangle, rounded corners=3pt, draw=popblue, fill=popblueL, text width=10.5cm, align=left, inner sep=5pt, font=\small}]
\node[block] (a) {\textbf{1. Ort und Datum :} Yaoundé, den 15. Mai};
\node[block, below=of a] (b) {\textbf{2. Anrede :} Liebe Lena, / Lieber Paul,};
\node[block, below=of b, fill=popgreenL, draw=popgreen] (c) {\textbf{3. Einleitung :} vielen Dank für deinen Brief. Bei uns ist alles gut.};
\node[block, below=of c, fill=poporangeL, draw=poporange] (d) {\textbf{4. Hauptteil :} Letzten Sonntag haben wir ... gefeiert. (Perfekt !)};
\node[block, below=of d, fill=poppinkL, draw=poppink] (e) {\textbf{5. Schluss und Gruß :} Schreib mir bald! Viele Grüße — Paul};
\draw[-{Stealth}, thick, popdark] (a) -- (b);
\draw[-{Stealth}, thick, popdark] (b) -- (c);
\draw[-{Stealth}, thick, popdark] (c) -- (d);
\draw[-{Stealth}, thick, popdark] (d) -- (e);
\end{tikzpicture}
\end{center}
\legende{Plan-type de la lettre personnelle en cinq blocs}
\end{popfigure}

\begin{attention}[La minuscule après la salutation]
En allemand, la salutation est suivie d'une \textbf{virgule}, et la ligne suivante commence par une \textbf{minuscule} — contrairement au français où l'on met une majuscule après « Chère Lena, ».

\all{Lieber Paul,}\\
\all{ich danke dir für deine Einladung.} « Cher Paul, je te remercie de ton invitation. »

Écrire \all{Lieber Paul, Ich danke dir...} avec majuscule est une faute classique du francophone. Seule exception : le pronom de politesse \all{Sie} garde toujours sa majuscule.
\end{attention}

\begin{exemplebox}[Formules utiles pour la lettre]
\all{Vielen Dank für deinen Brief.} « Merci beaucoup pour ta lettre. »\\
\all{Wie geht es dir? Mir geht es gut.} « Comment vas-tu ? Je vais bien. »\\
\all{Ich habe eine tolle Nachricht!} « J'ai une super nouvelle ! »\\
\all{Schreib mir bald!} « Écris-moi vite ! »\\
\all{Viele Grüße / Herzliche Grüße / Liebe Grüße} « Amitiés / Bien cordialement / Grosses bises »
\end{exemplebox}

\courssub{La description de famille}

Pour décrire sa famille, on présente chaque membre selon quatre axes : \textbf{l'identité et l'âge}, \textbf{le physique}, \textbf{le caractère}, \textbf{les activités}. On utilise surtout le présent.

\begin{center}
\begin{tabularx}{\linewidth}{|X|X|X|}
\hline
\rowcolor{popblueL} \textbf{Axe} & \textbf{Deutsch} & \textbf{Français} \\
\hline
Identité, âge & \all{Mein Vater ist 45 Jahre alt.} & Mon père a 45 ans. \\
\hline
Physique & \all{Er ist groß und hat kurze Haare.} & Il est grand et a les cheveux courts. \\
\hline
Caractère & \all{Sie ist freundlich und lustig.} & Elle est gentille et drôle. \\
\hline
Activité & \all{Er arbeitet als Lehrer.} & Il travaille comme professeur. \\
\hline
Loisir & \all{Mein Bruder spielt gern Fußball.} & Mon frère aime jouer au football. \\
\hline
\end{tabularx}
\end{center}

\begin{exemplebox}[Une courte description modèle]
\all{Meine Familie ist nicht groß. Mein Vater ist 45 Jahre alt. Er ist groß und freundlich. Er arbeitet als Lehrer in Douala. Meine Mutter ist 40 Jahre alt. Sie ist klein und sehr nett. Sie kocht sehr gern. Meine Schwester heißt Amina und ist 12 Jahre alt. Sie ist lustig und tanzt gern.} « Ma famille n'est pas grande. Mon père a 45 ans. Il est grand et gentil. Il travaille comme professeur à Douala. Ma mère a 40 ans. Elle est petite et très gentille. Elle aime beaucoup cuisiner. Ma sœur s'appelle Amina et a 12 ans. Elle est drôle et aime danser. »
\end{exemplebox}

\begin{attention}[Faux amis et pièges de la description]
\begin{itemize}
\item \all{groß} = grand (taille), pas « gros » ; « gros » se dit \all{dick}.
\item L'âge : \all{Er ist 45 Jahre alt} (verbe \all{sein}, pas \all{haben} comme le français « il a 45 ans »).
\item \all{arbeiten als} + métier sans article : \all{Er arbeitet als Lehrer} (jamais \all{als ein Lehrer}).
\end{itemize}
\end{attention}

\courssub{Le dialogue : inviter, accepter, refuser, féliciter}

Le dialogue est un échange oral reconstitué à l'écrit. Il faut maîtriser les formules de base de la vie sociale.

\begin{center}
\begin{tabularx}{\linewidth}{|X|X|X|}
\hline
\rowcolor{popblueL} \textbf{Fonction} & \textbf{Deutsch} & \textbf{Français} \\
\hline
Inviter & \all{Kommst du zu meiner Party?} & Tu viens à ma fête ? \\
\hline
Inviter & \all{Hast du Lust, ins Kino zu gehen?} & Ça te dit d'aller au cinéma ? \\
\hline
Accepter & \all{Ja, gern! / Gute Idee!} & Oui, avec plaisir ! / Bonne idée ! \\
\hline
Refuser & \all{Leider kann ich nicht kommen.} & Malheureusement, je ne peux pas venir. \\
\hline
Féliciter & \all{Herzlichen Glückwunsch zum Geburtstag!} & Joyeux anniversaire ! \\
\hline
Remercier & \all{Danke schön! / Vielen Dank!} & Merci beaucoup ! \\
\hline
\end{tabularx}
\end{center}

\begin{exemplebox}[Mini-dialogue modèle]
\all{Lena: Hallo Paul! Kommst du am Samstag zu meiner Party?} « Lena : Salut Paul ! Tu viens à ma fête samedi ? »\\
\all{Paul: Ja, gern! Wann beginnt sie?} « Paul : Oui, avec plaisir ! À quelle heure ça commence ? »\\
\all{Lena: Um 18 Uhr. Bringst du deine Schwester mit?} « Lena : À 18 heures. Tu amènes ta sœur ? »\\
\all{Paul: Leider kann sie nicht kommen, sie ist krank.} « Paul : Malheureusement elle ne peut pas venir, elle est malade. »\\
\all{Lena: Oh, schade! Dann bis Samstag!} « Lena : Oh, dommage ! Alors à samedi ! »
\end{exemplebox}

\begin{savaistu}[Du oder Sie ?]
En Allemagne, on \textbf{tutoie} (\all{du}) la famille, les amis et les jeunes entre eux ; le \textbf{vouvoiement} (\all{Sie}, toujours avec majuscule, même au milieu d'une phrase) est réservé aux inconnus, aux adultes et aux situations formelles. Dans un dialogue entre amis, écrire \all{Kommen Sie zu meiner Party?} est une erreur de registre : on dit \all{Kommst du zu meiner Party?}. À l'inverse, un élève qui s'adresse à son directeur dira \all{Sie}. Bien choisir le registre fait partie de la compétence de communication !
\end{savaistu}

\courssub{Les connecteurs pour enrichir la production}

Une bonne production n'est pas une suite de phrases isolées : les phrases doivent être \textbf{reliées}.

\begin{center}
\begin{tabularx}{\linewidth}{|X|X|X|}
\hline
\rowcolor{popblueL} \textbf{Connecteur} & \textbf{Français} & \textbf{Exemple} \\
\hline
\all{und} & et & \all{Wir haben gegessen und getanzt.} (Nous avons mangé et dansé.) \\
\hline
\all{aber} & mais & \all{Ich war müde, aber glücklich.} (J'étais fatigué mais heureux.) \\
\hline
\all{denn} & car & \all{Ich komme nicht, denn ich bin krank.} (Je ne viens pas car je suis malade.) \\
\hline
\all{deshalb} & c'est pourquoi & \all{Ich bin krank, deshalb bleibe ich zu Hause.} (Je suis malade, c'est pourquoi je reste à la maison.) \\
\hline
\all{zuerst} & d'abord & \all{Zuerst haben wir gegessen.} (D'abord nous avons mangé.) \\
\hline
\all{dann / danach} & puis / ensuite & \all{Dann haben wir getanzt.} (Puis nous avons dansé.) \\
\hline
\all{schließlich} & finalement & \all{Schließlich sind alle nach Hause gegangen.} (Finalement tout le monde est rentré.) \\
\hline
\end{tabularx}
\end{center}

\begin{attention}[deshalb et la place du verbe]
\all{deshalb} occupe la première position : le verbe conjugué vient juste après, en position 2, et le sujet passe derrière : \all{Ich bin krank, deshalb bleibe ich zu Hause.} Ne dites jamais \all{deshalb ich bleibe} (calque du français « c'est pourquoi je reste »). Rappel : \all{und}, \all{aber}, \all{denn} ne comptent pas comme position : après eux, l'ordre reste sujet + verbe.
\end{attention}

\courssub{Méthode de production guidée en 4 étapes}

\begin{methode}[Réussir une courte lettre au Bac]
\begin{enumerate}
\item \textbf{Brainstormer le lexique} : sur le brouillon, lister 8 à 10 mots du chapitre (famille, fête, sentiments) que l'on veut réinvestir : \all{der Geburtstag, feiern, die Familie, der Kuchen, tanzen, schön...}
\item \textbf{Rédiger des phrases simples mais correctes} : 5 à 8 phrases, verbe en position 2, au moins 2 phrases au parfait avec le bon auxiliaire (\all{haben} ou \all{sein}).
\item \textbf{Relier avec des connecteurs} : ajouter \all{und}, \all{aber}, \all{zuerst}, \all{dann}, \all{deshalb} pour donner du rythme au texte.
\item \textbf{Relire avec la check-list} (voir ci-dessous) et respecter le plan en 5 blocs : lieu/date, salutation + virgule, introduction, corps, conclusion + formule de politesse.
\end{enumerate}
\end{methode}

\begin{aretenir}[Check-list de relecture]
Avant de rendre votre copie, vérifiez :
\begin{itemize}
\item[$\square$] Majuscule à \textbf{tous les noms} : \all{die Familie, der Brief, das Fest}.
\item[$\square$] Verbe conjugué en \textbf{position 2} dans chaque phrase principale.
\item[$\square$] Participe passé et infinitif \textbf{en fin de phrase} : \all{Wir haben getanzt.} / \all{Ich kann nicht kommen.}
\item[$\square$] Bon auxiliaire au parfait : \all{sein} pour les verbes de mouvement (\all{ist gekommen}), \all{haben} pour les autres (\all{hat gebacken}).
\item[$\square$] Cas corrects après prépositions et verbes à datif : \all{für deinen Brief} (accusatif), \all{ich danke dir} (datif), \all{zu meiner Party} (datif).
\item[$\square$] Minuscule après la salutation, virgule après \all{Liebe Lena}.
\end{itemize}
\end{aretenir}

\courssub{Commentaire modèle d'un texte}

À l'examen, on peut aussi vous demander de \textbf{commenter} un texte. La démarche comporte cinq points : le type de texte, le thème, la structure, les temps employés, le ton.

\begin{exemplebox}[Commentaire du texte de Lena]
\emph{Type de texte :} il s'agit d'une lettre personnelle (\all{ein persönlicher Brief}) : on reconnaît le lieu et la date, la salutation \all{Liebe Lena} et la formule finale \all{Viele Grüße}.

\emph{Thème :} la vie familiale — Paul raconte la fête d'anniversaire de son père.

\emph{Structure :} cinq blocs classiques (date, salutation, introduction avec remerciement, corps du récit, conclusion avec invitation à répondre).

\emph{Temps employés :} le présent pour les nouvelles actuelles (\all{Bei uns ist alles gut}) et le parfait pour raconter les événements passés (\all{haben wir gefeiert}, \all{ist gekommen}, \all{hat gebacken}).

\emph{Ton :} chaleureux et familier — tutoiement (\all{deinen Brief}, \all{Schreib mir bald}), exclamation \all{Es war sehr schön!}.
\end{exemplebox}

\courssub{Exercice résolu et entraînement}

\begin{exoresolu}[Remettre une lettre en ordre]
\textbf{Énoncé.} Remettez les éléments suivants dans le bon ordre pour reconstituer la lettre de Paul :
(a) \all{Viele Grüße — Paul} ; (b) \all{vielen Dank für deinen Brief. Bei uns ist alles gut.} ; (c) \all{Yaoundé, den 15. Mai} ; (d) \all{Letzten Sonntag haben wir den Geburtstag meines Vaters gefeiert.} ; (e) \all{Liebe Lena,} ; (f) \all{Schreib mir bald!}
\tcblower
\textbf{Solution détaillée.} On applique le plan en 5 blocs :
\begin{enumerate}
\item Lieu et date d'abord : (c) \all{Yaoundé, den 15. Mai}.
\item Salutation avec virgule : (e) \all{Liebe Lena,}.
\item Introduction (remerciement) : (b) — attention, elle commence par une minuscule : \all{vielen Dank...}.
\item Corps du récit au parfait : (d) \all{Letzten Sonntag haben wir ... gefeiert.} (participe \all{gefeiert} bien en fin de phrase).
\item Conclusion : (f) \all{Schreib mir bald!} puis formule de politesse et signature : (a) \all{Viele Grüße — Paul}.
\end{enumerate}
Ordre correct : \textbf{c – e – b – d – f – a}.
\end{exoresolu}

\begin{exercice}[À vous de produire]
\begin{enumerate}
\item \textbf{Lettre.} Écrivez une courte lettre (6 à 8 phrases) à votre ami allemand Thomas : remerciez-le pour sa lettre, racontez une fête de famille à Yaoundé (au moins 2 phrases au parfait) et posez-lui une question. Respectez le plan en 5 blocs.
\item \textbf{Description.} Décrivez deux membres de votre famille (âge, physique, caractère, activité) en 4 phrases.
\item \textbf{Dialogue.} Complétez : \all{A: Kommst du morgen zu meiner Geburtstagsparty? B: Ja, ...! Wann beginnt sie? A: Um 17 Uhr. B: Gut, dann bis morgen!}
\end{enumerate}
\end{exercice}

\begin{corrige}[Exercice : éléments de correction]
\begin{enumerate}
\item Modèle possible : \all{Yaoundé, den 20. Mai. Lieber Thomas, vielen Dank für deinen Brief. Mir geht es gut. Letzten Samstag haben wir eine große Party gemacht. Meine ganze Familie ist gekommen und wir haben zusammen gegessen. Es war sehr schön! Und du, was hast du gemacht? Schreib mir bald! Viele Grüße, [prénom].} — Vérifier : minuscule après la salutation, auxiliaire \all{sein} avec \all{kommen}, participes en fin de phrase.
\item Modèle : \all{Mein Bruder ist 20 Jahre alt. Er ist groß und sportlich. Er ist sehr lustig. Er studiert in Yaoundé.}
\item \all{B: Ja, gern! Wann beginnt sie?}
\end{enumerate}
\end{corrige}

En résumé : une production réussie repose sur un \textbf{plan clair}, des \textbf{phrases simples et correctes}, des \textbf{connecteurs} et une \textbf{relecture méthodique}. Ces outils vous serviront directement dans la section suivante, consacrée à la traduction — thème et version — où la rigueur grammaticale sera encore plus décisive.

\coursec{Traduction : thème et version}

Dans la section précédente, tu as appris à \emph{produire} des textes courts : une lettre, une description, un dialogue. La traduction est l'étape suivante logique : il s'agit maintenant de passer d'une langue à l'autre en gardant le sens, la grammaire et le naturel. Au baccalauréat comme dans la vie professionnelle, savoir traduire fidèlement un court texte sur la vie familiale et sociale est une compétence précieuse. Distinguons d'abord les deux directions.

\begin{definition}[Thème et version]
La \cle{version} est la traduction de l'\textbf{allemand vers le français} : tu lis un texte allemand et tu le restitues en français correct.
Le \cle{thème} est la traduction du \textbf{français vers l'allemand} : tu pars d'une phrase française et tu la reconstruis en allemand correct.
Dans les deux cas, on traduit des \emph{idées}, jamais des mots isolés.
\end{definition}

\courssub{La version : de l'allemand vers le français}

\textbf{Observation.} Lis d'abord cette phrase extraite du texte de Lena (section 1) :

\all{Gestern hat meine Familie das Fest im Garten gefeiert.}

Un élève qui traduit mot à mot obtient : « Hier a ma famille la fête dans le jardin fêté. » Ce n'est pas du français ! Le bon réflexe consiste à repérer d'abord le verbe conjugué (\all{hat ... gefeiert}, Perfekt), puis le sujet (\all{meine Familie}), puis à reconstruire : « Hier, ma famille a fêté la fête dans le jardin. »

\begin{methode}[Les 5 étapes de la version]
\begin{enumerate}
\item \textbf{Lire tout le texte} une première fois, sans traduire, pour comprendre le sens global (qui ? quoi ? où ? quand ?).
\item \textbf{Repérer le verbe conjugué} (2\textsuperscript{e} position) et son \textbf{sujet}. Attention : le participe ou l'infinitif se trouve en fin de phrase.
\item \textbf{Identifier les temps} : Präsens $\rightarrow$ présent français ; Perfekt (\all{hat gemacht}, \all{ist gekommen}) $\rightarrow$ passé composé.
\item \textbf{Traduire par blocs de sens} (sujet + verbe + compléments), pas mot à mot.
\item \textbf{Rédiger un français correct et naturel} : on doit pouvoir lire ta phrase à voix haute sans qu'elle « sonne allemand ».
\end{enumerate}
\end{methode}

\begin{exemplebox}[Exemples de version]
\all{Meine Eltern haben uns zum Essen eingeladen.} $\rightarrow$ « Mes parents nous ont invités à manger. » (verbe à particule \all{einladen} : le participe \all{eingeladen} se traduit simplement par « invités ».)

\all{Am Wochenende besuchen wir unsere Großeltern in Bamako.} $\rightarrow$ « Le week-end, nous rendons visite à nos grands-parents à Bamako. » (\all{am Wochenende} = « le week-end » ; inversion allemande sans conséquence en français.)

\all{Lena ist nach dem Fest zu Hause geblieben.} $\rightarrow$ « Lena est restée à la maison après la fête. » (auxiliaire \all{sein} : verbe de déplacement ou d'état.)
\end{exemplebox}

\begin{attention}[Ne jamais laisser de mot sans traduction]
En version, \textbf{chaque mot allemand doit avoir un écho} dans ta phrase française. Si tu « oublies » \all{gestern}, tu perds un point. Vérifie en comparant : nombre de compléments de temps, de lieu, de personnes. En revanche, ne traduis jamais mot à mot les structures : \all{es gibt} = « il y a », \all{zum Essen} = « à manger / pour le repas ».
\end{attention}

\courssub{Le thème : du français vers l'allemand}

Le thème est plus difficile : c'est toi qui \emph{construis} la phrase allemande. Il faut donc appliquer consciemment toutes les règles révisées dans ce chapitre (présent, parfait, ordre des mots).

\begin{methode}[Les 5 étapes du thème]
\begin{enumerate}
\item \textbf{Analyser la phrase française} : qui est le sujet ? quel est le verbe ? quels compléments ?
\item \textbf{Choisir le temps} : présent $\rightarrow$ Präsens ; passé composé $\rightarrow$ Perfekt ; futur proche $\rightarrow$ Präsens + complément de temps.
\item \textbf{Placer le verbe conjugué en position 2} : si la phrase commence par un complément (\all{hier}, \all{morgen}), le sujet passe après le verbe (inversion).
\item \textbf{Vérifier cas et auxiliaire} : accusatif ou datif après le verbe et les prépositions ; \all{haben} ou \all{sein} au Perfekt ; participe en fin de phrase.
\item \textbf{Relire} : majuscules aux noms, orthographe (ä, ö, ü, ß), place du participe.
\end{enumerate}
\end{methode}

\begin{popfigure}
\begin{tikzpicture}[
  node distance=0.55cm,
  etape/.style={rectangle, rounded corners=4pt, draw=popblue, fill=popblueL, text width=4.6cm, align=center, font=\small, minimum height=0.9cm},
  fleche/.style={-{Stealth[length=2.5mm]}, thick, poporange}
]
\node[etape, align=center] (e1){\textbf{1. Analyser}\\ sujet, verbe, compléments};
\node[etape, below=of e1, align=center] (e2){\textbf{2. Choisir le temps}\\ Präsens ou Perfekt};
\node[etape, below=of e2, align=center] (e3){\textbf{3. Positionner le verbe}\\ verbe conjugué en 2\textsuperscript{e} position};
\node[etape, below=of e3, align=center] (e4){\textbf{4. Vérifier cas et auxiliaire}\\ Akkusativ/Dativ, haben/sein};
\node[etape, below=of e4, align=center] (e5){\textbf{5. Relire}\\ majuscules, orthographe, participe final};
\draw[fleche] (e1) -- (e2);
\draw[fleche] (e2) -- (e3);
\draw[fleche] (e3) -- (e4);
\draw[fleche] (e4) -- (e5);
\end{tikzpicture}
\legende{Organigramme : les 5 étapes du thème (français $\rightarrow$ allemand)}
\end{popfigure}

\begin{propriete}[Correspondance des temps français / allemand]
\begin{center}
\begin{tabularx}{\linewidth}{|X|X|X|}
\hline
\rowcolor{popblueL} Français & Deutsch & Exemple \\
\hline
Présent & Präsens & « J'aide ma mère. » $\rightarrow$ \all{Ich helfe meiner Mutter.} \\
\hline
Passé composé & Perfekt & « Nous avons fêté. » $\rightarrow$ \all{Wir haben gefeiert.} \\
\hline
Imparfait (description) & Präteritum de \all{sein}/\all{haben} (sinon Perfekt) & « La fête était belle. » $\rightarrow$ \all{Das Fest war schön.} \\
\hline
Futur proche & Präsens + complément de temps & « Demain, j'invite mes amis. » $\rightarrow$ \all{Morgen lade ich meine Freunde ein.} \\
\hline
\end{tabularx}
\end{center}
L'allemand n'a pas d'imparfait de narration obligatoire au niveau A2 : au passé, utilise le \textbf{Perfekt}, sauf pour \all{sein} (\all{war}) et \all{haben} (\all{hatte}), dont le Präteritum est très courant.
\end{propriete}

\begin{exemplebox}[Exemple de thème commenté]
« Ma sœur a reçu beaucoup de cadeaux. »
\begin{itemize}
\item Étape 1 : sujet = « ma sœur » ; verbe = « a reçu » (passé composé) ; complément = « beaucoup de cadeaux ».
\item Étape 2 : passé composé $\rightarrow$ Perfekt.
\item Étape 3 : sujet en tête, verbe en position 2 : \all{Meine Schwester hat ...}
\item Étape 4 : auxiliaire \all{haben} ; participe de \all{bekommen} = \all{bekommen} ; \all{viele Geschenke} à l'accusatif pluriel.
\item Étape 5 : \all{Meine Schwester hat viele Geschenke bekommen.}
\end{itemize}
\end{exemplebox}

\begin{attention}[Les pièges du thème]
\begin{itemize}
\item \textbf{L'inversion} : « hier » en tête de phrase déclenche l'inversion : « Hier, nous sommes restés à la maison. » $\rightarrow$ \all{Hier sind wir zu Hause geblieben.} (et non ✗ \all{Hier wir sind...}).
\item \textbf{Les verbes pronominaux} : « Nous nous entendons bien avec nos voisins. » $\rightarrow$ \all{Wir verstehen uns gut mit unseren Nachbarn.} — ne pas oublier le pronom réfléchi \all{uns}.
\item \textbf{Les verbes à datif} : \all{helfen}, \all{danken}, \all{gefallen} : « J'aide mon frère. » $\rightarrow$ \all{Ich helfe meinem Bruder.} (datif !).
\item \textbf{Ne jamais inventer une structure inconnue} : si tu ne sais pas dire une phrase, \emph{reformule} avec ce que tu sais. « Ma tante a préparé un repas délicieux » peut devenir \all{Meine Tante hat sehr gut gekocht.}
\end{itemize}
\end{attention}

\begin{attention}[Les faux-amis à connaître par cœur]
\begin{center}
\begin{tabularx}{\linewidth}{|X|X|X|}
\hline
\rowcolor{popblueL} Mot allemand & Vrai sens français & Piège \\
\hline
\all{bekommen} & recevoir, obtenir & « devenir » = \all{werden} \\
\hline
\all{aktuell} & actuel, en ce moment & « actuel » se dit bien \all{aktuell}, mais « actually » ≠ \all{aktuell} \\
\hline
\all{der Chef} & le patron, le chef d'entreprise & « le chef » (cuisinier) = \all{der Koch} \\
\hline
\all{die Mappe} & la pochette, le classeur & « la carte » = \all{die Karte} \\
\hline
\all{das Gymnasium} & le lycée & « le gymnase » = \all{die Turnhalle} \\
\hline
\all{der Onkel wohnt bei uns} & l'oncle habite chez nous & \all{bei} + datif, pas « avec » \\
\hline
\end{tabularx}
\end{center}
\end{attention}

\courssub{Version guidée : trois phrases du texte de Lena}

\begin{exoresolu}[Version guidée n° 1]
Énoncé. Traduis en français : \all{Am Samstag sind alle Gäste um 18 Uhr gekommen.}
\tcblower
Solution. Étape 1 : phrase au passé (auxiliaire \all{sind} + participe \all{gekommen}). Étape 2 : verbe conjugué = \all{sind}, sujet = \all{alle Gäste} (inversion après \all{am Samstag}). Étape 3 : Perfekt $\rightarrow$ passé composé. Traduction : « Samedi, tous les invités sont arrivés à 18 heures. » Justification : \all{kommen} est un verbe de déplacement, d'où l'auxiliaire \all{sein} ; en français, « arriver » se conjugue aussi avec « être ».
\end{exoresolu}

\begin{exoresolu}[Version guidée n° 2]
Énoncé. Traduis : \all{Meine Mutter hat für die ganze Familie gekocht.}
\tcblower
Solution. Verbe conjugué : \all{hat} ; participe final : \all{gekocht} ; sujet : \all{meine Mutter}. \all{für} + accusatif = « pour ». Traduction : « Ma mère a cuisiné pour toute la famille. » Justification : \all{die ganze Familie} = « toute la famille » ; \all{ganz} devant un article défini signifie « entier/tout ».
\end{exoresolu}

\begin{exoresolu}[Version guidée n° 3]
Énoncé. Traduis : \all{Nach dem Essen haben wir zusammen gesungen und getanzt.}
\tcblower
Solution. \all{Nach} + datif = « après » ; deux participes coordonnés : \all{gesungen} et \all{getanzt} partagent le même auxiliaire \all{haben}. Traduction : « Après le repas, nous avons chanté et dansé ensemble. » Justification : un seul auxiliaire pour deux participes, comme en français (« nous avons chanté et dansé »).
\end{exoresolu}

\courssub{Thème guidé : la vie familiale camerounaise}

\begin{exoresolu}[Thème guidé]
Énoncé. Traduis en allemand : « Hier, ma famille a fêté le mariage de mon oncle. »
\tcblower
Solution. Étape 1 : sujet = « ma famille » ; verbe = « a fêté » (passé composé). Étape 2 : Perfekt avec \all{haben} : \all{feiern} $\rightarrow$ \all{gefeiert}. Étape 3 : « hier » en tête $\rightarrow$ inversion : \all{Gestern hat meine Familie ...} Étape 4 : « le mariage de mon oncle » = \all{die Hochzeit meines Onkels} (génitif : \all{meines Onkels}). Résultat : \all{Gestern hat meine Familie die Hochzeit meines Onkels gefeiert.}
\end{exoresolu}

\begin{exercice}[Thème : quatre phrases]
Traduis en allemand en appliquant les 5 étapes :
\begin{enumerate}
\item Le week-end dernier, nous avons rendu visite à notre grand-mère à Douala.
\item Ma sœur aide ma mère à la maison.
\item Demain, mon père invite ses collègues à manger.
\item Après la fête, les enfants sont restés chez leurs cousins.
\end{enumerate}
\end{exercice}

\begin{corrige}[Exercice : thème, quatre phrases]
\begin{enumerate}
\item \all{Letztes Wochenende haben wir unsere Großmutter in Douala besucht.} (inversion après le complément de temps ; \all{besuchen} = « rendre visite à » + accusatif.)
\item \all{Meine Schwester hilft meiner Mutter zu Hause.} (\all{helfen} + datif : \all{meiner Mutter}.)
\item \all{Morgen lädt mein Vater seine Kollegen zum Essen ein.} (futur proche = Präsens + \all{morgen} ; verbe à particule \all{einladen} : \all{ein} en fin de phrase.)
\item \all{Nach dem Fest sind die Kinder bei ihren Cousins geblieben.} (\all{bleiben} $\rightarrow$ auxiliaire \all{sein} ; \all{bei} + datif : \all{bei ihren Cousins}.)
\end{enumerate}
\end{corrige}

\courssub{Mini-texte complet : thème puis version}

\begin{textebox}[Mini-texte de départ (français)]
Dimanche dernier, toute la famille s'est réunie chez mes parents à Yaoundé. Ma tante a préparé un grand repas et mon oncle a apporté des boissons. Après le déjeuner, nous avons écouté de la musique et les enfants ont joué dans le jardin. Le soir, tout le monde était fatigué mais très content.
\end{textebox}

\begin{exoresolu}[Thème : le mini-texte]
Énoncé. Traduis le mini-texte ci-dessus en allemand.
\tcblower
Solution détaillée, phrase par phrase :
\begin{enumerate}
\item « Dimanche dernier » en tête $\rightarrow$ inversion ; « se réunir » = \all{sich treffen} (Perfekt : \all{hat sich getroffen}) ; « chez mes parents » = \all{bei meinen Eltern}. $\rightarrow$ \all{Letzten Sonntag hat sich die ganze Familie bei meinen Eltern in Jaunde getroffen.}
\item Deux actions au Perfekt : \all{zubereiten} (préparer) et \all{mitbringen} (apporter). $\rightarrow$ \all{Meine Tante hat ein großes Essen zubereitet und mein Onkel hat Getränke mitgebracht.}
\item \all{Musik hören}, \all{im Garten spielen}. $\rightarrow$ \all{Nach dem Mittagessen haben wir Musik gehört und die Kinder haben im Garten gespielt.}
\item « était » $\rightarrow$ Präteritum de \all{sein} : \all{war}. $\rightarrow$ \all{Am Abend waren alle müde, aber sehr zufrieden.}
\end{enumerate}
\end{exoresolu}

\begin{exoresolu}[Version : le même texte en sens inverse]
Énoncé. Traduis en français : \all{Letztes Wochenende haben wir unsere Verwandten in Bafoussam besucht. Mein Bruder hat seiner Tante beim Kochen geholfen. Nach dem Essen sind wir lange im Garten geblieben und haben viel gelacht. Am Abend sind wir müde nach Hause gefahren.}
\tcblower
Solution. « Le week-end dernier, nous avons rendu visite à nos parents/parentés à Bafoussam. » — attention, \all{die Verwandten} = « les membres de la famille, la parenté ». « Mon frère a aidé sa tante à cuisiner. » (\all{helfen} + datif : \all{seiner Tante} ; \all{beim Kochen} = « à cuisiner ».) « Après le repas, nous sommes restés longtemps dans le jardin et nous avons beaucoup ri. » « Le soir, nous sommes rentrés à la maison, fatigués. » (\all{nach Hause fahren} = « rentrer à la maison » ; auxiliaire \all{sein} car déplacement.)
\end{exoresolu}

\begin{methode}[Grille d'auto-correction après une traduction]
Relis ta production en cochant mentalement chaque point :
\begin{enumerate}
\item \textbf{Sens} : ai-je traduit toutes les informations (temps, lieux, personnes) ?
\item \textbf{Temps} : passé composé français $\rightarrow$ Perfekt ; le participe est-il en fin de phrase ?
\item \textbf{Ordre des mots} : le verbe conjugué est-il bien en 2\textsuperscript{e} position ? Inversion après un complément initial ?
\item \textbf{Cas} : accusatif ou datif correct après les verbes (\all{helfen} + datif) et les prépositions (\all{bei}, \all{nach}, \all{mit} + datif) ?
\item \textbf{Orthographe} : majuscule à tous les noms, Umlaute (ä, ö, ü), \all{ß}, pronoms réfléchis (\all{sich}, \all{uns}) présents ?
\end{enumerate}
\end{methode}

\begin{aretenir}[Barème-type au baccalauréat]
La traduction est généralement notée sur trois critères :
\begin{itemize}
\item \textbf{Sens global} (environ 40 \%) : toutes les idées du texte sont rendues, sans omission ni contresens.
\item \textbf{Exactitude grammaticale} (environ 40 \%) : temps corrects, verbe en position 2, participe final, cas et auxiliaires justes.
\item \textbf{Orthographe et présentation} (environ 20 \%) : majuscules des noms, Umlaute, ponctuation, français ou allemand naturel.
\end{itemize}
Une phrase grammaticalement parfaite mais au sens faux ne rapporte presque rien ; une phrase au sens juste avec une petite faute de cas perd peu. Priorité au sens, ensuite à la forme.
\end{aretenir}

\begin{savaistu}[Traduire, un métier et une tradition]
Entre le Cameroun et les pays germanophones, la traduction est partout : les lycéens camerounais qui partent étudier en Allemagne grâce aux bourses du DAAD doivent souvent faire traduire officiellement leurs relevés de notes par un \all{vereidigter Übersetzer} (traducteur assermenté). Et savais-tu que la Bible de Luther (1534) est considérée comme l'une des premières grandes traductions qui ont façonné la langue allemande moderne ? Traduire, c'est construire des ponts entre les cultures — exactement ce que tu fais dans cette leçon.
\end{savaistu}

\begin{exercice}[Pour aller plus loin]
Traduis ce court dialogue en allemand, puis vérifie avec la grille d'auto-correction :
« — Tu es allé à la fête de famille samedi ? — Oui, nous avons beaucoup dansé et ma grand-mère a raconté des histoires. — Demain, je rends visite à mon oncle à Bafoussam. Tu viens avec moi ? »
\end{exercice}

\begin{corrige}[Exercice : pour aller plus loin]
\all{— Bist du am Samstag zum Familienfest gegangen? — Ja, wir haben viel getanzt und meine Großmutter hat Geschichten erzählt. — Morgen besuche ich meinen Onkel in Bafoussam. Kommst du mit?}
Points de vigilance : \all{gehen} $\rightarrow$ auxiliaire \all{sein} (\all{bist ... gegangen}) ; \all{erzählen} = « raconter » ; futur proche rendu par le Präsens + \all{morgen} ; \all{mitkommen} = « venir avec ».
\end{corrige}

\newpage

\coursec{Activité d'intégration}

\courssub{Objectifs de l'activité}

À la fin de cette activité d'intégration, tu seras capable de :

\begin{itemize}
\item rédiger une courte lettre personnelle en allemand (8 à 10 phrases) sur le thème de la vie familiale, en respectant la structure d'une lettre (lieu et date, salutation, corps, formule de politesse, signature) ;
\item réinvestir correctement le lexique de la famille et de la vie sociale (der Vater, die Schwester, das Fest, feiern, gratulieren\ldots) ;
\item employer à bon escient le présent (se présenter, décrire sa famille) et le parfait (raconter un événement passé : une fête familiale) ;
\item respecter l'ordre des mots allemand : verbe conjugué en deuxième position, participe passé en fin de phrase ;
\item traduire une lettre de l'allemand vers le français (version) et vérifier une production avec une check-list de relecture ;
\item lire un texte à voix haute avec une prononciation claire.
\end{itemize}

\courssub{Situation}

Ton lycée de Yaoundé participe à un échange scolaire avec un Gymnasium de Munich, en Bavière. Chaque élève de ta classe de Première a reçu le nom d'un correspondant ou d'une correspondante allemand(e). Le professeur d'allemand de Munich demande à ses élèves d'écrire d'abord ; mais cette semaine, c'est à ton tour de prendre la plume ! Tu as reçu la lettre suivante de ta correspondante, Lena, et tu dois maintenant lui répondre.

\begin{textebox}[La lettre de Lena — le document de départ]
München, den 12. März\par
\medskip
Lieber Freund, liebe Freundin,\par
\medskip
ich heiße Lena und ich bin sechzehn Jahre alt. Ich wohne mit meinen Eltern und meinem Bruder in München. Mein Vater ist Ingenieur und meine Mutter ist Lehrerin. Mein Bruder Tom ist erst zehn Jahre alt.\par
\medskip
Letztes Wochenende haben wir den Geburtstag meiner Großmutter gefeiert. Sie ist achtzig Jahre alt geworden! Die ganze Familie ist gekommen: meine Tanten, meine Onkel und meine Cousins. Wir haben einen großen Kuchen gebacken und wir haben zusammen gesungen und gelacht. Es war sehr schön.\par
\medskip
Und du? Wie ist deine Familie? Hast du ein Familienfest schon gefeiert? Schreib mir bald!\par
\medskip
Viele Grüße\par
Deine Lena
\end{textebox}

\textbf{Glossaire :} der Ingenieur, -e (l'ingénieur) ; die Großmutter, -mütter (la grand-mère) ; feiern (fêter) ; der Kuchen, - (le gâteau) ; backen (cuire, faire un gâteau) ; letztes Wochenende (le week-end dernier).

\courssub{Consignes}

Travaillez par binômes. Suivez les étapes dans l'ordre.

\begin{enumerate}
\item \textbf{Compréhension.} Lisez la lettre de Lena à voix haute, chacun votre tour. Puis répondez par écrit, en français, à ces trois questions : Qui est Lena ? Avec qui habite-t-elle ? Quelle fête la famille a-t-elle célébrée et qui était présent ?
\item \textbf{Repérage grammatical.} Dans la lettre de Lena, soulignez en bleu trois verbes au présent et en rouge trois verbes au parfait. Pour chaque verbe au parfait, entourez l'auxiliaire (\all{haben} ou \all{sein}) et le participe passé. Que remarquez-vous sur la place du participe passé dans la phrase ?
\item \textbf{Préparation du lexique.} Complétez individuellement une petite fiche : les membres de votre famille (au moins cinq mots avec leur article), trois verbes utiles pour raconter une fête (par exemple \all{feiern}, \all{einladen}, \all{essen}, \all{tanzen}, \all{schenken}) et deux adjectifs pour décrire l'ambiance (par exemple \all{schön}, \all{lustig}, \all{groß}).
\item \textbf{Production écrite.} Rédigez votre réponse à Lena : une lettre de 8 à 10 phrases. Votre lettre doit contenir obligatoirement : le lieu et la date (par exemple \all{Yaoundé, den 20. März}) ; une salutation (\all{Liebe Lena,}) ; une présentation de vous et de votre famille au présent (4 phrases) ; le récit d'une fête familiale récente au parfait (3 à 4 phrases : un mariage, un anniversaire, une fête de fin d'année, un baptême\ldots) ; deux questions à Lena ; une formule finale (\all{Viele Grüße}) et votre prénom.
\item \textbf{Échange et version.} Échangez votre lettre avec celle de votre camarade. Traduisez la lettre reçue en français, phrase par phrase, sur votre cahier. Attention aux pièges : \all{meine Eltern} signifie « mes parents », \all{das Fest} signifie « la fête » et non « le feste » !
\item \textbf{Relecture croisée.} Vérifiez la lettre de votre camarade avec la check-list ci-dessous. Notez les erreurs au crayon, sans corriger à sa place : c'est l'auteur qui corrige.
\item \textbf{Correction finale et lecture.} Corrigez votre lettre, recopiez-la proprement, puis les binômes volontaires lisent leur lettre devant la classe. Les meilleures lettres seront affichées au mur de la classe et envoyées (en photo) au Gymnasium de Munich.
\end{enumerate}

\begin{methode}[La check-list de relecture]
Avant de rendre ta lettre, vérifie dans l'ordre :
\begin{enumerate}
\item \textbf{Majuscules} : tous les noms commencent par une majuscule (\all{der Bruder}, \all{das Fest}, \all{München}).
\item \textbf{Position du verbe} : le verbe conjugué est en deuxième position, même après un complément de temps : \all{Letztes Wochenende \textbf{haben} wir gefeiert.}
\item \textbf{Parfait} : auxiliaire correct (\all{sein} avec les verbes de mouvement ou de changement : \all{ist gekommen} ; \all{haben} ailleurs : \all{haben gefeiert}) et participe passé en fin de phrase.
\item \textbf{Structure de la lettre} : lieu et date, salutation avec virgule et minuscule à la ligne suivante, formule finale, signature.
\item \textbf{Orthographe} : Umlaute (ä, ö, ü), \all{ß} ou \all{ss}, vocabulaire vérifié dans le glossaire.
\end{enumerate}
\end{methode}

\courssub{Ressources à mobiliser}

\begin{aretenir}[Le présent et le parfait : les bons réflexes]
\textbf{Présent} : pour se présenter et décrire sa famille : \all{Ich heiße Alima. Ich wohne in Yaoundé. Mein Vater arbeitet in einer Bank.} « Je m'appelle Alima. J'habite à Yaoundé. Mon père travaille dans une banque. »\par
\medskip
\textbf{Parfait} : pour raconter la fête : auxiliaire (\all{haben} ou \all{sein}) conjugué en deuxième position + participe passé à la fin : \all{Wir haben ein Fest gefeiert.} « Nous avons fêté une fête. » ; \all{Meine Tante ist aus Douala gekommen.} « Ma tante est venue de Douala. »
\end{aretenir}

\begin{center}
\begin{tabularx}{\linewidth}{|X|X|X|}\hline
\rowcolor{popblueL} Français & Deutsch & Remarque \\ \hline
mon père / ma mère & der Vater / die Mutter & \all{mein Vater}, \all{meine Mutter} \\ \hline
les parents & die Eltern & toujours au pluriel \\ \hline
le frère / la sœur & der Bruder / die Schwester & pluriel : Brüder / Schwestern \\ \hline
fêter & feiern & \all{wir haben gefeiert} \\ \hline
inviter & einladen & séparable : \all{wir haben eingeladen} \\ \hline
venir & kommen & \all{ist gekommen} (auxiliaire \all{sein}) \\ \hline
offrir & schenken & \all{ich habe ein Geschenk geschenkt} \\ \hline
c'était très beau & es war sehr schön & Präteritum de \all{sein}, très courant \\ \hline
\end{tabularx}
\end{center}

\begin{attention}[Erreurs fréquentes des francophones]
\begin{itemize}
\item Ne dis jamais \all{Ich habe gekommen} : \all{kommen} exige l'auxiliaire \all{sein} : \all{Ich bin gekommen.} « Je suis venu(e). »
\item Après la salutation \all{Liebe Lena,} la phrase suivante commence par une \textbf{minuscule} en allemand (sauf nom ou \all{ich} en début absolu).
\item Le participe passé se place à la fin : pas \all{Wir haben gefeiert ein Fest}, mais \all{Wir haben ein Fest gefeiert.}
\item \all{bekommen} ne signifie pas « devenir » mais « recevoir » ; « devenir » se dit \all{werden} : \all{Sie ist achtzig geworden.} « Elle a eu quatre-vingts ans. »
\end{itemize}
\end{attention}

\courssub{Proposition de production et corrigé commenté}

\begin{exoresolu}[Une lettre modèle : la réponse d'Alima]
Rédige une lettre de 8 à 10 phrases en réponse à Lena : présente-toi, décris ta famille et raconte une fête familiale récente au parfait.
\tcblower
\textbf{Production modèle :}\par
\medskip
Yaoundé, den 20. März\par
\medskip
Liebe Lena,\par
\medskip
vielen Dank für deinen Brief. Ich heiße Alima und ich bin sechzehn Jahre alt. Ich wohne mit meiner Familie in Yaoundé, der Hauptstadt von Kamerun. Mein Vater ist Lehrer und meine Mutter verkauft Stoffe auf dem Markt. Ich habe zwei Brüder und eine kleine Schwester.\par
\medskip
Letzten Monat haben wir den Geburtstag meines Großvaters gefeiert. Die ganze Familie ist zu uns gekommen, auch meine Tante aus Douala. Meine Mutter hat einen guten Fisch gekocht, und wir haben zusammen gegessen, getanzt und gelacht. Ich habe meinem Großvater ein Geschenk gegeben. Es war ein sehr schönes Fest!\par
\medskip
Und du, Lena? Hast du Geschwister? Was habt ihr am Wochenende gemacht? Schreib mir bald!\par
\medskip
Viele Grüße\par
Deine Alima\par
\medskip
\textbf{Commentaire des choix de langue :}\par
\begin{itemize}
\item \textbf{Structure} : lieu et date en haut (\all{Yaoundé, den 20. März}), salutation suivie d'une virgule et d'une minuscule (\all{vielen Dank\ldots}), formule finale \all{Viele Grüße} et signature avec \all{Deine Alima}.
\item \textbf{Présent} pour la présentation : \all{ich heiße}, \all{ich wohne}, \all{mein Vater ist}, \all{ich habe}. Le verbe est toujours en deuxième position.
\item \textbf{Parfait} pour le récit de la fête : \all{haben wir gefeiert}, \all{hat gekocht}, \all{haben gegessen, getanzt und gelacht}, \all{habe gegeben} ; avec l'auxiliaire \all{sein} pour le verbe de mouvement : \all{ist gekommen}.
\item \textbf{Ordre des mots} : après le complément de temps \all{Letzten Monat}, le verbe reste en deuxième position et le sujet suit : \all{Letzten Monat \textbf{haben} \textbf{wir} gefeiert.} Le participe passé est bien en fin de phrase.
\item \textbf{Lexique du chapitre} : \all{die Hauptstadt} (la capitale), \all{die Stoffe} (les tissus), \all{der Markt} (le marché), \all{die Geschwister} (les frères et sœurs), \all{das Geschenk} (le cadeau).
\item \textbf{Interaction} : la lettre se termine par deux questions au correspondant, comme dans la lettre de Lena : cela donne un échange authentique.
\end{itemize}
\end{exoresolu}

\begin{exercice}[Version : traduis la lettre modèle]
Traduis en français le deuxième paragraphe de la lettre d'Alima (de \all{Letzten Monat} à \all{sehr schönes Fest}). Fais attention aux temps : le parfait allemand se rend souvent par un passé composé français.
\end{exercice}

\begin{corrige}[Version : traduis la lettre modèle]
« Le mois dernier, nous avons fêté l'anniversaire de mon grand-père. Toute la famille est venue chez nous, même ma tante de Douala. Ma mère a cuisiné un bon poisson, et nous avons mangé, dansé et ri ensemble. J'ai offert un cadeau à mon grand-père. C'était une très belle fête ! »\par
Remarques : \all{zu uns} se traduit par « chez nous » ; \all{auch} signifie ici « même / aussi » ; \all{es war} (Präteritum) se traduit par l'imparfait « c'était ».
\end{corrige}

\courssub{Bilan de l'activité}

Dans cette activité d'intégration, tu as mobilisé l'ensemble des savoir-faire de la leçon : comprendre une lettre authentique, repérer les temps de verbe, rédiger une production personnelle structurée, traduire de l'allemand vers le français et relire un texte avec une méthode précise. Tu as constaté que la lettre personnelle allemande obéit à des règles fixes (date, salutation, minuscule après la virgule, formule finale) et que le choix entre présent et parfait dépend du sens : le présent décrit une situation stable, le parfait raconte un événement terminé. L'échange par binômes t'a enfin appris à regarder un texte avec l'œil du correcteur : majuscules des noms, verbe en deuxième position, auxiliaire \all{haben} ou \all{sein}, participe en fin de phrase. Ces réflexes de relecture te serviront pour toutes tes productions futures, en classe comme à l'examen.

\newpage

\coursec{Méthodes et exercices résolus}

\courssub{Méthode 1 : Rédiger une courte lettre sur la vie familiale}

\begin{methode}[Rédiger une lettre personnelle en allemand]
Pour rédiger une lettre simple (invitation, remerciement, nouvelles de la famille), suis toujours les mêmes étapes :
\begin{enumerate}
\item \textbf{Repérer la consigne} : à qui écris-tu ? Quel sujet ? Combien de mots ? Souligne les mots-clés de l'énoncé.
\item \textbf{Poser le cadre} : lieu et date en haut à droite (\all{Yaoundé, den 15. März 2025}), salutation (\all{Liebe Anna,} / \all{Lieber Paul,}).
\item \textbf{Construire le corps en 3 temps} : ouverture (\all{vielen Dank für deinen Brief} / \all{wie geht es dir?}), développement (2 ou 3 idées au présent ou au Perfekt), conclusion (\all{ich freue mich auf deine Antwort}).
\item \textbf{Surveiller la grammaire} : verbe conjugué en 2\textsuperscript{e} position, participe passé en fin de phrase au Perfekt, majuscule à tous les noms.
\item \textbf{Terminer} par la formule de politesse (\all{Viele Grüße} / \all{Deine} + prénom), puis \textbf{relire} : orthographe, cas, ordre des mots.
\end{enumerate}
Structures à retenir : \all{Ich schreibe dir, weil...} (weil + verbe en fin), \all{Letztes Wochenende haben wir... gefeiert.} (Perfekt), \all{Kommst du mit?}
\end{methode}

\begin{exoresolu}[Lettre : une fête de famille]
Énoncé : Tu as participé à une fête de famille le week-end dernier. Écris une lettre de 60 à 80 mots à ton ami allemand Paul : dis où et quand la fête a eu lieu, qui était là et ce que vous avez fait.
\tcblower
\textbf{Étape 1 — Analyse de la consigne} : destinataire = Paul (tutoiement) ; temps principal = Perfekt (événement passé) ; 3 informations obligatoires : lieu/date, personnes, activités.

\textbf{Étape 2 — Plan} : salutation + ouverture ; phrase 1 : \all{Letzten Samstag haben wir ein Familienfest gefeiert.} ; phrase 2 : personnes présentes ; phrase 3 : activités ; conclusion + formule finale.

\textbf{Étape 3 — Rédaction} :

\all{Yaoundé, den 12. Mai 2025}

\all{Lieber Paul,}

\all{wie geht es dir? Letzten Samstag haben wir bei meiner Großmutter ein großes Familienfest gefeiert. Meine Eltern, meine Geschwister, meine Onkel und meine Cousins waren alle da. Wir haben zusammen gegessen, getanzt und viel gelacht. Meine Großmutter hat ein leckeres Essen gekocht. Es war ein wunderbarer Tag!}

\all{Schreib mir bald!}

\all{Viele Grüße}

\all{Junior}

\textbf{Étape 4 — Justifications grammaticales} : Perfekt avec \all{haben} + participes en fin de phrase (\all{gefeiert, gegessen, getanzt, gelacht, gekocht}) ; \all{letzten Samstag} : accusatif de temps, adjectif en \emph{-en} ; \all{bei meiner Großmutter} : datif après \all{bei} ; verbe toujours en 2\textsuperscript{e} position, y compris après le complément de temps en tête.

\textbf{Conclusion} : la lettre fait 75 mots, respecte les 3 points demandés et utilise le Perfekt correctement.
\end{exoresolu}

\courssub{Méthode 2 : Décrire une personne ou une scène familiale}

\begin{methode}[Rédiger une description]
\begin{enumerate}
\item \textbf{Choisir l'angle} : description physique (taille, cheveux, yeux) puis caractère (gentil, drôle, sérieux).
\item \textbf{Utiliser le présent} et les verbes de description : \all{sein, haben, aussehen, wohnen, arbeiten}.
\item \textbf{Accorder les adjectifs} : après \all{sein}, pas de terminaison (\all{Mein Bruder ist groß.}) ; devant un nom, terminaison obligatoire (\all{mein großer Bruder}).
\item \textbf{Varier le vocabulaire} : famille (\all{der Vater, die Mutter, die Geschwister}), qualités (\all{freundlich, hilfsbereit, lustig, streng}).
\item \textbf{Relire} : majuscules des noms, position du verbe, accords.
\end{enumerate}
\end{methode}

\begin{exoresolu}[Description : ma famille]
Énoncé : Décris ta famille en 5 à 6 phrases (présent) : les membres, l'apparence et le caractère d'une personne que tu admires.
\tcblower
\textbf{Étape 1} : je choisis de décrire ma mère ; temps = présent ; verbes : \all{sein, haben, arbeiten, helfen}.

\textbf{Étape 2 — Rédaction} :

\all{Meine Familie ist nicht groß: mein Vater, meine Mutter, meine Schwester und ich. Meine Mutter ist 45 Jahre alt. Sie ist klein und hat schwarze Haare und braune Augen. Sie arbeitet als Lehrerin in einer Schule in Douala. Meine Mutter ist sehr freundlich und hilfsbereit: sie hilft immer den Nachbarn. Ich bewundere sie, weil sie für die Familie alles tut.}

\textbf{Étape 3 — Justifications} : \all{ist nicht groß} : adjectif sans terminaison après \all{sein} ; \all{hat schwarze Haare} : accusatif pluriel, adjectif en \emph{-e} ; \all{arbeitet als Lehrerin} : présent, 3\textsuperscript{e} personne en \emph{-t} ; \all{den Nachbarn} : datif pluriel après \all{helfen} (verbe à datif !) avec \emph{-n} au pluriel datif ; \all{weil sie ... tut} : subordonnée avec \all{weil}, verbe conjugué en dernière position.

\textbf{Conclusion} : description complète (membres, physique, caractère), présent correct, un verbe à datif bien employé.
\end{exoresolu}

\courssub{Méthode 3 : Rédiger un dialogue}

\begin{methode}[Construire un dialogue naturel]
\begin{enumerate}
\item \textbf{Identifier la situation} : qui parle ? Où ? Quel est le but (inviter, demander, proposer) ?
\item \textbf{Alterner les répliques courtes} : une phrase par réplique suffit au niveau A2 ; nom du personnage suivi de deux-points.
\item \textbf{Utiliser les formules de communication} : question (\all{Kommst du mit? Was machst du am Samstag?}), réponse (\all{Ja, gern! / Leider nicht.}), proposition (\all{Wollen wir...?}).
\item \textbf{Placer le verbe correctement} : position 2 dans les déclaratives, position 1 dans les questions sans mot interrogatif (\all{Hast du Zeit?}).
\item \textbf{Relire} : chaque réplique doit avoir un verbe conjugué et une majuscule en début de phrase.
\end{enumerate}
\end{methode}

\begin{exoresolu}[Dialogue : une invitation]
Énoncé : Écris un dialogue de 8 répliques entre toi et ton amie Clara : tu l'invites à l'anniversaire de ta sœur ; elle pose des questions (quand, où) et accepte.
\tcblower
\textbf{Étape 1} : situation = invitation ; formules utiles : \all{Kommst du zu...? Wann? Wo? Ich komme gern.}

\textbf{Étape 2 — Rédaction} :

\all{Ich: Hallo Clara! Meine Schwester hat am Samstag Geburtstag. Kommst du zu ihrer Party?}

\all{Clara: Gern! Wann beginnt die Party?}

\all{Ich: Sie beginnt um 16 Uhr.}

\all{Clara: Und wo feiert ihr?}

\all{Ich: Wir feiern bei uns zu Hause, in Yaoundé.}

\all{Clara: Was kann ich mitbringen?}

\all{Ich: Bring bitte etwas zu trinken mit.}

\all{Clara: Kein Problem! Ich freue mich schon. Bis Samstag!}

\textbf{Étape 3 — Justifications} : \all{Kommst du...?} : question fermée, verbe en position 1 ; \all{Wann beginnt die Party?} : question avec mot interrogatif, verbe en position 2 ; \all{um 16 Uhr} : heure avec \all{um} ; \all{bei uns zu Hause} : datif après \all{bei} ; \all{Bring ... mit} : impératif du verbe à particule \all{mitbringen}, particule \all{mit} renvoyée en fin de phrase.

\textbf{Conclusion} : 8 répliques équilibrées, invitation + 2 questions + acceptation, ordre des mots respecté.
\end{exoresolu}

\courssub{Méthode 4 : Traduire du français vers l'allemand (thème)}

\begin{methode}[Réussir le thème]
\begin{enumerate}
\item \textbf{Lire toute la phrase} et repérer le temps (présent ou passé → Perfekt), le sujet et le verbe.
\item \textbf{Traduire le verbe d'abord} et le placer en position 2 ; au Perfekt, l'auxiliaire en position 2 et le participe en fin.
\item \textbf{Traiter les groupes nominaux} : article + cas (qui ? nominatif ; quoi ? accusatif ; à qui ? datif) + majuscule au nom.
\item \textbf{Attention aux connecteurs} : « parce que » = \all{weil} (verbe en fin), « mais » = \all{aber} (verbe en position 2).
\item \textbf{Relire mot à mot} : chaque mot français doit avoir un correspondant ; vérifier Umlaute, ß, majuscules.
\end{enumerate}
\end{methode}

\begin{exoresolu}[Thème : phrases sur la famille]
Énoncé : Traduis en allemand : 1. Ma grand-mère habite avec nous à Yaoundé. 2. Hier, nous avons rendu visite à nos voisins. 3. Mon frère ne vient pas à la fête parce qu'il est malade.
\tcblower
\textbf{Phrase 1} : temps = présent ; verbe \all{wohnen} ; « avec nous » = \all{bei uns} (datif). \all{Meine Großmutter wohnt bei uns in Yaoundé.} — Verbe en position 2, majuscules à \all{Großmutter} et \all{Yaoundé}.

\textbf{Phrase 2} : « hier » + passé → Perfekt ; « rendre visite à » = \all{besuchen} (accusatif) ; « nos voisins » = \all{unsere Nachbarn}. \all{Gestern haben wir unsere Nachbarn besucht.} — Auxiliaire \all{haben} en position 2 (le complément de temps \all{gestern} occupe la position 1), participe \all{besucht} en fin de phrase.

\textbf{Phrase 3} : « ne vient pas » = \all{kommt ... nicht} ; « parce qu'il est malade » = subordonnée \all{weil}, verbe en fin. \all{Mein Bruder kommt nicht zum Fest, weil er krank ist.} — \all{zum} = \all{zu dem} (datif neutre) ; dans la subordonnée, \all{ist} va en dernière position.

\textbf{Conclusion} : trois pièges classiques évités : place du participe, datif après \all{bei}/\all{zu}, verbe final avec \all{weil}.
\end{exoresolu}

\courssub{Méthode 5 : Traduire de l'allemand vers le français (version)}

\begin{methode}[Réussir la version]
\begin{enumerate}
\item \textbf{Identifier le verbe conjugué} (position 2) et le participe éventuel en fin de phrase : cela donne le temps.
\item \textbf{Repérer sujet et compléments} grâce aux articles (\all{der/die/das} = sujet ; \all{den/die/das} = COD ; \all{dem/der/dem} = COI).
\item \textbf{Traduire d'abord le sens global}, en français naturel, sans calquer l'ordre allemand.
\item \textbf{Attention aux verbes à particule} : \all{Er lädt seine Freunde ein.} = « Il invite ses amis » (\all{einladen}), pas « il charge ».
\item \textbf{Relire en français} : la phrase doit être correcte et fluide en français, avec tous les éléments du texte allemand.
\end{enumerate}
\end{methode}

\begin{exoresolu}[Version : un court texte]
Énoncé : Traduis en français : \all{Am Sonntag hat Familie Mbarga ihre Verwandten aus Bafoussam eingeladen. Die Kinder haben im Garten gespielt, und die Erwachsenen haben zusammen gekocht. Alle waren sehr zufrieden.}
\tcblower
\textbf{Étape 1 — Analyse} : \all{hat ... eingeladen} = Perfekt de \all{einladen} (inviter) ; \all{ihre Verwandten} = accusatif pluriel (COD) ; \all{haben gespielt / haben gekocht} = Perfekt ; \all{die Erwachsenen} = « les adultes » (adjectif nominalisé, majuscule !) ; \all{waren} = Präteritum de \all{sein}.

\textbf{Étape 2 — Traduction} : « Dimanche, la famille Mbarga a invité ses parents (ses proches) de Bafoussam. Les enfants ont joué dans le jardin, et les adultes ont cuisiné ensemble. Tout le monde était très content. »

\textbf{Étape 3 — Vérifications} : \all{Verwandte} se traduit « parents/proches » au sens de « membres de la famille », jamais « parents » au sens de père et mère (faux ami !) ; \all{zufrieden} = « content, satisfait » ; l'ordre français est naturel, sans calque.

\textbf{Conclusion} : version fidèle (aucun élément oublié) et française correcte.
\end{exoresolu}

\begin{attention}[Les 5 erreurs qui coûtent des points au Bac]
\begin{enumerate}
\item \textbf{Oublier la majuscule aux noms} : \all{meine mutter} est faux → \all{meine Mutter}. En allemand, TOUT nom prend une majuscule (\all{das Fest, die Kinder, die Erwachsenen}).
\item \textbf{Mal placer le verbe} : le verbe conjugué est TOUJOURS en 2\textsuperscript{e} position : \all{Gestern haben wir gefeiert.} et jamais \emph{Gestern wir haben gefeiert}. Avec \all{weil}, le verbe va en fin : \all{weil er krank ist}.
\item \textbf{Se tromper d'auxiliaire au Perfekt} : les verbes de déplacement ou de changement prennent \all{sein} : \all{Wir sind nach Douala gefahren.} (et non \emph{haben gefahren}).
\item \textbf{Tomber dans les faux amis} : \all{die Verwandten} = les proches (pas « les parents ») ; \all{bekommen} = recevoir (pas « devenir » = \all{werden}) ; \all{aktuelle} = actuel (pas « actuel » au sens de vrai).
\item \textbf{Calquer le français} : « j'ai 15 ans » = \all{Ich bin 15 Jahre alt.} (et non \emph{ich habe 15 Jahre}) ; « rendre visite à mes grands-parents » = \all{ich besuche meine Großeltern} (verbe simple, pas de préposition calquée).
\end{enumerate}
Réflexe final : relis toujours ta copie en ne regardant QUE le verbe de chaque phrase (place, conjugaison, auxiliaire), puis QUE les majuscules des noms.
\end{attention}

\newpage

\coursec{Exercices}

\courssub{Je vérifie mes connaissances}

\begin{exercice}[QCM : la bonne forme]
Entoure la bonne réponse.
\begin{enumerate}
\item \all{Gestern ... wir unsere Großeltern besucht.} \quad a) \all{haben} \quad b) \all{sind} \quad c) \all{hatten}
\item \all{Meine Schwester ... gestern zu Hause geblieben.} \quad a) \all{hat} \quad b) \all{ist} \quad c) \all{wird}
\item \all{Am Samstag ... wir ein großes Fest.} \quad a) \all{feiern} \quad b) \all{feiert} \quad c) \all{feiere}
\item \all{Ich gratuliere ... zum Geburtstag.} \quad a) \all{du} \quad b) \all{dich} \quad c) \all{dir}
\item \all{Meine Tante ... uns gestern eingeladen.} \quad a) \all{hat} \quad b) \all{ist} \quad c) \all{hast}
\end{enumerate}
\end{exercice}

\begin{exercice}[Vrai ou faux ? Justifie ta réponse]
Dis si chaque phrase est vraie (V) ou fausse (F), puis corrige les phrases fausses.
\begin{enumerate}
\item Au parfait, le verbe \all{kommen} se conjugue avec l'auxiliaire \all{haben}.
\item Le participe passé de \all{essen} est \all{gegessen}.
\item Dans une phrase au parfait, le participe passé se place à la fin de la phrase.
\item Le verbe \all{helfen} se construit avec l'accusatif.
\item En allemand, tous les noms communs prennent une majuscule.
\end{enumerate}
\end{exercice}

\begin{exercice}[Vocabulaire : la famille et les fêtes]
Complète chaque phrase avec le mot allemand qui convient : \all{die Hochzeit}, \all{der Geburtstag}, \all{die Einladung}, \all{das Geschenk}, \all{die Verwandten}.
\begin{enumerate}
\item Meine Großmutter wird 80 Jahre alt : wir feiern ihren ...
\item Paul und Lena heiraten im Juni : die ... findet in der Kirche statt.
\item Ich schreibe eine ... an meinen Freund in Douala.
\item Alle ... kommen aus Bamenda und Bafoussam zum Fest.
\item Für meine Mutter kaufe ich ein schönes ...
\end{enumerate}
\end{exercice}

\begin{exercice}[Conjugaison : présent ou parfait ?]
Complète le texte avec les verbes entre parenthèses au temps demandé.
\begin{enumerate}
\item Meine Familie \textbf{(wohnen, Präsens)} in Yaoundé.
\item Letztes Wochenende \textbf{(besuchen, Perfekt)} wir unsere Tante.
\item Mein Bruder \textbf{(helfen, Präsens)} meiner Mutter in der Küche.
\item Wir \textbf{(essen, Perfekt)} gestern einen großen Kuchen.
\item Meine Eltern \textbf{(sein, Perfekt)} sehr zufrieden.
\end{enumerate}
\end{exercice}

\courssub{Je m'entraîne}

\begin{exercice}[L'ordre des mots : remets dans l'ordre]
Reconstitue des phrases correctes. Attention à la place du verbe et du participe.
\begin{enumerate}
\item \all{eingeladen / hat / uns / meine Tante / gestern}
\item \all{meine Großeltern / am Sonntag / wir / besucht / haben}
\item \all{ein Geschenk / ich / meiner Schwester / kaufe / heute}
\item \all{die Familie / ist / nach München / gefahren / letzte Woche}
\item \all{steht / mein Bruder / jeden Tag / früh / auf}
\end{enumerate}
\end{exercice}

\begin{exercice}[Texte à trous : le parfait]
Complète le texte avec l'auxiliaire (\all{haben} ou \all{sein}) et le participe passé des verbes entre parenthèses.

\begin{textebox}[Das Fest]
Gestern ... wir den Geburtstag meines Vaters \textbf{(feiern)}. Viele Gäste ... zu uns \textbf{(kommen)}. Meine Mutter ... ein gutes Essen \textbf{(kochen)}. Mein Bruder ... lange in der Küche \textbf{(bleiben)} und er ... meiner Mutter \textbf{(helfen)}. Nach dem Essen ... wir alle zusammen \textbf{(singen)}. Mein Vater ... sich sehr \textbf{(freuen)}.
\end{textebox}
\end{exercice}

\begin{exercice}[Appariements : invitations et réponses]
Relie chaque phrase d'invitation (1--4) à la réponse qui convient (a--d).
\begin{enumerate}
\item \all{Möchtest du zu meiner Geburtstagsfeier kommen?}
\item \all{Kannst du am Samstag zu unserer Hochzeit kommen?}
\item \all{Hast du Lust, mit uns ins Restaurant zu gehen?}
\item \all{Kommst du morgen zu unserem Familienfest?}
\end{enumerate}
\begin{itemize}
\item[a)] \all{Ja, gern! Wann fängt die Feier an?}
\item[b)] \all{Leider kann ich nicht, ich muss meiner Mutter helfen.}
\item[c)] \all{Natürlich! Ich freue mich schon sehr darauf.}
\item[d)] \all{Gute Idee! Um wie viel Uhr treffen wir uns?}
\end{itemize}
\end{exercice}

\begin{exercice}[Thème : phrases à traduire]
Traduis les phrases suivantes en allemand.
\begin{enumerate}
\item Ma grand-mère habite à Douala avec mon oncle.
\item Hier, nous avons visité nos grands-parents.
\item Ma sœur a aidé ma mère dans la cuisine.
\item Dimanche dernier, toute la famille est venue chez nous.
\item J'ai offert un beau cadeau à mon frère pour son anniversaire.
\end{enumerate}
\end{exercice}

\courssub{Je me prépare au Bac}

\begin{exercice}[Compréhension et étude de texte (type Bac)]
Lis le texte suivant, puis réponds aux questions \textbf{en allemand et en phrases complètes}.

\begin{textebox}[Ein Familienfest in München]
Letzten Samstag hat die Familie Berger in München ein großes Fest gefeiert. Die Großmutter ist 75 Jahre alt geworden. Alle Verwandten sind gekommen : die Kinder, die Enkel und auch die Tante aus Österreich. Am Morgen haben die Frauen in der Küche gearbeitet. Sie haben einen großen Kuchen gebacken und viele gute Sachen gekocht. Die Männer haben den Garten vorbereitet und die Tische aufgestellt. Um 14 Uhr haben alle zusammen gegessen und getrunken. Nach dem Essen hat die Großmutter ihre Geschenke geöffnet. Sie hat sich sehr gefreut. Am Abend haben die jungen Leute Musik gemacht und getanzt. Das Fest war sehr schön, und alle sind spät nach Hause gefahren.
\end{textebox}

\textbf{A. Questions de compréhension}
\begin{enumerate}
\item Warum hat die Familie Berger gefeiert ?
\item Wer ist zum Fest gekommen ? Nenne zwei Personen.
\item Was haben die Frauen am Morgen gemacht ?
\item Wann hat die Großmutter ihre Geschenke geöffnet ?
\item Was haben die jungen Leute am Abend gemacht ?
\end{enumerate}

\textbf{B. Questions de langue}
\begin{enumerate}
\item Relève dans le texte trois participes passés et donne l'infinitif de chaque verbe.
\item Pourquoi dit-on \all{ist gekommen} mais \all{hat gefeiert} ? Explique la règle.
\item Mets la phrase \all{Die Großmutter hat ihre Geschenke geöffnet} au présent.
\end{enumerate}
\end{exercice}

\begin{exercice}[Thème et version (type Bac)]
\textbf{A. Thème.} Traduis le texte suivant en allemand. Fais attention au choix de l'auxiliaire, à la place du participe et aux majuscules.

« Dimanche dernier, nous avons fêté l'anniversaire de ma grand-mère. Toute la famille est venue et ma mère a préparé un bon repas. Mon frère a aidé ma mère dans la cuisine. Après le repas, nous avons chanté et dansé. Ma grand-mère était très contente. »

\textbf{B. Version.} Traduis le texte suivant en français.

\begin{textebox}[Zum Übersetzen]
Letztes Wochenende hat meine Familie ein Fest gefeiert. Viele Gäste sind zu uns gekommen. Meine Tante hat meiner Mutter in der Küche geholfen. Wir haben zusammen gegessen und viel gelacht. Am Abend sind alle Gäste glücklich nach Hause gefahren.
\end{textebox}
\end{exercice}

\begin{exercice}[Production écrite guidée (type Bac)]
\textbf{Sujet.} Tu organises une fête familiale (un anniversaire, un mariage ou une fête de fin d'année) chez toi à Yaoundé. Écris une lettre à ton ami allemand Thomas pour l'inviter.

\textbf{Consignes :}
\begin{itemize}
\item Rédige une lettre d'environ \textbf{80 à 100 mots} (au moins 8 phrases).
\item Respecte la structure d'une lettre : lieu et date, formule d'appel (\all{Lieber Thomas,}), introduction, corps, conclusion, formule de politesse (\all{Viele Grüße}) et signature.
\item Utilise \textbf{au moins 3 connecteurs} (par exemple : \all{zuerst}, \all{dann}, \all{danach}, \all{außerdem}, \all{zum Schluss}).
\item Emploie \textbf{au moins 2 phrases au Perfekt} et 1 phrase avec un verbe à particule séparable.
\item Donne la date, l'heure et le lieu de la fête, et dis ce que vous allez faire.
\item Relis-toi : majuscules des noms, place du verbe, orthographe.
\end{itemize}
\end{exercice}

\newpage

\coursec{Corrigés des exercices}

\begin{corrige}[Exercice 1 — QCM : la bonne forme]
\begin{enumerate}
\item \textbf{a) \all{haben}} — \all{Gestern haben wir unsere Großeltern besucht.} « Hier, nous avons rendu visite à nos grands-parents. » Le verbe \all{besuchen} (transitif) se conjugue au Perfekt avec \all{haben} ; à la 1\iere{} personne du pluriel : \all{haben}.
\item \textbf{b) \all{ist}} — \all{Meine Schwester ist gestern zu Hause geblieben.} « Ma sœur est restée à la maison hier. » \all{bleiben} est un verbe d'état/de non-mouvement qui se conjugue avec \all{sein} au Perfekt.
\item \textbf{a) \all{feiern}} — \all{Am Samstag feiern wir ein großes Fest.} « Samedi, nous fêtons une grande fête. » Sujet \all{wir} : terminaison \all{-en}. Le complément de temps en tête déclenche l'inversion : verbe en 2\ieme{} position.
\item \textbf{c) \all{dir}} — \all{Ich gratuliere dir zum Geburtstag.} « Je te félicite pour ton anniversaire. » \all{jemandem gratulieren} se construit avec le \textbf{datif} : \all{dir} (et non \all{dich}).
\item \textbf{a) \all{hat}} — \all{Meine Tante hat uns gestern eingeladen.} « Ma tante nous a invités hier. » \all{einladen} est transitif (auxiliaire \all{haben}) ; 3\ieme{} personne du singulier : \all{hat}.
\end{enumerate}
\end{corrige}

\begin{corrige}[Exercice 2 — Vrai ou faux ?]
\begin{enumerate}
\item \textbf{Faux.} \all{kommen} est un verbe de mouvement : il se conjugue avec \all{sein}. Correction : \all{Ich bin gekommen.} « Je suis venu. »
\item \textbf{Vrai.} \all{essen} est un verbe fort : \all{essen -- aß -- hat gegessen}. Exemple : \all{Wir haben einen Kuchen gegessen.} « Nous avons mangé un gâteau. »
\item \textbf{Vrai.} Au Perfekt, l'auxiliaire conjugué occupe la 2\ieme{} position et le participe passé se place \textbf{à la fin} de la phrase : \all{Wir haben gestern ein Fest gefeiert.}
\item \textbf{Faux.} \all{helfen} se construit avec le \textbf{datif} : \all{jemandem helfen}. Correction : \all{Ich helfe meiner Mutter.} « J'aide ma mère » (et non \all{meine Mutter}).
\item \textbf{Vrai.} En allemand, tous les noms communs prennent une majuscule : \all{das Fest}, \all{die Familie}, \all{der Kuchen}. C'est une différence majeure avec le français.
\end{enumerate}
\end{corrige}

\begin{corrige}[Exercice 3 — Vocabulaire : la famille et les fêtes]
\begin{enumerate}
\item \all{Meine Großmutter wird 80 Jahre alt : wir feiern ihren \textbf{Geburtstag}.} « Ma grand-mère va avoir 80 ans : nous fêtons son anniversaire. »
\item \all{Paul und Lena heiraten im Juni : die \textbf{Hochzeit} findet in der Kirche statt.} « Paul et Lena se marient en juin : le mariage a lieu à l'église. »
\item \all{Ich schreibe eine \textbf{Einladung} an meinen Freund in Douala.} « J'écris une invitation à mon ami à Douala. »
\item \all{Alle \textbf{Verwandten} kommen aus Bamenda und Bafoussam zum Fest.} « Tous les parents viennent de Bamenda et Bafoussam pour la fête. »
\item \all{Für meine Mutter kaufe ich ein schönes \textbf{Geschenk}.} « Pour ma mère, j'achète un beau cadeau. »
\end{enumerate}
\textbf{Rappel des articles et pluriels :} \all{der Geburtstag, -e} ; \all{die Hochzeit, -en} ; \all{die Einladung, -en} ; \all{der Verwandte, -n} (ici pluriel \all{die Verwandten}) ; \all{das Geschenk, -e}.
\end{corrige}

\begin{corrige}[Exercice 4 — Conjugaison : présent ou parfait ?]
\begin{enumerate}
\item \all{Meine Familie \textbf{wohnt} in Yaoundé.} « Ma famille habite à Yaoundé. » (Präsens, 3\ieme{} pers. sg. : \all{-t})
\item \all{Letztes Wochenende \textbf{haben} wir unsere Tante \textbf{besucht}.} « Le week-end dernier, nous avons rendu visite à notre tante. » (verbe faible transitif : \all{haben} + \all{besucht})
\item \all{Mein Bruder \textbf{hilft} meiner Mutter in der Küche.} « Mon frère aide ma mère dans la cuisine. » (verbe fort à voyelle changeante : \all{helfen} $\rightarrow$ \all{hilft} ; \all{meiner Mutter} = datif)
\item \all{Wir \textbf{haben} gestern einen großen Kuchen \textbf{gegessen}.} « Nous avons mangé un grand gâteau hier. » (verbe fort : participe \all{gegessen})
\item \all{Meine Eltern \textbf{sind} sehr zufrieden \textbf{gewesen}.} « Mes parents ont été très contents. » (\all{sein} se conjugue avec \all{sein} ; participe \all{gewesen})
\end{enumerate}
\end{corrige}

\begin{corrige}[Exercice 5 — L'ordre des mots]
\begin{enumerate}
\item \all{Meine Tante hat uns gestern eingeladen.} « Ma tante nous a invités hier. » (auxiliaire en 2\ieme{} position, participe \all{eingeladen} à la fin)
\item \all{Am Sonntag haben wir meine Großeltern besucht.} « Dimanche, nous avons rendu visite à mes grands-parents. » (le complément de temps en tête provoque l'inversion sujet-verbe)
\item \all{Ich kaufe meiner Schwester heute ein Geschenk.} « J'achète un cadeau à ma sœur aujourd'hui. » (\all{meiner Schwester} = datif, placé avant l'accusatif)
\item \all{Die Familie ist letzte Woche nach München gefahren.} « La famille est allée à Munich la semaine dernière. » (\all{fahren} = mouvement : auxiliaire \all{sein})
\item \all{Mein Bruder steht jeden Tag früh auf.} « Mon frère se lève tôt tous les jours. » (verbe à particule séparable \all{aufstehen} : la particule \all{auf} va à la fin)
\end{enumerate}
\textbf{Point de vigilance :} dans tous les cas, le verbe conjugué reste en \textbf{2\ieme{} position} ; c'est la règle d'or de la phrase déclarative allemande.
\end{corrige}

\begin{corrige}[Exercice 6 — Texte à trous : le parfait]
\all{Gestern \textbf{haben} wir den Geburtstag meines Vaters \textbf{gefeiert}. Viele Gäste \textbf{sind} zu uns \textbf{gekommen}. Meine Mutter \textbf{hat} ein gutes Essen \textbf{gekocht}. Mein Bruder \textbf{ist} lange in der Küche \textbf{geblieben} und er \textbf{hat} meiner Mutter \textbf{geholfen}. Nach dem Essen \textbf{haben} wir alle zusammen \textbf{gesungen}. Mein Vater \textbf{hat} sich sehr \textbf{gefreut}.}

« Hier, nous avons fêté l'anniversaire de mon père. Beaucoup d'invités sont venus chez nous. Ma mère a cuisiné un bon repas. Mon frère est resté longtemps dans la cuisine et il a aidé ma mère. Après le repas, nous avons tous chanté ensemble. Mon père s'est beaucoup réjoui. »

\textbf{Justifications :}
\begin{itemize}
\item \all{feiern, kochen, helfen, singen, sich freuen} : verbes transitifs ou sans mouvement $\rightarrow$ auxiliaire \all{haben}.
\item \all{kommen} (mouvement) et \all{bleiben} (exception) $\rightarrow$ auxiliaire \all{sein}.
\item Participes : \all{gefeiert}, \all{gekommen}, \all{gekocht}, \all{geblieben}, \all{geholfen}, \all{gesungen}, \all{gefreut}. Attention : \all{singen} et \all{helfen} sont des verbes forts (participe en \all{-en}).
\end{itemize}
\end{corrige}

\begin{corrige}[Exercice 7 — Appariements : invitations et réponses]
\begin{center}
\begin{tabular}{|l|l|}
\hline
\rowcolor{popblueL} Invitation & Réponse \\
\hline
1. \all{Möchtest du zu meiner Geburtstagsfeier kommen?} & \textbf{c)} \all{Natürlich! Ich freue mich schon sehr darauf.} \\
\hline
2. \all{Kannst du am Samstag zu unserer Hochzeit kommen?} & \textbf{b)} \all{Leider kann ich nicht, ich muss meiner Mutter helfen.} \\
\hline
3. \all{Hast du Lust, mit uns ins Restaurant zu gehen?} & \textbf{d)} \all{Gute Idee! Um wie viel Uhr treffen wir uns?} \\
\hline
4. \all{Kommst du morgen zu unserem Familienfest?} & \textbf{a)} \all{Ja, gern! Wann fängt die Feier an?} \\
\hline
\end{tabular}
\end{center}
\textbf{Traductions :} 1. « Veux-tu venir à ma fête d'anniversaire ? » — c) « Bien sûr ! Je me réjouis déjà beaucoup. » ; 2. « Peux-tu venir à notre mariage samedi ? » — b) « Malheureusement je ne peux pas, je dois aider ma mère. » ; 3. « As-tu envie d'aller au restaurant avec nous ? » — d) « Bonne idée ! À quelle heure est-ce qu'on se retrouve ? » ; 4. « Viens-tu demain à notre fête de famille ? » — a) « Oui, volontiers ! Quand commence la fête ? »
\end{corrige}

\begin{corrige}[Exercice 8 — Thème : phrases à traduire]
\begin{enumerate}
\item \all{Meine Großmutter wohnt in Douala bei meinem Onkel.} (\all{bei} + datif : \all{meinem Onkel} ; « habiter chez » = \all{wohnen bei})
\item \all{Gestern haben wir unsere Großeltern besucht.} (Perfekt avec \all{haben} ; participe \all{besucht} à la fin)
\item \all{Meine Schwester hat meiner Mutter in der Küche geholfen.} (\all{helfen} + datif : \all{meiner Mutter} ; participe fort \all{geholfen})
\item \all{Letzten Sonntag ist die ganze Familie zu uns gekommen.} (\all{kommen} = mouvement : auxiliaire \all{sein} ; « dimanche dernier » = \all{letzten Sonntag}, accusatif de temps)
\item \all{Ich habe meinem Bruder ein schönes Geschenk zum Geburtstag gegeben.} ou \all{... geschenkt.} (\all{geben/schenken} + datif de la personne : \all{meinem Bruder} ; « pour son anniversaire » = \all{zum Geburtstag})
\end{enumerate}
\textbf{Points de vigilance :} majuscule à tous les noms (\all{Großmutter}, \all{Küche}, \all{Geschenk}) ; verbe conjugué en 2\ieme{} position ; participe passé rejeté à la fin ; ne pas traduire « à ma mère » par un accusatif après \all{helfen}.
\end{corrige}

\begin{corrige}[Exercice 9 — Compréhension et étude de texte (type Bac)]
\textbf{Barème indicatif :} A. compréhension : 5 $\times$ 2 pts = 10 pts ; B. questions de langue : 4 + 3 + 3 = 10 pts.

\textbf{A. Questions de compréhension — réponses modèles :}
\begin{enumerate}
\item \all{Die Familie Berger hat gefeiert, weil die Großmutter 75 Jahre alt geworden ist.} « La famille Berger a fêté parce que la grand-mère a eu 75 ans. »
\item \all{Zum Fest sind alle Verwandten gekommen, zum Beispiel die Kinder und die Tante aus Österreich.} « Tous les parents sont venus, par exemple les enfants et la tante d'Autriche. »
\item \all{Am Morgen haben die Frauen in der Küche gearbeitet : sie haben einen großen Kuchen gebacken und gute Sachen gekocht.} « Le matin, les femmes ont travaillé dans la cuisine : elles ont cuit un grand gâteau et préparé de bonnes choses. »
\item \all{Nach dem Essen hat die Großmutter ihre Geschenke geöffnet.} « Après le repas, la grand-mère a ouvert ses cadeaux. »
\item \all{Am Abend haben die jungen Leute Musik gemacht und getanzt.} « Le soir, les jeunes ont fait de la musique et ont dansé. »
\end{enumerate}

\textbf{B. Questions de langue :}
\begin{enumerate}
\item Trois participes possibles parmi : \all{gefeiert} (\all{feiern}), \all{geworden} (\all{werden}), \all{gekommen} (\all{kommen}), \all{gearbeitet} (\all{arbeiten}), \all{gebacken} (\all{backen}), \all{gekocht} (\all{kochen}), \all{vorbereitet} (\all{vorbereiten}), \all{aufgestellt} (\all{aufstellen}), \all{gegessen} (\all{essen}), \all{getrunken} (\all{trinken}), \all{geöffnet} (\all{öffnen}), \all{gefreut} (\all{sich freuen}), \all{getanzt} (\all{tanzen}), \all{gefahren} (\all{fahren}).
\item On dit \all{ist gekommen} parce que \all{kommen} exprime un \textbf{mouvement, un déplacement} : au Perfekt, les verbes de mouvement (et \all{sein}, \all{bleiben}, \all{werden}) se conjuguent avec l'auxiliaire \all{sein}. On dit \all{hat gefeiert} parce que \all{feiern} est un verbe \textbf{transitif} (on fête quelque chose) : il se conjugue avec \all{haben}, comme la majorité des verbes.
\item Au présent : \all{Die Großmutter öffnet ihre Geschenke.} « La grand-mère ouvre ses cadeaux. » (verbe régulier, 3\ieme{} pers. sg. : \all{öffnet} ; le participe disparaît et le verbe reprend la 2\ieme{} position).
\end{enumerate}
\textbf{Points de vigilance :} répondre par des phrases complètes commençant par le verbe en 2\ieme{} position ; ne pas recopier le texte sans adapter la personne ; soigner les majuscules des noms.
\end{corrige}

\begin{corrige}[Exercice 10 — Thème et version (type Bac)]
\textbf{Barème indicatif :} thème 10 pts (2 pts par phrase : choix de l'auxiliaire, participe, ordre des mots, majuscules) ; version 10 pts (fidélité et qualité du français).

\textbf{A. Thème — traduction modèle :}

\all{Letzten Sonntag haben wir den Geburtstag meiner Großmutter gefeiert. Die ganze Familie ist gekommen, und meine Mutter hat ein gutes Essen vorbereitet. Mein Bruder hat meiner Mutter in der Küche geholfen. Nach dem Essen haben wir gesungen und getanzt. Meine Großmutter war sehr froh.}

\textbf{Points de vigilance :}
\begin{itemize}
\item « Dimanche dernier » = \all{letzten Sonntag} (accusatif masculin : \all{letzten}).
\item « est venue » : \all{kommen} + \all{sein} $\rightarrow$ \all{ist gekommen}.
\item « a aidé ma mère » : datif obligatoire $\rightarrow$ \all{hat meiner Mutter geholfen}.
\item « était contente » : ici le Präteritum de \all{sein} est préférable $\rightarrow$ \all{war sehr froh} (\all{ist sehr froh gewesen} est accepté).
\item Participes en fin de phrase : \all{gefeiert}, \all{vorbereitet} (pas de \all{ge-} malgré la particule séparable \all{vor-} : \all{vorbereiten} $\rightarrow$ \all{vorbereitet}, cas particulier à retenir).
\end{itemize}

\textbf{B. Version — traduction modèle :}

« Le week-end dernier, ma famille a fêté une fête. Beaucoup d'invités sont venus chez nous. Ma tante a aidé ma mère dans la cuisine. Nous avons mangé ensemble et nous avons beaucoup ri. Le soir, tous les invités sont rentrés chez eux, heureux. »

\textbf{Points de vigilance :} \all{zu uns gekommen} = « venus chez nous » ; \all{glücklich nach Hause gefahren} = « rentrés chez eux heureux / joyeux » ; rendre le Perfekt allemand par un passé composé français.
\end{corrige}

\begin{corrige}[Exercice 11 — Production écrite guidée (type Bac)]
\textbf{Barème indicatif (20 pts) :} respect de la tâche et de la longueur (4 pts) ; structure de la lettre (4 pts) ; connecteurs et contraintes grammaticales respectées (4 pts) ; correction grammaticale : ordre des mots, Perfekt, cas (4 pts) ; orthographe, majuscules, ponctuation (4 pts).

\textbf{Proposition de rédaction modèle :}

\begin{textebox}[Eine Einladung]
Yaoundé, den 15. Mai

Lieber Thomas,

wie geht es dir? Mir geht es gut. Ich schreibe dir, weil ich dich zu einem großen Fest einladen möchte. Am Samstag, am 25. Mai, feiern wir den Geburtstag meiner Großmutter bei uns in Yaoundé. Das Fest fängt um 15 Uhr an.

Zuerst essen wir zusammen : meine Mutter hat gestern schon einen großen Kuchen gebacken. Dann singen und tanzen wir. Danach öffnet meine Großmutter ihre Geschenke. Außerdem gibt es Musik und viele gute Sachen zu trinken. Letztes Jahr haben wir sehr viel gelacht, und alle Gäste sind spät nach Hause gefahren.

Kommst du? Ich freue mich sehr auf dich. Zum Schluss sage ich : Bitte ruf mich an!

Viele Grüße

Dein Freund Junior
\end{textebox}

\textbf{Vérification des consignes :}
\begin{itemize}
\item Structure : lieu et date (\all{Yaoundé, den 15. Mai}), appel (\all{Lieber Thomas,}), corps, formule finale (\all{Viele Grüße}), signature. Après l'appel, on va à la ligne et on commence par une \textbf{minuscule} (\all{wie geht es dir?}).
\item Connecteurs : \all{zuerst}, \all{dann}, \all{danach}, \all{außerdem}, \all{zum Schluss} (5 connecteurs, minimum 3 exigés).
\item Perfekt : \all{meine Mutter hat gestern schon einen großen Kuchen gebacken} ; \all{wir haben sehr viel gelacht} ; \all{alle Gäste sind spät nach Hause gefahren}.
\item Verbe à particule séparable : \all{das Fest fängt um 15 Uhr an} (\all{anfangen}) ; \all{bitte ruf mich an} (\all{anrufen}).
\item Informations pratiques : date (25 mai), heure (15 h), lieu (chez moi à Yaoundé), programme.
\end{itemize}
\textbf{Points de vigilance :} virgule après l'appel puis minuscule ; \all{den 15. Mai} (accusatif de date) ; inversion après un complément en tête (\all{Letztes Jahr haben wir...}) ; ne jamais écrire \all{Ich habe gebacken einen Kuchen} : le participe reste à la fin.
\end{corrige}

\newpage

\coursec{Fiche bilan}

\begin{popfigure}
\begin{tikzpicture}[mindmap, grow cyclic, every node/.style={concept, text=white, font=\small\bfseries, align=center},
root concept/.append style={concept color=popblue, font=\large\bfseries},
level 1/.append style={level distance=3.4cm, sibling angle=51.4},
level 2/.append style={level distance=2.1cm, sibling angle=45, font=\scriptsize}]
\node[root concept, align=center]{Vie familiale et sociale\\ ALA04}
  child[concept color=poporange]{ node[align=center]{Lexique\\ Familie \& Gesellschaft}
    child{ node[align=center]{der Verwandte\\ die Feier} }
    child{ node[align=center]{sich freuen\\ einladen} }
  }
  child[concept color=popgreen]{ node[align=center]{Présent\\ (Révision)}
    child{ node{verbe en 2\up{e} place} }
  }
  child[concept color=poppurple]{ node{Perfekt}
    child{ node[align=center]{haben/sein\\ + Partizip II} }
    child{ node[align=center]{participe\\ en fin} }
  }
  child[concept color=poppink]{ node[align=center]{Ordre\\ des mots}
    child{ node{TeKaMoLo} }
    child{ node[align=center]{subordonnée :\\ verbe final} }
  }
  child[concept color=popgold]{ node{Production}
    child{ node[align=center]{lettre\\ description\\ dialogue} }
  }
  child[concept color=popdark]{ node[align=center]{Traduction\\ thème \& version}
    child{ node[align=center]{identifier\\ temps \& cas} }
  };
\end{tikzpicture}
\legende{Carte mentale de la leçon ALA04 : la vie en famille et en société}
\end{popfigure}

\begin{bilan}
\textbf{1. Lexique clé (à connaître par cœur)}
\begin{itemize}
\item \all{die Familie, -n} : la famille ; \all{die Eltern} (pl.) : les parents ; \all{der Bruder, „-} : le frère ; \all{die Schwester, -n} : la sœur ; \all{die Großeltern} (pl.) : les grands-parents ; \all{der Verwandte, -n} : le parent (membre de la famille).
\item \all{das Fest, -e} : la fête ; \all{die Feier, -n} : la célébration ; \all{der Geburtstag, -e} : l'anniversaire ; \all{die Hochzeit, -en} : le mariage ; \all{das Geschenk, -e} : le cadeau ; \all{der Gast, „-e} : l'invité.
\item \all{einladen (lädt ein, lud ein, hat eingeladen)} : inviter ; \all{feiern} : fêter ; \all{sich freuen (über + A / auf + A)} : se réjouir (de / à l'idée de) ; \all{schenken} : offrir ; \all{helfen (hilft, half, hat geholfen) + D} : aider.
\item \all{der Nachbar, -n} : le voisin ; \all{die Gemeinde, -n} : la communauté ; \all{zusammen} : ensemble ; \all{sich treffen (trifft, traf, hat getroffen)} : se retrouver.
\end{itemize}

\textbf{2. Grammaire : les règles à connaître par cœur}
\begin{center}
\begin{tabularx}{\linewidth}{|l|X|X|}
\hline
\rowcolor{popblueL} \textbf{Point} & \textbf{Règle} & \textbf{Exemple} \\
\hline
Présent & Radical + -e, -st, -t, -en, -t, -en ; verbes forts : changement de voyelle à la 2\up{e}/3\up{e} pers. sg. & \all{er hilft, sie lädt ein} \\
\hline
Perfekt & Auxiliaire \all{haben} ou \all{sein} (mouvement, changement d'état) + Partizip II en fin de phrase & \all{Wir haben gefeiert.} / \all{Sie ist gekommen.} \\
\hline
Partizip II & \all{ge- + radical + -t} (faibles) ; \all{ge- + -en} (forts) ; pas de \all{ge-} pour les verbes en \all{-ieren} et à préfixe inséparable & \all{gemacht, gegessen, studiert, besucht} \\
\hline
Verbe à particule & Au présent : particule en fin ; au Perfekt : \all{ge-} entre particule et radical & \all{Ich lade ein.} / \all{Ich habe eingeladen.} \\
\hline
Ordre des mots & Verbe conjugué en 2\up{e} position ; compléments : TeKaMoLo (temps, cause, manière, lieu) & \all{Gestern habe ich meine Tante besucht.} \\
\hline
Subordonnée & Verbe conjugué rejeté à la fin (avec \all{weil, dass, wenn}...) & \all{..., weil sie Geburtstag hat.} \\
\hline
\end{tabularx}
\end{center}

\textbf{3. Méthodes express}
\begin{itemize}
\item \textbf{Lettre courte} : \all{Liebe/Lieber ...}, introduction (\all{wie geht es dir?}), corps, salutation (\all{Viele Grüße / Deine ...}).
\item \textbf{Description} : présent + adjectifs ; ordre : qui ? où ? quand ? que fait-on ?
\item \textbf{Dialogue} : tirets ou répliques nommées ; verbes de parole au présent ; registre familier (\all{du}).
\item \textbf{Thème (F $\rightarrow$ D)} : repérer le temps, conjuguer, placer le verbe en 2\up{e} position, vérifier les cas et les majuscules.
\item \textbf{Version (D $\rightarrow$ F)} : identifier sujet/verbe/compléments, repérer le participe final, traduire naturellement en français.
\item \textbf{Relecture} : verbe en 2\up{e} place ? participe en fin ? majuscules aux noms ? Umlaute ?
\end{itemize}
\end{bilan}

\begin{aretenir}[Checklist avant le Bac]
Avant de rendre ma copie, je vérifie :
\begin{itemize}
\item[$\square$] Le verbe conjugué est en deuxième position dans chaque phrase principale.
\item[$\square$] Au Perfekt, l'auxiliaire \all{haben}/\all{sein} est correct et le Partizip II est en fin de phrase.
\item[$\square$] Les verbes à particule sont séparés au présent (\all{ich lade ein}) et soudés au Perfekt (\all{eingeladen}).
\item[$\square$] Tous les noms allemands portent une majuscule.
\item[$\square$] Les cas sont corrects : accusatif après \all{für, durch} ; datif après \all{mit, zu, von} et après \all{helfen, danken}.
\item[$\square$] Dans les subordonnées (\all{weil, dass, wenn}), le verbe est à la fin.
\item[$\square$] Les Umlaute (\all{ä, ö, ü}) et le \all{ß} sont écrits correctement.
\item[$\square$] Ma lettre contient une formule d'appel et une formule de politesse.
\item[$\square$] Mon dialogue utilise le tutoiement et des répliques claires.
\item[$\square$] En version, ma traduction française est naturelle et complète (rien n'est oublié).
\item[$\square$] En thème, j'ai respecté l'ordre TeKaMoLo pour les compléments.
\item[$\square$] Je me suis relu au moins une fois pour l'orthographe et les accords.
\end{itemize}
\end{aretenir}

\findecours

\end{document}
